% Le Droit et la Justice — Philosophie 1re Tech — Cours signé OnBuch+
% Programme officiel MINESEC. Compiler avec : tectonic N08-droit-justice.tex
\newcommand{\DOCMATIERE}{Philosophie}
\newcommand{\DOCNIVEAU}{1re Tech}
\newcommand{\DOCMODULE}{Philosophie – classe de Première technique}
\newcommand{\DOCLECON}{N08}
\newcommand{\DOCTITRE}{Le Droit et la Justice}
% OnBuch+ — préambule « cours signés » ENRICHI (compilé avec tectonic / XeLaTeX)
% Système visuel « Pop » d'OnBuch+ : fond crème (jamais blanc), bordures encre
% épaisses, ombres « dures » décalées, titres Archivo Black en capitales.
% Version MATHÉMATIQUES : graphes (pgfplots/tikz), boîtes pédagogiques
% (définition, propriété, méthode, attention, le savais-tu, exercice résolu,
% exercice, TP, bilan), page de garde et signature OnBuch+ ancrée en bas.
%
% Macros attendues dans main.tex AVANT \input{preamble} :
%   \DOCMATIERE  \DOCNIVEAU  \DOCMODULE  \DOCLECON (numéro)  \DOCTITRE
\documentclass[11pt]{article}

\usepackage{fontspec}
\usepackage[french]{babel}
\usepackage[a4paper,margin=2cm,top=3.4cm,bottom=2.8cm,headheight=1.4cm,footskip=1.3cm]{geometry}
\usepackage{amsmath,amssymb,amsfonts}
\usepackage{mathtools} % \xmapsto, \coloneqq... souvent utilisés par les modèles
\usepackage{cancel} % simplifications barrées (bilans de forces)
% \abs{z} (module, valeur absolue) et \norm{u} (norme), utilisés par les modèles
\providecommand{\abs}{}\let\abs\relax\DeclarePairedDelimiter{\abs}{\lvert}{\rvert}
\providecommand{\norm}{}\let\norm\relax\DeclarePairedDelimiter{\norm}{\lVert}{\rVert}
\usepackage[table]{xcolor} % [table] : fournit aussi \rowcolors (alternance de lignes)
\usepackage{pagecolor}
\usepackage{tikz}
\usetikzlibrary{arrows.meta,positioning,calc,shapes.geometric,shapes.misc,shapes.arrows,decorations.pathmorphing,decorations.pathreplacing,decorations.markings,patterns,shadows,mindmap,trees,fit,backgrounds,matrix,intersections,angles,quotes}
\usepackage{pgfplots}
\usepackage[version=4]{mhchem} % équations nucléaires \ce{^{235}_{92}U}
\usepackage[european,siunitx]{circuitikz} % schémas électriques
\pgfplotsset{compat=1.18}
\usepgfplotslibrary{fillbetween,groupplots} % aires entre courbes (calcul intégral)
\usepackage{enumitem}
\usepackage{fancyhdr}
% [most] charge toutes les bibliothèques tcolorbox (listings, minted, raster,
% documentation...) et gonfle inutilement la mémoire TeX sur un document long
% (~300 boîtes) : on ne charge que ce qui est réellement utilisé.
\usepackage[skins,breakable]{tcolorbox}
\usepackage{graphicx}
\usepackage{colortbl}
\usepackage{booktabs}
\usepackage{tabularx}
\usepackage{multirow}
\usepackage{diagbox} % case d'en-tête diagonale des tableaux à double entrée
% Colonnes extensibles centrée / gauche / droite (C, L, R), souvent utilisées par les modèles
\newcolumntype{C}{>{\centering\arraybackslash}X}
\newcolumntype{L}{>{\raggedright\arraybackslash}X}
\newcolumntype{R}{>{\raggedleft\arraybackslash}X}
\usepackage{array}
\usepackage{multicol}
\usepackage{siunitx}
% parse-numbers=false : accepte aussi \SI{2,5 \times 10^{-3}}{...}
\sisetup{locale=FR,output-decimal-marker={,},group-minimum-digits=4,parse-numbers=false}
\DeclareSIUnit\ohm{\ensuremath{\symup{\Omega}}} % U+2126 absent de la police
\DeclareSIUnit\division{div} % réglages d'oscilloscope (ms/div, V/div)
\DeclareSIUnit\tr{tr} % tours (vitesse de rotation, ex. \SI{1500}{\tr\per\minute})
\DeclareSIUnit\molar{M} % concentration molaire (ex. \SI{3}{\molar}, \SI{1}{\micro\molar})
\DeclareSIUnit\an{an} % année (le modèle confond parfois avec \year, un compteur TeX) — ex. \SI{5}{\kilo\meter\per\an}
\DeclareSIUnit\FCFA{FCFA} % franc CFA (ex. \SI{500}{\FCFA\per\kilo\gram})
\DeclareSIUnit\degreeBrix{\textdegree Brix} % teneur en sucre d'un sirop/jus (réfractométrie)
\DeclareSIUnit\month{mois} % le modèle utilise parfois \month, un compteur TeX, comme unité de durée
\DeclareSIUnit\microliter{\micro\liter} % le modèle écrit parfois \microliter en un seul token
\DeclareSIUnit\million{millions} % dénombrement (ex. \SI{15}{\million\per\mL} spermatozoïdes)
\usepackage{qrcode}
% \degree : commande fréquemment utilisée par les modèles pour un angle ou une
% température en dehors de \SI{}{} (non fournie par les paquets chargés).
\providecommand{\degree}{^{\circ}}
% \qed (fin de démonstration), utilisé par les modèles sans amsthm
\providecommand{\qed}{\hfill$\square$}
% \barre : double barre des valeurs interdites dans un tableau de variations (tkz-tab)
\providecommand{\barre}{\ensuremath{\|}}
% environnement proof (amsthm non chargé) utilisé par les modèles
\ifdefined\proof\else\newenvironment{proof}[1][Démonstration]{\par\noindent\textit{#1.}\ }{\qed\par}\fi
\usepackage{lastpage}
\usepackage{hyperref}
\hypersetup{hidelinks}

% ============================================================
% Palette « Pop » OnBuch+ — mêmes hex que la classe P (plus_ui.dart)
% ============================================================
\definecolor{poporange}{HTML}{F59321}
\definecolor{popdark}{HTML}{E07A0C}
\definecolor{popcream}{HTML}{FFF6E8}
\definecolor{popink}{HTML}{1C1714}
\definecolor{poppurple}{HTML}{7A5AE0}
\definecolor{popgreen}{HTML}{2E9E6A}
\definecolor{poppink}{HTML}{E0457B}
\definecolor{popblue}{HTML}{2F7DE1}
\definecolor{popgold}{HTML}{C98A12}
\definecolor{popmuted}{HTML}{8A7F76}
\definecolor{popline}{HTML}{EADFD2}
\definecolor{popgray}{HTML}{8A7F76}
\colorlet{popgrey}{popgray}
% Alias tolérants (le modèle nomme parfois les couleurs différemment)
\colorlet{popwhite}{white}
\colorlet{popblack}{popink}
\colorlet{popred}{poppink}
\colorlet{popyellow}{popgold}
% Teintes claires (fonds de tableaux/encadrés) — le modèle nomme parfois
% les couleurs avec un suffixe L (ex. popblueL) sans les avoir définies.
\colorlet{poporangeL}{poporange!20}
\colorlet{popdarkL}{popdark!20}
\colorlet{poppurpleL}{poppurple!20}
\colorlet{popgreenL}{popgreen!20}
\colorlet{poppinkL}{poppink!20}
\colorlet{popblueL}{popblue!20}
\colorlet{popgoldL}{popgold!20}
\colorlet{popredL}{popred!20}
\colorlet{popgrayL}{popgray!20}
% Teintes claires pour fonds de tableaux / zones
\colorlet{popblueL}{popblue!10!white}
\colorlet{poporangeL}{poporange!14!white}
\colorlet{popgreenL}{popgreen!12!white}
\colorlet{poppurpleL}{poppurple!12!white}
\colorlet{poppinkL}{poppink!12!white}
\colorlet{popgoldL}{popgold!14!white}

\pagecolor{popcream}

% ============================================================
% Typographie
% ============================================================
\defaultfontfeatures{Ligatures=TeX}
\setmainfont{PlusJakartaSans-Medium.ttf}[Path=../../../fonts/,BoldFont=PlusJakartaSans-Bold.ttf]
% Maths en sans-serif (Fira Math) avec chiffres et lettres droites de Plus
% Jakarta, pour que les formules se fondent dans le texte.
\usepackage[math-style=ISO,bold-style=ISO]{unicode-math}
\setmathfont{FiraMath-Regular.otf}[Path=../../../fonts/]
\setmathfont{PlusJakartaSans-Medium.ttf}[Path=../../../fonts/,range={up/num,up/{Latin,latin},\mathup}]
\usepackage{icomma} % virgule décimale sans espace en mode math
% Caractères mathématiques Unicode absents des polices de texte (Plus Jakarta,
% Archivo Black) : rendus par leur équivalent mathématique, en texte comme en
% formule — sinon ils s'affichent en carré vide (ex. « Arithmétique dans ℤ »).
\usepackage{newunicodechar}
\newunicodechar{ℝ}{\ensuremath{\mathbb{R}}}
\newunicodechar{ℤ}{\ensuremath{\mathbb{Z}}}
\newunicodechar{ℂ}{\ensuremath{\mathbb{C}}}
\newunicodechar{ℕ}{\ensuremath{\mathbb{N}}}
\newunicodechar{ℚ}{\ensuremath{\mathbb{Q}}}
\newunicodechar{𝔻}{\ensuremath{\mathbb{D}}}
\newunicodechar{α}{\ensuremath{\alpha}}
\newunicodechar{β}{\ensuremath{\beta}}
\newunicodechar{γ}{\ensuremath{\gamma}}
\newunicodechar{δ}{\ensuremath{\delta}}
\newunicodechar{ε}{\ensuremath{\varepsilon}}
\newunicodechar{θ}{\ensuremath{\theta}}
\newunicodechar{λ}{\ensuremath{\lambda}}
\newunicodechar{μ}{\ensuremath{\mu}}
\newunicodechar{σ}{\ensuremath{\sigma}}
\newunicodechar{τ}{\ensuremath{\tau}}
\newunicodechar{φ}{\ensuremath{\varphi}}
\newunicodechar{ψ}{\ensuremath{\psi}}
\newunicodechar{ω}{\ensuremath{\omega}}
\newunicodechar{Σ}{\ensuremath{\Sigma}}
% majuscules produites par \MakeUppercase dans les titres (α-aminés -> Α)
\newunicodechar{Α}{\ensuremath{\alpha}}
\newunicodechar{Β}{\ensuremath{\beta}}
\newunicodechar{Γ}{\ensuremath{\Gamma}}
\newunicodechar{Δ}{\ensuremath{\Delta}}
\newunicodechar{Ε}{\ensuremath{\varepsilon}}
\newunicodechar{Θ}{\ensuremath{\Theta}}
\newunicodechar{Λ}{\ensuremath{\Lambda}}
\newunicodechar{Μ}{\ensuremath{\mu}}
\newunicodechar{Τ}{\ensuremath{\tau}}
\newunicodechar{Φ}{\ensuremath{\Phi}}
\newunicodechar{Ψ}{\ensuremath{\Psi}}
\newunicodechar{Ω}{\ensuremath{\Omega}}
\newunicodechar{↦}{\ensuremath{\mapsto}}
\newunicodechar{∈}{\ensuremath{\in}}
\newunicodechar{∉}{\ensuremath{\notin}}
\newunicodechar{∧}{\ensuremath{\wedge}}
\newunicodechar{⇔}{\ensuremath{\Leftrightarrow}}
\newunicodechar{⟺}{\ensuremath{\Longleftrightarrow}}
\newunicodechar{⇒}{\ensuremath{\Rightarrow}}
\newunicodechar{⟹}{\ensuremath{\Longrightarrow}}
\newunicodechar{∘}{\ensuremath{\circ}}
\newunicodechar{∪}{\ensuremath{\cup}}
\newunicodechar{∩}{\ensuremath{\cap}}
\newunicodechar{≡}{\ensuremath{\equiv}}
\newunicodechar{⊂}{\ensuremath{\subset}}
\newunicodechar{⊥}{\ensuremath{\perp}}
\newunicodechar{∃}{\ensuremath{\exists}}
\newunicodechar{∀}{\ensuremath{\forall}}
\newunicodechar{∛}{\ensuremath{\sqrt[3]{\,}}}
\newunicodechar{∜}{\ensuremath{\sqrt[4]{\,}}}
\newunicodechar{⃗}{\ensuremath{{}^{\to}}}
\newunicodechar{ⁿ}{\ifmmode^{n}\else\textsuperscript{\ensuremath{n}}\fi}
\newunicodechar{ˣ}{\ifmmode^{x}\else\textsuperscript{\ensuremath{x}}\fi}
\newunicodechar{ᵇ}{\ifmmode^{b}\else\textsuperscript{\ensuremath{b}}\fi}
\newunicodechar{ʳ}{\ifmmode^{r}\else\textsuperscript{\ensuremath{r}}\fi}
\newunicodechar{ᵘ}{\ifmmode^{u}\else\textsuperscript{\ensuremath{u}}\fi}
\newunicodechar{ʸ}{\ifmmode^{y}\else\textsuperscript{\ensuremath{y}}\fi}
\newunicodechar{ᵃ}{\ifmmode^{a}\else\textsuperscript{\ensuremath{a}}\fi}
\newunicodechar{ᵉ}{\ifmmode^{e}\else\textsuperscript{\ensuremath{e}}\fi}
\newunicodechar{ᵏ}{\ifmmode^{k}\else\textsuperscript{\ensuremath{k}}\fi}
\newunicodechar{ᵖ}{\ifmmode^{p}\else\textsuperscript{\ensuremath{p}}\fi}
\newunicodechar{ᴺ}{\ifmmode^{N}\else\textsuperscript{\ensuremath{N}}\fi}
\newunicodechar{ᶜ}{\ifmmode^{c}\else\textsuperscript{\ensuremath{c}}\fi}
\newunicodechar{ʰ}{\ifmmode^{h}\else\textsuperscript{\ensuremath{h}}\fi}
\newunicodechar{ᵠ}{\ifmmode^{q}\else\textsuperscript{\ensuremath{q}}\fi}
\newunicodechar{ₙ}{\ifmmode_{n}\else\textsubscript{\ensuremath{n}}\fi}
\newunicodechar{ₐ}{\ifmmode_{a}\else\textsubscript{\ensuremath{a}}\fi}
\newunicodechar{ᵢ}{\ifmmode_{i}\else\textsubscript{\ensuremath{i}}\fi}
\newunicodechar{ᵣ}{\ifmmode_{r}\else\textsubscript{\ensuremath{r}}\fi}
\newunicodechar{ₖ}{\ifmmode_{k}\else\textsubscript{\ensuremath{k}}\fi}
\newunicodechar{ᵦ}{\ifmmode_{\beta}\else\textsubscript{\ensuremath{\beta}}\fi}
\newunicodechar{ₓ}{\ifmmode_{x}\else\textsubscript{\ensuremath{x}}\fi}
\newfontfamily\popdisplay{ArchivoBlack-Regular.ttf}[Path=../../../fonts/]
\newfontfamily\poplogo{SpaceGrotesk-Bold.ttf}[Path=../../../fonts/]
\newfontfamily\popsemi{PlusJakartaSans-SemiBold.ttf}[Path=../../../fonts/]
\newfontfamily\popxbold{PlusJakartaSans-ExtraBold.ttf}[Path=../../../fonts/]
\newfontfamily\popxboldls{PlusJakartaSans-ExtraBold.ttf}[Path=../../../fonts/,LetterSpace=3.0]

\setlist[enumerate]{leftmargin=1.7em,itemsep=5pt,topsep=4pt}
\setlist[itemize]{leftmargin=1.7em,itemsep=4pt,topsep=4pt,label={\color{poporange}\textbullet}}
\setlength{\parindent}{0pt}
\setlength{\parskip}{5pt}
\renewcommand{\arraystretch}{1.25}

% pgfplots : style Pop par défaut
\pgfplotsset{
  every axis/.append style={axis line style={popink,line width=1pt},tick style={popink},
    grid style={popline,line width=0.5pt},label style={font=\small},tick label style={font=\footnotesize}},
  popcurve/.style={poporange,line width=1.8pt,no markers,smooth},
}
% Même style disponible dans un \draw TikZ simple (hors pgfplots)
\tikzset{popcurve/.style={poporange,line width=1.8pt}}
\tikzset{popgrid/.style={popline,very thin}}
\colorlet{popcurve}{poporange} % le nom du style est parfois utilisé comme couleur

% ============================================================
% Pill / squiggle
% ============================================================
\newcommand{\poppill}[3]{%
  \tikz[baseline=-0.55ex]\node[draw=popink,line width=1.2pt,rounded corners=2.6pt,%
    fill=#2,inner xsep=6.5pt,inner ysep=3pt]{%
    \popxboldls\fontsize{7}{7}\selectfont\color{#3}\MakeUppercase{#1}};%
}
\newcommand{\popsquiggle}[1]{%
  \tikz[baseline=-0.4ex]{
    \draw[#1,line width=2.6pt,line cap=round]
      plot [smooth] coordinates {(0,0.09) (0.34,0.01) (0.68,0.21) (1.02,0.01) (1.36,0.10)};
  }%
}
% Mot-clé surligné (marqueur orange sous le texte)
\newcommand{\cle}[1]{{\popxbold\color{popdark}#1}}

% ============================================================
% En-tête / pied
% ============================================================
\pagestyle{fancy}
\fancyhf{}
\renewcommand{\headrulewidth}{0pt}
\renewcommand{\footrulewidth}{0pt}
\lhead{\raisebox{-0.15ex}{\poplogo\fontsize{13}{13}\selectfont\color{popink}OnBuch\color{poporange}+}}
\rhead{\poppill{\DOCMATIERE\ \textbullet\ \DOCNIVEAU\ \textbullet\ Leçon \DOCLECON}{white}{popink}}
\lfoot{\popxbold\fontsize{7.6}{7.6}\selectfont\color{popmuted}\MakeUppercase{Cours signé OnBuch+}\ \textbullet\ {\popsemi\DOCTITRE}}
\rfoot{\tikz[baseline=-0.6ex]\node[rounded corners=8.5pt,fill=popink,minimum height=17pt,minimum width=17pt,inner xsep=5pt,inner sep=0]{\popxbold\fontsize{8}{8}\selectfont\color{popcream}\thepage\,/\,\pageref*{LastPage}};}

% ============================================================
% Titres
% ============================================================
\newcommand{\chaptitre}[2]{%
  \begin{tcolorbox}[enhanced,colback=poporange,colframe=popink,boxrule=2.6pt,arc=11pt,
      drop shadow=popink,left=15pt,right=15pt,top=12pt,bottom=11pt,before skip=3pt,after skip=18pt]
    \poppill{#2}{popink}{popcream}\par\vspace{9pt}
    {\raggedright\hyphenpenalty=10000\exhyphenpenalty=10000\popdisplay\fontsize{22}{24}\selectfont\color{popcream}\MakeUppercase{#1}\par}
  \end{tcolorbox}
}
% Section numérotée : pastille orange avec le numéro + titre Archivo
\newcounter{popsec}
\newcommand{\coursec}[1]{%
  \stepcounter{popsec}\setcounter{popsub}{0}%
  \par\vspace{16pt}\noindent
  \begin{minipage}{\linewidth}
  \tikz[baseline=-0.7ex]\node[draw=popink,line width=1.6pt,fill=poporange,rounded corners=4pt,minimum size=22pt,inner sep=0]{\popdisplay\fontsize{12}{12}\selectfont\color{popcream}\thepopsec};%
  \hspace{8pt}{\popdisplay\fontsize{15}{17}\selectfont\color{popink}\MakeUppercase{#1}}\par
  \vspace{4pt}\noindent{\color{poporange}\rule{36pt}{3.6pt}}
  \end{minipage}\par\nopagebreak\vspace{8pt}
}
\newcounter{popsub}
\newcommand{\courssub}[1]{%
  \stepcounter{popsub}%
  \par\vspace{9pt}\noindent{\popxbold\fontsize{11.5}{13}\selectfont\color{popink}{\color{poporange}\thepopsec.\thepopsub}\hspace{6pt}#1}\par\nopagebreak\vspace{4pt}
}

% ============================================================
% Boîtes pédagogiques — même recette : fond blanc, bordure épaisse colorée,
% ombre dure encre, badge plein. Titre optionnel [#1].
% ============================================================
\tcbset{popbase/.style={breakable,enhanced,colback=white,boxrule=2.3pt,arc=9pt,
  shadow={3pt}{-3pt}{0pt}{popink},left=14pt,right=14pt,top=11pt,bottom=11pt,before skip=11pt,after skip=15pt,
  fontupper=\popsemi\fontsize{10.6}{14.5}\selectfont\color{popink},
  fontlower=\popsemi\fontsize{10.6}{14.5}\selectfont\color{popink}}}
% Coche dessinée (le glyphe ✓ n'existe pas dans les polices du document)
\newcommand{\popcheck}[1]{\tikz[baseline=-0.5ex]\draw[#1,line width=1.6pt,line cap=round,line join=round](-0.13,0.0)--(-0.04,-0.09)--(0.13,0.1);}
\providecommand{\coche}{\popcheck{popgreen}}
\providecommand{\resultat}[1]{\par\smallskip\noindent\fcolorbox{popgreen}{popgreenL}{\parbox{\dimexpr\linewidth-2\fboxsep-2\fboxrule\relax}{\textbf{Résultat.} #1}}\par\smallskip}
\newcommand{\popboxhead}[3]{\poppill{#1}{#2}{white}\ifx\relax#3\relax\else\hspace{6pt}{\popxbold\fontsize{10.6}{12}\selectfont\color{popink}#3}\fi\par\vspace{7pt}}

\newtcolorbox{aretenir}[1][]{popbase,colframe=popgreen,before upper={\popboxhead{À retenir}{popgreen}{#1}}}
\newtcolorbox{exemplebox}[1][]{popbase,colframe=poporange,before upper={\popboxhead{Exemple}{poporange}{#1}}}
\newtcolorbox{definition}[1][]{popbase,colframe=popblue,before upper={\popboxhead{Définition}{popblue}{#1}}}
\newtcolorbox{propriete}[1][]{popbase,colframe=poppurple,before upper={\popboxhead{Propriété}{poppurple}{#1}}}
\newtcolorbox{methode}[1][]{popbase,colframe=popgold,before upper={\popboxhead{Méthode}{popgold}{#1}}}
\newtcolorbox{attention}[1][]{popbase,colframe=poppink,before upper={\popboxhead{Attention}{poppink}{#1}}}
\newtcolorbox{experience}[1][]{popbase,colframe=popblue,colback=popblueL,before upper={\popboxhead{Expérience}{popblue}{#1}}}
% « Le savais-tu ? » : carte violette pleine (contexte, culture, Cameroun)
\newtcolorbox{savaistu}[1][]{popbase,colback=poppurple,colframe=popink,
  fontupper=\popsemi\fontsize{10.4}{14.2}\selectfont\color{white},
  before upper={\poppill{Le savais-tu\,?}{white}{poppurple}\ifx\relax#1\relax\else\hspace{6pt}{\popxbold\color{white}#1}\fi\par\vspace{7pt}}}
% Exercice résolu : énoncé en haut, \tcblower puis la solution sur fond crème
\newcounter{popexo}\newcounter{popexr}
\newtcolorbox{exoresolu}[1][]{popbase,colframe=poporange,colbacklower=popcream,
  segmentation style={popink,line width=1.2pt,dashed},
  before upper={\stepcounter{popexr}\popboxhead{Exercice résolu \thepopexr}{poporange}{#1}},
  before lower={\poppill{Solution}{popink}{popcream}\par\vspace{6pt}}}
% Exercice (non résolu en place) — la correction est dans la partie Corrigés
\newtcolorbox{exercice}[1][]{popbase,colframe=popink,
  before upper={\stepcounter{popexo}\popboxhead{Exercice \thepopexo}{popink}{#1}}}
% Corrigé
\newtcolorbox{corrige}[1][]{popbase,colframe=popgreen,colback=popgreenL,
  before upper={\popboxhead{Corrigé}{popgreen}{#1}}}
% Objectifs / prérequis / bilan
\newtcolorbox{objectifs}{popbase,colframe=popink,colback=poporangeL,
  before upper={\popboxhead{Objectifs de la leçon}{popink}{}}}
\newtcolorbox{prerequis}{popbase,colframe=popink,colback=popblueL,
  before upper={\popboxhead{Prérequis}{popblue}{}}}
\newtcolorbox{bilan}{popbase,colframe=popink,colback=white,boxrule=2.8pt,
  before upper={\popboxhead{Fiche bilan}{poporange}{L'essentiel à connaître pour l'examen}}}
% Figure encadrée avec légende
\newtcolorbox{popfigure}[1][]{enhanced,breakable=false,colback=white,colframe=popline,boxrule=1.4pt,arc=8pt,
  left=10pt,right=10pt,top=10pt,bottom=8pt,before skip=10pt,after skip=12pt,halign=center,
  lower separated=false,halign lower=center,
  fontlower=\popsemi\fontsize{9}{11.5}\selectfont\color{popmuted},
  #1}
\newcommand{\legende}[1]{\par\vspace{6pt}{\popsemi\fontsize{9}{11.5}\selectfont\color{popmuted}#1}}

% ============================================================
% Signature OnBuch+
% ============================================================
\newcommand{\poplogoinline}[1]{{\poplogo\fontsize{#1}{#1}\selectfont\color{popink}OnBuch\color{poporange}+}}
\newcommand{\signatureonbuch}{%
  \begin{tcolorbox}[enhanced,colback=popink,colframe=popink,boxrule=2.2pt,arc=10pt,
      drop shadow=poporange,left=14pt,right=14pt,top=10pt,bottom=10pt,before skip=0pt,after skip=0pt]
    \tikz[baseline=-0.7ex]\node[circle,fill=popgreen,draw=popcream,line width=1.3pt,minimum size=22pt,inner sep=0]{\popcheck{white}};%
    \hspace{9pt}\begin{minipage}[c]{0.62\linewidth}
      {\popxbold\fontsize{12}{13}\selectfont\color{popcream}Signé\ \textbullet\ {\poplogo OnBuch\color{poporange}+}}\par\vspace{2pt}
      {\raggedright\popsemi\fontsize{8}{10.5}\selectfont\color{popline}\MakeUppercase{Équipe pédagogique OnBuch+}\\\MakeUppercase{Conforme au programme officiel MINESEC}\par}
    \end{minipage}\hfill
    {\poplogo\fontsize{22}{22}\selectfont\color{popcream}OnBuch\color{poporange}+}
  \end{tcolorbox}%
}

% ============================================================
% Page de garde
% #1 titre · #2 matière · #3 niveau · #4 module · #5 n° leçon · #6 durée ·
% #7 description / objectifs (texte ou itemize)
% ============================================================
\newcommand{\pagegarde}[7]{%
  \thispagestyle{empty}
  \newgeometry{a4paper,margin=1.7cm,top=1.6cm,bottom=1.4cm}%
  \begin{tikzpicture}[remember picture,overlay]
    \fill[poporange!18] ([xshift=-3.2cm,yshift=-3.6cm]current page.north east) circle (4.6cm);
    \fill[poppurple!14] ([xshift=2.4cm,yshift=7.6cm]current page.south west) circle (3.8cm);
  \end{tikzpicture}%
  {\poplogo\fontsize{30}{30}\selectfont\color{popink}OnBuch\color{poporange}+}\hfill
  \raisebox{0.6ex}{\poppill{Cours signé \textbullet\ Premium}{popink}{popcream}}\par
  \vspace{1pt}\popsquiggle{poppurple}\par
  \vspace{8pt}
  \begin{tcolorbox}[enhanced,colback=poporange,colframe=popink,boxrule=2.8pt,arc=13pt,
      drop shadow=popink,left=18pt,right=18pt,top=12pt,bottom=13pt,after skip=10pt]
    \poppill{#2 \textbullet\ #3}{popink}{popcream}\hspace{5pt}\poppill{#4}{white}{popink}\par\vspace{8pt}
    {\popdisplay\fontsize{14}{14}\selectfont\color{popink}LEÇON #5}\par\vspace{4pt}
    {\raggedright\hyphenpenalty=10000\exhyphenpenalty=10000\popdisplay\fontsize{25}{27}\selectfont\color{popcream}\MakeUppercase{#1}\par}
    \vspace{8pt}
    {\popxbold\fontsize{9.5}{9.5}\selectfont\color{popink}\MakeUppercase{Durée conseillée : #6}}
  \end{tcolorbox}
  \begin{tcolorbox}[enhanced,colback=white,colframe=popink,boxrule=2.2pt,arc=10pt,
      drop shadow=popink,left=16pt,right=16pt,top=10pt,bottom=10pt,after skip=8pt]
    \poppill{Dans cette leçon}{poppurple}{white}\par\vspace{5pt}
    \popsemi\fontsize{9.3}{12.6}\selectfont\color{popink}#7
  \end{tcolorbox}
  \noindent\begin{minipage}[t]{0.487\linewidth}
    \begin{tcolorbox}[enhanced,colback=popgreen,colframe=popink,boxrule=2.2pt,arc=10pt,
        drop shadow=popink,left=11pt,right=11pt,top=10pt,bottom=10pt,height=3.1cm,valign=center]
      \tcbox[colback=white,colframe=popink,boxrule=1.3pt,arc=5pt,boxsep=0pt,nobeforeafter,
        left=3pt,right=3pt,top=3pt,bottom=3pt,valign=center]{\qrcode[height=1.9cm]{https://play.google.com/store/apps/details?id=com.luvvix.onbuch&hl=fr}}%
      \hspace{8pt}\begin{minipage}[c]{0.5\linewidth}
        {\popdisplay\fontsize{11}{12.5}\selectfont\color{white}\MakeUppercase{Télécharge OnBuch}}\par\vspace{4pt}
        {\raggedright\popsemi\fontsize{8.4}{11}\selectfont\color{white}Gratuit \textbullet\ Google Play\\Tuteur IA, annales, concours, résultats\par}
      \end{minipage}
    \end{tcolorbox}
  \end{minipage}\hfill
  \begin{minipage}[t]{0.487\linewidth}
    \begin{tcolorbox}[enhanced,colback=poppurple,colframe=popink,boxrule=2.2pt,arc=10pt,
        drop shadow=popink,left=11pt,right=11pt,top=10pt,bottom=10pt,height=3.1cm,valign=center]
      \tikz[baseline=-0.7ex]\node[circle,fill=white,draw=popink,line width=1.3pt,minimum size=1.6cm,inner sep=0]{\popdisplay\fontsize{24}{24}\selectfont\color{poppurple}+};%
      \hspace{8pt}\begin{minipage}[c]{0.6\linewidth}
        {\popdisplay\fontsize{11}{12.5}\selectfont\color{white}\MakeUppercase{Passe à OnBuch+}}\par\vspace{4pt}
        {\raggedright\popsemi\fontsize{8.4}{11}\selectfont\color{white}Tous les cours signés, Tuteur IA illimité, fiches PDF — onglet OnBuch+ de l'app\par}
      \end{minipage}
    \end{tcolorbox}
  \end{minipage}
  \vspace{8pt}
  \popcontenu
  \vfill
  \signatureonbuch
  \restoregeometry
  \newpage
}

% Rangée « Ce cours contient » (page de garde)
\newcommand{\poptile}[3]{% #1 couleur #2 gros texte #3 légende
  \begin{tcolorbox}[enhanced,colback=#1,colframe=popink,boxrule=2pt,arc=9pt,shadow={2.5pt}{-2.5pt}{0pt}{popink},
      left=6pt,right=6pt,top=9pt,bottom=9pt,halign=center,valign=center,height=1.75cm,width=\linewidth,nobeforeafter]
    {\popdisplay\fontsize{12.5}{13}\selectfont\color{white}\MakeUppercase{#2}}\par\vspace{3pt}
    {\centering\popsemi\fontsize{7.8}{9.5}\selectfont\color{white}#3\par}
  \end{tcolorbox}}
\newcommand{\popcontenu}{%
  \par\noindent{\popxboldls\fontsize{7.5}{7.5}\selectfont\color{popmuted}CE COURS SIGNÉ CONTIENT}\par\vspace{7pt}
  \noindent\begin{minipage}[t]{0.235\linewidth}\poptile{popblue}{Cours}{complet, illustré, pas à pas}\end{minipage}\hfill
  \begin{minipage}[t]{0.235\linewidth}\poptile{poporange}{Méthodes}{et exercices résolus}\end{minipage}\hfill
  \begin{minipage}[t]{0.235\linewidth}\poptile{poppink}{Exercices}{type examen + corrigés}\end{minipage}\hfill
  \begin{minipage}[t]{0.235\linewidth}\poptile{popgreen}{Fiche bilan}{l'essentiel à retenir}\end{minipage}\par}

% Sommaire
\newtcolorbox{sommaire}{enhanced,colback=white,colframe=popink,boxrule=2.2pt,arc=10pt,
  drop shadow=popink,left=18pt,right=18pt,top=13pt,bottom=13pt,before skip=6pt,after skip=20pt,
  before upper={\poppill{Sommaire}{poppurple}{white}\par\vspace{9pt}\popsemi\fontsize{10.6}{14.5}\selectfont\color{popink}}}

% Signature de fin de cours — poussée tout en bas de la dernière page
\newcommand{\findecours}{%
  \par\vfill\nopagebreak
  \begin{center}\popsquiggle{poporange}\end{center}\vspace{4pt}
  \signatureonbuch
}
\AtBeginDocument{\def\llbracket{\mathopen{[\mkern-3.5mu[}}\def\rrbracket{\mathclose{]\mkern-3.5mu]}}} % intervalles d'entiers (glyphe ⟦ absent de la police)
% Glyphes absents de Fira Math : redéfinis à partir de glyphes disponibles
\AtBeginDocument{%
  \def\setminus{\mathbin{\backslash}}\let\smallsetminus\setminus
  \def\vdots{\mathord{\vbox{\baselineskip3.2pt\lineskiplimit0pt\kern4pt\hbox{.}\hbox{.}\hbox{.}}}}%
  \def\ddots{\mathinner{\mkern1mu\raise6.5pt\hbox{.}\mkern2mu\raise3.6pt\hbox{.}\mkern2mu\raise0.7pt\hbox{.}\mkern1mu}}%
  \def\triangle{\mathord{\scalebox{0.9}{$\symup{\Delta}$}}}%
  \def\nabla{\mathord{\raisebox{\depth}{\scalebox{1}[-1]{$\symup{\Delta}$}}}}%
  \def\checkmark{\popcheck{popdark}}%
  \def\star{\mathbin{\ast}}\def\bigstar{\mathord{\ast}}%
  \def\diamond{\mathbin{\scalebox{0.6}{$\Diamond$}}}%
  \def\bigcup{\mathop{\mathchoice{\raisebox{-0.3em}{\scalebox{1.7}{$\cup$}}}{\scalebox{1.25}{$\cup$}}{\cup}{\cup}}\displaylimits}%
  \def\bigcap{\mathop{\mathchoice{\raisebox{-0.3em}{\scalebox{1.7}{$\cap$}}}{\scalebox{1.25}{$\cap$}}{\cap}{\cap}}\displaylimits}%
  \def\lesseqqgtr{\mathrel{\lessgtr}}\def\bigoplus{\mathop{\oplus}\displaylimits}%
}
\providecommand{\ssi}{si et seulement si} % abréviation fréquente
\AtBeginDocument{\providecommand{\wideparen}{\overparen}} % arc de cercle

%% FIGFIX-BEGIN v4 — correctifs globaux des figures (docs/onbuch-generation/tools/figfix)
\usepackage{adjustbox}\usepackage{etoolbox}
% toute figure TikZ/pgfplots plus large que la ligne est réduite à la largeur disponible
\BeforeBeginEnvironment{tikzpicture}{\begin{adjustbox}{max width=\linewidth}}
\AfterEndEnvironment{tikzpicture}{\end{adjustbox}}
% légendes pgfplots lisibles par-dessus les courbes
\pgfplotsset{every axis/.append style={legend style={fill=white,fill opacity=0.92,text opacity=1,draw=black!15,inner sep=3pt}}}
% étiquettes d'arêtes (syntaxe edge["..."]) avec fond blanc
\tikzset{every edge quotes/.append style={fill=white,inner sep=1.2pt}}
%% FIGFIX-END
\begin{document}

\pagegarde{Le Droit et la Justice}{Philosophie}{1re Tech}{Philosophie – classe de Première technique}{N08}{4 séances (fiche de progression harmonisée)}{Cette leçon interroge les notions de droit et de justice : leurs définitions, les fondements du droit (naturel ou positif), le rapport entre justice et égalité, et le passage des formes concrètes de justice à l'idée absolue de justice. Elle s'appuie sur des situations concrètes de la vie juridique et sociale camerounaise.
\vspace{4pt}
\begin{multicols}{2}\small
\begin{itemize}
\item À la fin de cette leçon, je sais définir les concepts de droit et de justice et les distinguer clairement
\item À la fin de cette leçon, je sais distinguer le droit objectif des droits subjectifs à partir d'exemples camerounais
\item À la fin de cette leçon, je sais comparer le droit naturel et le droit positif et identifier leurs fondements
\item À la fin de cette leçon, je sais analyser une situation concrète (procès, sanction, conflit foncier) à l'aide des notions de la leçon
\item À la fin de cette leçon, je sais expliquer la distinction entre justice distributive et justice commutative
\item À la fin de cette leçon, je sais discuter le rapport entre justice et égalité (égalité arithmétique / égalité proportionnelle)
\end{itemize}\end{multicols}}

\begin{sommaire}
\begin{multicols}{2}
\begin{enumerate}
\item Définition des concepts de droit et de justice
\item Les fondements du droit
\item Justice et égalité
\item Les formes de justice
\item De la justice relative à l'idée absolue de justice
\item Méthode et application aux épreuves
\item Activité d'intégration
\item Méthodes et exercices résolus
\item Exercices
\item Corrigés des exercices
\item Fiche bilan
\end{enumerate}
\end{multicols}
\end{sommaire}

\begin{objectifs}
\begin{itemize}
\item À la fin de cette leçon, je sais définir les concepts de droit et de justice et les distinguer clairement
\item À la fin de cette leçon, je sais distinguer le droit objectif des droits subjectifs à partir d'exemples camerounais
\item À la fin de cette leçon, je sais comparer le droit naturel et le droit positif et identifier leurs fondements
\item À la fin de cette leçon, je sais analyser une situation concrète (procès, sanction, conflit foncier) à l'aide des notions de la leçon
\item À la fin de cette leçon, je sais expliquer la distinction entre justice distributive et justice commutative
\item À la fin de cette leçon, je sais discuter le rapport entre justice et égalité (égalité arithmétique / égalité proportionnelle)
\item À la fin de cette leçon, je sais construire un schéma ou un tableau comparatif synthétisant les formes de justice
\item À la fin de cette leçon, je sais rédiger et justifier une réponse argumentée à une question de dissertation ou d'explication de texte
\end{itemize}
\end{objectifs}

\begin{prerequis}
\begin{itemize}
\item Notion de règle et de sanction (vie de classe, règlement intérieur)
\item Distinction entre morale et coutume (notions abordées en classe de Seconde)
\item Notions de liberté et d'État (leçons antérieures de Première)
\item Connaissance élémentaire des institutions camerounaises (tribunaux, Constitution de 1996)
\end{itemize}
\end{prerequis}

\newpage

\coursec{Définition des concepts de droit et de justice}

\courssub{Situation d'accroche et problématique}

Imaginez une scène banale à un carrefour de Yaoundé, un matin de semaine. Un jeune motocycliste, pressé d'aller au travail, est intercepté par un agent de police. Il a oublié sa pièce d'identité à la maison. L'agent dresse un procès-verbal et exige le paiement immédiat d'une amende de 10\,000 FCFA, sans reçu officiel. Le conducteur s'exécute, la rage au cœur : « C'est légal, dit-il, mais ce n'est pas juste ! »

Quelques centaines de kilomètres plus au nord, dans un village de l'Extrême-Nord, deux familles se disputent un champ hérité de leurs ancêtres. Aucun titre foncier n'existe. Plutôt que de saisir le tribunal de grande instance de Maroua, les parties vont voir le Lamido. Après une séance de palabre sous l'arbre à palabres, le chef tranche selon la coutume : le champ revient à la famille qui l'exploite depuis trois générations, l'autre recevant une compensation en bétail. Les deux parties acceptent la décision et partagent le kola.

Ces deux situations soulèvent le cœur de notre leçon : \textbf{Pourquoi certains conflits relèvent-ils du tribunal d'État et d'autres de la coutume ? Le droit est-il toujours juste, et la justice exige-t-elle de traiter tout le monde de la même manière ?} Pour y répondre, nous devons d'abord distinguer avec rigueur les concepts de \cle{droit} et de \cle{justice}, souvent confondus dans le langage courant.

\courssub{Sens étymologique et approche initiale}

\begin{methode}[Méthode pour définir une notion philosophique]
\begin{enumerate}
\item \textbf{Étymologie :} remonter au sens originel du mot (latin, grec) pour en saisir l'intuition première.
\item \textbf{Sens commun :} recenser les usages quotidiens, souvent flous ou métaphoriques.
\item \textbf{Définition conceptuelle :} proposer une définition rigoureuse, opératoire, qui distingue la notion de ses voisines.
\item \textbf{Exemple concret :} illustrer par un fait social, un texte de loi ou une situation vécue (ici, au Cameroun).
\end{enumerate}
\end{methode}

Appliquons cette méthode.

\begin{itemize}
\item \textbf{Droit} : du latin \emph{directum} (participe passé de \emph{dirigere}), « ce qui est droit, ce qui est réglé, ce qui va en ligne droite ». L'idée première est celle de la \cle{rectitude}, de la \cle{règle} qui redresse le comportement humain. En français, le mot est polysémique : « le Droit » (l'ensemble des règles) vs « un droit » (une prérogative).
\item \textbf{Justice} : du latin \emph{justitia}, dérivé de \emph{jus} (le droit, la loi) et de la racine \emph{justus} (juste, réglé). Étymologiquement, la justice est l'art de \cle{rendre à chacun ce qui lui est dû} (\emph{suum cuique tribuere}, formule d'Ulpien, juriste romain). Elle implique une idée d'\cle{équilibre}, de \cle{proportion} et de \cle{reconnaissance}.
\end{itemize}

\begin{aretenir}[Intuition fondatrice]
Le \textbf{Droit} est l'ensemble des \textbf{règles} qui organisent la vie en société (l'ordre). La \textbf{Justice} est l'idéal de justesse dans l'application de ces règles ou dans les rapports entre les hommes (l'équité). Le droit est un \emph{fait} (des textes, des institutions) ; la justice est une \emph{valeur} (un horizon normatif).
\end{aretenir}

\courssub{Le droit objectif : l'ordre juridique}

Dans sa première acception, le droit n'est pas un « avoir », c'est un « devoir-être ». C'est le \cle{Droit objectif}.

\begin{definition}[Droit objectif]
Ensemble des \textbf{règles juridiques} (lois, décrets, coutumes, traités) qui s'imposent à tous les membres d'une société, sanctionnées par la \textbf{puissance publique} (État, tribunaux, forces de l'ordre). Il est \emph{objectif} car il existe indépendamment de la volonté individuelle : on le \emph{subit} avant de pouvoir s'en prévaloir.
\end{definition}

\textbf{Caractères des règles de droit objectif (selon la doctrine juridique) :}
\begin{enumerate}
\item \textbf{Généralité :} elles s'adressent à des catégories de personnes (ex. « tout citoyen », « tout conducteur »), pas à un individu nommé.
\item \textbf{Abstraction :} elles régissent des situations types (ex. « tout vol »), pas un fait unique.
\item \textbf{Obligativité :} elles commandent, interdisent ou autorisent de manière impérative.
\item \textbf{Sanction étatique :} leur violation entraîne une réaction organisée de la puissance publique (amende, prison, nullité d'acte, dommages-intérêts).
\end{enumerate}

\begin{exemplebox}[Exemples camerounais de droit objectif]
\begin{itemize}
\item Le \textbf{Code pénal} (Loi n°2016/007 du 12 juillet 2016) : définit les infractions (vol, escroquerie, homicide) et fixe les peines.
\item Le \textbf{Code civil} (inspiré du Code Napoléon, adapté) : régit les contrats, la filiation, les successions, la propriété.
\item Le \textbf{Code du travail} (Loi n°92/007) : encadre les relations employeur/salarié (contrat, licenciement, syndicats).
\item La \textbf{Constitution du 18 janvier 1996} (révisée en 2008) : norme suprême, source de toute légalité.
\item La \textbf{coutume} : reconnue par l'article 1er de la loi n°2006/015 du 29 décembre 2006 portant organisation judiciaire, pourvu qu'elle ne soit pas contraire à l'ordre public (ex. droit coutumier foncier, successoral dans certaines régions).
\end{itemize}
\end{exemplebox}

\begin{popfigure}
\begin{tikzpicture}[
    box/.style={draw=popdark, thick, rounded corners=6pt, fill=popcream, minimum height=1.8cm, text width=0.42\linewidth, align=center, font=\small},
    title/.style={font=\bfseries\large, text=popblue, fill=popblue!10, rounded corners=4pt, inner sep=4pt},
    arrow/.style={->, >=Stealth, thick, poporange, bend left=15},
    ex/.style={font=\footnotesize, text=popmuted, align=left, text width=0.38\linewidth}
]
% Colonne Gauche : Droit Objectif
\node[box, align=center] (obj) at (0,0){
    {\bfseries\color{popblue}DROIT OBJECTIF}\\[4pt]
    \textbf{Structure :} Ensemble de règles impersonnelles.\\
    \textbf{Source :} Loi, Coutume, Traité, Jurisprudence.\\
    \textbf{Sanction :} Contrainte publique (État).
};
\node[ex, anchor=north west, align=center] at (obj.south west){
    $\bullet$ Code pénal camerounais (2016)\\
    $\bullet$ Constitution 1996 (Art. 1er : « Le Cameroun est une République décentralisée... »)\\
    $\bullet$ Coutume Bamiléké (héritage patrilinéaire)\\
    $\bullet$ Traité CEMAC / OHADA
};
% Colonne Droite : Droits Subjectifs
\node[box, align=center] (subj) at (8.5,0){
    {\bfseries\color{popblue}DROITS SUBJECTIFS}\\[4pt]
    \textbf{Structure :} Prérogatives / Pouvoirs individuels.\\
    \textbf{Titulaire :} Personne physique ou morale.\\
    \textbf{Protection :} Action en justice (tribunaux).
};
\node[ex, anchor=north west, align=center] at (subj.south west){
    $\bullet$ Droit de vote (Art. 2 Const.)\\
    $\bullet$ Droit de propriété (Préambule Const. + DDH 1789)\\
    $\bullet$ Droit à l'éducation (Art. 25 Const.)\\
    $\bullet$ Droit au salaire (Code du travail)
};
% Flèche de lien
\draw[arrow] (obj.east) to node[midway, above, font=\small\bfseries, text=poppurple] {FONDE / GÉNÈRE} (subj.west);
% Titre global
\node[font=\bfseries\Large, text=popdark, anchor=south] at (4.25, 3.2) {La distinction fondamentale : Droit Objectif vs Droits Subjectifs};
\end{tikzpicture}
\legende{Schéma conceptuel : le Droit objectif (l'ordre normatif) est la source qui fonde et délimite les droits subjectifs (les prérogatives individuelles). Exemples tirés du contexte juridique camerounais.}
\end{popfigure}

\begin{attention}[Erreur fréquente : confondre « avoir le droit de » et « le Droit »]
\begin{itemize}
\item \textbf{« J'ai le droit de voter »} $\rightarrow$ \cle{Droit subjectif} (prérogative, pouvoir, faculté reconnue à \emph{moi}).
\item \textbf{« Le Droit électoral impose la majorité à 20 ans »} $\rightarrow$ \cle{Droit objectif} (règle générale, impersonnelle, s'imposant à \emph{tous}).
\item \textbf{Piège à l'examen :} Si le sujet porte sur « Le Droit », ne parlez pas seulement de vos droits personnels. Si le sujet porte sur « Les droits de l'homme », analysez des prérogatives subjectives garanties par un droit objectif (Constitution, traités).
\end{itemize}
\end{attention}

\courssub{Les droits subjectifs : les prérogatives de la personne}

Le droit objectif ne flotte pas dans l'air : il \textbf{attribue} des pouvoirs aux sujets de droit (personnes physiques : citoyens, résidents ; personnes morales : État, entreprises, associations).

\begin{definition}[Droit subjectif]
Prérogative, pouvoir ou faculté reconnue par le \textbf{droit objectif} à un \textbf{titulaire} (personne physique ou morale), lui permettant d'exiger un comportement d'autrui (droit de créance), de disposer d'un bien (droit réel) ou d'agir librement dans une sphère protégée (droit de liberté).
\end{definition}

On distingue classiquement trois grandes catégories (classification de la doctrine romaniste) :

\begin{center}
\begin{tabularx}{\linewidth}{|l|X|X|}\hline
\rowcolor{popblueL}
\textbf{Type de droit subjectif} & \textbf{Objet / Contenu} & \textbf{Exemple camerounais} \\ \hline
\textbf{Droits réels} (sur une chose) & Pouvoir direct et immédiat sur un bien (propriété, usufruit, servitude). & Titre foncier n°12345 à Douala : droit de propriété sur une parcelle. \\ \hline
\textbf{Droits de créance} (contre une personne) & Droit d'exiger une prestation (donner, faire, ne pas faire) d'un débiteur. & Contrat de travail : le salarié a une créance de salaire contre l'employeur. \\ \hline
\textbf{Droits de liberté / personnalité} (attachés à la personne) & Sphère d'autonomie protégée (vie, intégrité, opinion, vote, vie privée). & Art. 2 Const. : « Droit de vote ». Préambule Const. : « Propriété inviolable et sacrée ». \\ \hline
\end{tabularx}
\end{center}

\begin{exemplebox}[Document d'étude : Extrait de la Constitution du Cameroun (1996, révisée 2008)]
\textbf{Préambule :} « La personne humaine possède des droits inaliénables et sacrés. [...] Tous les citoyens ont les mêmes droits et les mêmes devoirs. [...] La propriété étant un droit inviolable et sacré, nul ne peut en être privé, si ce n'est pour cause d'utilité publique, légalement constatée, et sous condition d'une juste et préalable indemnité. »
\textbf{Article 1er, al. 2 :} « Le Cameroun est une République décentralisée. [...] La souveraineté nationale appartient au peuple qui l'exerce soit par référendum, soit par l'intermédiaire de ses représentants élus. »
\textbf{Article 2 :} « Le suffrage est universel, égal et secret. Tout citoyen des deux sexes, âgé de vingt ans révolus, jouissant de ses droits civiques, est électeur. » \tcblower \textbf{Analyse :} Ce texte fonde des \textbf{droits subjectifs} (vote, propriété, égalité) en les inscrivant dans le \textbf{droit objectif} suprême (la Constitution). Le citoyen camerounais devient \emph{titulaire} de ces prérogatives qu'il peut opposer à l'État lui-même (recours devant le Conseil constitutionnel ou la Cour suprême).
\end{exemplebox}

\courssub{La justice : vertu et institution}

Si le droit est l'outil (la règle), la justice est à la fois le \textbf{but} visé et le \textbf{moteur} de l'action. Elle se décline en deux sens indissociables.

\begin{definition}[Justice : double acception]
\begin{enumerate}
\item \textbf{Justice comme vertu (éthique) :} Disposition stable de l'âme (Aristote) ou de la volonté (Kant) qui porte à \textbf{rendre à chacun ce qui lui est dû} (\emph{suum cuique}). L'homme juste respecte la loi, mais aussi l'équité (\emph{epikeia}) quand la loi est trop générale pour le cas particulier. C'est une \cle{qualité morale}.
\item \textbf{Justice comme institution (politique/juridique) :} Appareil d'État chargé de \textbf{dire le droit} et de \textbf{faire appliquer la loi}. Au Cameroun : Cour Suprême, Cours d'Appel, Tribunaux de Grande Instance, Tribunaux de Première Instance, Tribunaux Coutumiers, Conseil Constitutionnel. Composée de \textbf{magistrats} (juges du siège et parquet) et d'auxiliaires (greffiers, huissiers, avocats). C'est une \cle{structure organisée}.
\end{enumerate}
\end{definition}

\begin{popfigure}
\begin{tikzpicture}[
    building/.style={draw=popdark, thick, fill=popblue!5},
    pillar/.style={draw=popdark, thick, fill=popblue!20},
    roof/.style={draw=popdark, thick, fill=poporange!60},
    lbl/.style={font=\small, text=popdark, align=center},
    title/.style={font=\bfseries\large, text=popblue}
]
% Base
\fill[popline!30] (-3.5,-0.5) rectangle (3.5,0);
% Escaliers
\foreach \y in {0,0.3,0.6} {
    \fill[popmuted!50] (-2.5,\y) rectangle (2.5,\y+0.15);
}
% Colonnes
\foreach \x in {-2.2, -0.7, 0.7, 2.2} {
    \node[pillar, minimum width=0.7cm, minimum height=3.5cm, anchor=south] at (\x,0.6) {};
}
% Fronton / Toit
\draw[roof] (-3,4.1) -- (0,5.5) -- (3,4.1) -- cycle;
% Entablement
\fill[popblue!30] (-3.2,4.1) rectangle (3.2,4.4);
% Inscription
\node[lbl, text=popblue] at (0,4.25) {REPUBLIQUE DU CAMEROUN};
\node[lbl, text=popdark] at (0,3.8) {COUR SUPRÊME};
% Porte
\draw[building, fill=popblue!30] (-1,0.6) rectangle (1,2.2);
\node[lbl, text=popdark] at (0,1.4) {ENTRÉE};
% Fenêtres
\foreach \x in {-2.8, -1.5, 1.5, 2.8} {
    \foreach \y in {1.5, 2.5} {
        \draw[building, fill=popblue!20] (\x-0.3,\y) rectangle (\x+0.3,\y+0.7);
    }
}
% Drapeau
\draw[thick, popdark] (3.2,5.5) -- (4.5,6.5);
\fill[popgreen] (3.2,5.5) rectangle (4.2,6.0);
\fill[poppink] (3.2,5.5) rectangle (4.2,5.83);
\fill[poppink] (3.2,5.17) rectangle (4.2,5.5);
\fill[poppink] (3.2,4.83) rectangle (4.2,5.17);
\fill[poporange] (3.2,4.83) rectangle (4.2,4.5);
\fill[popgold] (3.7,5.33) circle (0.15);
% Légende interne
\node[title, anchor=south] at (0,6.5) {Façade schématique de la Cour Suprême -- Yaoundé};
\node[font=\footnotesize, text=popmuted, anchor=north] at (0,-1.2) {Siège de la plus haute juridiction de l'ordre judiciaire, administrative et des comptes. Symbole de la Justice institutionnelle.};
\end{tikzpicture}
\legende{La Cour Suprême du Cameroun (Yaoundé) : incarnation architecturale de la Justice institutionnelle. Au sommet de la pyramide juridictionnelle, elle assure l'unité de la jurisprudence et le contrôle de la légalité administrative.}
\end{popfigure}

\begin{savaistu}[Le saviez-vous ? : Dualité juridictionnelle au Cameroun]
Le Cameroun connaît une \textbf{double organisation juridictionnelle} (héritage franco-anglais) :
\begin{itemize}
\item \textbf{Ordre judiciaire :} Tribunaux de Première Instance (TPI) $\rightarrow$ Tribunaux de Grande Instance (TGI) $\rightarrow$ Cours d'Appel $\rightarrow$ \textbf{Cour Suprême} (Chambre Judiciaire). Juge les litiges entre particuliers et les infractions pénales.
\item \textbf{Ordre administratif :} Tribunaux Administratifs $\rightarrow$ Cours d'Appel Administratif $\rightarrow$ \textbf{Cour Suprême} (Chambre Administrative). Juge les litiges impliquant l'État ou les collectivités (excès de pouvoir, responsabilité).
\item \textbf{Juridictions spécialisées :} Tribunal Militaire, Cour des Comptes (chambre des comptes à la Cour Suprême), Tribunaux Coutumiers (1er degré pour la coutume).
\end{itemize}
Cette dualité explique pourquoi un conflit foncier peut aller au Tribunal Administratif (si l'État a exproprié) ou au Tribunal Coutumier (si c'est un litige entre familles sur terre ancestrale).
\end{savaistu}

\courssub{Distinction cruciale : le légal et le légitime}

C'est le cœur philosophique de cette section. La situation du motocycliste yaoundéen l'illustre parfaitement.

\begin{propriete}[Légal vs Légitime]
\begin{itemize}
\item \textbf{Le Légal (Legalitas) :} Ce qui est \textbf{conforme à la loi écrite}, à la procédure, au texte en vigueur. C'est la validité \emph{formelle}. L'amende du motocycliste est \emph{légale} \textbf{SI} : le Code de la route prévoit une sanction pour défaut de papiers, l'agent est assermenté, la procédure verbale est respectée.
\item \textbf{Le Légitime (Legitimitas) :} Ce qui est \textbf{conforme à la justice, à l'équité, à la morale, au consentement}. C'est la validité \emph{matérielle} ou \emph{axiologique}. L'amende est \emph{illégitime} \textbf{SI} : l'agent exige un bakchich sans reçu (corruption), le montant est disproportionné, le motif est un prétexte pour racketter, le conducteur n'avait pas moyen matériel de présenter le papier sur l'instant.
\end{itemize}
\end{propriete}

\begin{exoresolu}[Analyse de la situation d'accroche : L'amende du motocycliste]
\textbf{Énoncé :} Un motocycliste arrêté sans pièce d'identité paie 10\,000 FCFA à un agent sans reçu. Il dit : « C'est légal mais pas juste. » Analysez cette affirmation.

\textbf{Solution détaillée :}
\begin{enumerate}
\item \textbf{Qualification juridique (Droit objectif) :} L'article 145 de la Loi n°2016/008 (Code de la route) sanctionne la non-présentation des documents. L'amende forfaitaire existe. \textbf{Donc l'infraction est légale (prévue par la loi).}
\item \textbf{Qualification de la procédure :} Le paiement direct à l'agent sans procès-verbal ni reçu officiel viole la procédure (Code de procédure pénale, Loi n°2005/007 du 27 juillet 2005, règlementation sur la perception des amendes). \textbf{Donc l'acte de l'agent est illégal (irrégulier).}
\item \textbf{Qualification morale (Justice) :} L'agent use de sa qualité pour extorquer de l'argent (corruption passive/active). Le conducteur paie sous la contrainte (peur de la garde à vue, perte de temps). \textbf{Donc la situation est illégitime (injuste, immorale).}
\item \textbf{Synthèse :} Le conducteur a raison sur le fond : il y a \textbf{légalité de la règle} (le texte existe) mais \textbf{illégalité de l'acte} (procédure violée) et surtout \textbf{illégitimité totale} (détournement de pouvoir, corruption). La \cle{justice} exige que la sanction soit prononcée par un juge, proportionnée, et que l'argent aille au Trésor public, pas dans la poche de l'agent.
\end{enumerate}
\textbf{Conclusion :} \emph{Tout ce qui est légal n'est pas nécessairement légitime} (lois injustes, application corruptrice). \emph{Tout ce qui est légitime n'est pas toujours légal} (désobéissance civile, résistance à l'oppression, coutume non codifiée mais acceptée comme juste -- cas du village de l'Extrême-Nord).
\end{exoresolu}

\begin{aretenir}[Synthèse conceptuelle : Droit, Justice, Légalité, Légitimité]
\begin{center}
\begin{tabular}{|c|c|c|}\hline
\rowcolor{popblueL}
\textbf{Concept} & \textbf{Nature} & \textbf{Question clé} \\ \hline
\textbf{Droit objectif} & Ensemble de règles (Fait social) & Quelle est la règle applicable ? \\ \hline
\textbf{Droit subjectif} & Prérogative individuelle (Pouvoir) & Qu'est-ce que \emph{je} peux exiger ? \\ \hline
\textbf{Justice (Vertu)} & Qualité morale (Valeur) & Suis-je juste dans mon jugement ? \\ \hline
\textbf{Justice (Institution)} & Appareil judiciaire (Structure) & Quel tribunal tranche ? \\ \hline
\textbf{Légal} & Conformité à la loi (Formel) & Est-ce que la procédure est respectée ? \\ \hline
\textbf{Légitime} & Conformité à la justice (Matériel) & Est-ce que la décision est équitable ? \\ \hline
\end{tabular}
\end{center}
\end{aretenir}

\courssub{Exercices d'application}

\begin{exercice}[Exercice 1 : Qualification juridique]
Pour chacune des situations ci-dessous, précisez s'il s'agit d'une manifestation du \textbf{Droit objectif}, d'un \textbf{Droit subjectif}, de la \textbf{Justice comme vertu} ou de la \textbf{Justice comme institution}. Justifiez en une phrase.
\begin{enumerate}
\item Le Code OHADA uniforme sur le droit des sociétés.
\item M. Fouda exige le paiement de son loyer auprès de son locataire.
\item Le juge du TPI de Ngaoundéré rend un jugement dans un litige successoral.
\item Une mère partage équitablement la part de poulet entre ses trois enfants.
\item La Constitution camerounaise garantit la liberté d'expression.
\item Le Lamido de Rey Bouba tranche un conflit de limite entre deux chefferies.
\end{enumerate}
\end{exercice}

\begin{corrige}[Exercice 1]
\begin{enumerate}
\item \textbf{Droit objectif} : Règle impersonnelle, générale, sanctionnée (traité international intégré).
\item \textbf{Droit subjectif (de créance)} : Prérogative de M. Fouda (créancier) contre son locataire (débiteur) née du contrat de bail.
\item \textbf{Justice comme institution} : Le juge est un organe de l'appareil judiciaire étatique qui « dit le droit ».
\item \textbf{Justice comme vertu} : La mère agit par équité (vertu de justice distributive), sans contrainte légale.
\item \textbf{Droit objectif (source) / Droit subjectif (résultat)} : La Constitution (règle suprême) fonde le droit subjectif de liberté d'expression pour chaque citoyen.
\item \textbf{Justice comme institution (coutumière)} : Le Lamido exerce une fonction juridictionnelle reconnue par l'État (loi sur l'organisation judiciaire) au nom de la coutume.
\end{enumerate}
\end{corrige}

\begin{exercice}[Exercice 2 : Légal vs Légitime -- Dissertation guidée]
Sujet : « \emph{Une loi injuste n'est pas une loi.} » (Saint Augustin). Discutez cette affirmation en vous appuyant sur la distinction entre légalité et légitimité et sur des exemples camerounais.
\textbf{Consignes :} Utilisez la méthode du plan dialectique (Thèse / Antithèse / Synthèse). Définissez les termes. Donnez au moins deux exemples concrets (un historique, un actuel).
\end{exercice}

\begin{corrige}[Exercice 2 -- Éléments de corrigé]
\textbf{Introduction :} Accroche (ex. lois raciales, lois d'apartheid, décret d'exception). Définition : Loi = règle juridique générale (légal) ; Juste = conforme à l'équité/droit naturel (légitime). Problématique : La validité formelle (légalité) suffit-elle à fonder l'obligation morale d'obéir ? Annonce du plan.
\textbf{Thèse :} La loi injuste n'est pas une loi (thèse du droit naturel / jusnaturalisme). Saint Augustin, Thomas d'Aquin, Locke, Déclaration des Droits de l'Homme 1789 (Art. 2 : résistance à l'oppression). Exemple : Lois coloniales d'indigénat (travail forcé, code de l'indigénat) -- légales formellement, illégitimes moralement. Désobéissance civile justifiée (mouvement nationaliste UPC, grève de 1945 à Douala).
\textbf{Antithèse :} La loi injuste reste une loi (thèse du positivisme juridique / Kelsen, Hart). Séparation stricte Droit / Morale. La loi tire sa validité de sa source (vote parlement, promulgation), pas de son contenu. Exemple : Code pénal actuel -- certaines peines jugées dures (peine de mort maintenue dans les textes bien que non appliquée depuis 1997) restent légales tant qu'elles ne sont pas abrogées. Stabilité de l'ordre juridique : si chacun juge de la justice de la loi, c'est l'anarchie.
\textbf{Synthèse :} Dépassement : La légalité est une condition nécessaire mais non suffisante de la légitimité. Le contrôle de constitutionnalité (Conseil Constitutionnel, opérationnel depuis 2018) et le contrôle de conventionnalité (traité internationaux, Charte africaine des droits de l'homme et des peuples) permettent d'écarter les lois illégitimes. Au Cameroun, l'art. 45 Const. : « Les traités régulièrement ratifiés ont une autorité supérieure à celle des lois. » Permet de sanctionner une loi nationale contraire aux droits de l'homme (légitimité supérieure).
\textbf{Conclusion :} La maxime d'Augustin exprime l'exigence éthique du droit. Le droit positif moderne intègre cette exigence par le bloc de constitutionnalité et le contrôle juridictionnel.
\end{corrige}

\coursec{Les fondements du droit}

Après avoir défini le droit comme un ensemble de normes obligatoires régissant la vie en société et la justice comme l'idéal de justesse et d'équité vers lequel le droit doit tendre, nous devons maintenant nous interroger sur l'origine de la force obligatoire du droit. Pourquoi obéit-on à la loi ? D'où tient-elle sa légitimité ? Ce questionnement nous conduit au cœur du débat philosophique sur les fondements du droit : le droit est-il inscrit dans la nature même de l'homme (droit naturel) ou est-il le fruit d'une construction humaine, historique et conventionnelle (droit positif) ?

\courssub{Le problème : nature ou convention}

L'opposition entre \emph{physis} (la nature) et \emph{nomos} (la loi, la coutume, la convention) structure la philosophie antique. Pour les sophistes, la loi est une contrainte artificielle qui brime les forts au profit des faibles ; pour Socrate, Platon et Aristote, la vraie loi découle d'un ordre naturel ou cosmique que la raison humaine peut découvrir.

\begin{aretenir}[La question fondamentale]
Le fondement du droit désigne le principe ultime qui confère à la norme juridique sa validité et son autorité. Deux réponses majeures s'affrontent :
\begin{itemize}
    \item \textbf{Le jusnaturalisme :} le droit précède l'État ; il existe des droits « antérieurs et supérieurs » aux lois positives (droit naturel).
    \item \textbf{Le positivisme juridique :} le droit est la loi posée par le souverain ; il n'y a de droit que le droit positif, créé par l'histoire et la volonté humaine.
\end{itemize}
\end{aretenir}

\courssub{Le droit naturel : universalité et imprescriptibilité}

\emph{Intuition :} Face à une loi inique (ex. lois raciales, esclavage), l'homme dit « c'est injuste ». Ce jugement suppose une norme supérieure à la loi écrite : la justice naturelle.

\begin{definition}[Droit naturel]
Ensemble de principes juridiques considérés comme universels, immuables et imprescriptibles, fondés sur la nature raisonnable de l'être humain (sa dignité, sa liberté, sa sociabilité). Ils s'imposent à tout législateur et servent de critère de validité aux lois positives.
\end{definition}

\textbf{Les grandes figures :}
\begin{itemize}
    \item \textbf{Aristote} (\emph{Ethique à Nicomaque}, V) : distingue le « juste naturel » (partout la même force) du « juste légal » (variable selon les cités). La loi naturelle est droite raison conforme à la nature.
    \item \textbf{Cicéron} (\emph{De Legibus}) : « La vraie loi est la droite raison, conforme à la nature, immuable, éternelle. » Elle ne peut être abrogée par le sénat ni par le peuple.
    \item \textbf{Les Lumières et 1948 :} Locke (droits naturels à la vie, liberté, propriété), Kant (droit inné à la liberté). Aboutissement : \textbf{Déclaration universelle des droits de l'homme (DUDH), 10 décembre 1948, 30 articles}. Son préambule proclame la « reconnaissance de la dignité inhérente à tous les membres de la famille humaine ».
\end{itemize}

\begin{exemplebox}[Exemple concret : L'esclavage et le droit naturel]
Au XVIIIe siècle, le \emph{Code Noir} (droit positif français) régissait l'esclavage aux Antilles. Il définissait l'esclave comme un « bien meuble ». Pour les jusnaturalistes (ex. l'abbé Raynal, Condorcet), cette loi positive était \textbf{radicalement injuste} car elle violait le droit naturel à la liberté et à la dignité. L'abolition de 1848 a été la victoire du droit naturel sur le droit positif inique.
\end{exemplebox}

\begin{savaistu}[Le préambule de la Constitution camerounaise de 1996]
Le Préambule de la Loi n°96-06 du 18 janvier 1996 portant révision de la Constitution du 2 juin 1972 affirme solennellement l'attachement du peuple camerounais « aux droits fondamentaux inscrits dans la Déclaration universelle des droits de l'homme de 1948 ». C'est une consécration constitutionnelle explicite de la supériorité hiérarchique du droit naturel (via la DUDH) sur la législation interne.
\end{savaistu}

\begin{exemplebox}[Exemple chiffré : La DUDH en chiffres]
\begin{itemize}
    \item \textbf{30 articles} adoptés le \textbf{10 décembre 1948} par l'Assemblée générale de l'ONU (Résolution 217 A III).
    \item \textbf{48 États} votants pour, 0 contre, 8 abstentions.
    \item Articles 1 et 2 : Fondement (dignité, liberté, égalité, non-discrimination).
    \item Articles 3 à 21 : Droits civils et politiques (vie, interdiction torture, procès équitable, liberté opinion).
    \item Articles 22 à 27 : Droits économiques, sociaux et culturels (travail, santé, éducation).
\end{itemize}
\end{exemplebox}

\courssub{Le droit positif : historicité et effectivité}

\emph{Intuition :} En pratique, le juge applique le Code pénal camerounais, le Code civil, la loi de finances. C'est le droit « en vigueur », identifiable, coercible.

\begin{definition}[Droit positif]
Ensemble des règles de droit effectivement en vigueur dans un État déterminé, à un moment donné de son histoire. Il comprend la Constitution, les lois votées par le Parlement, les décrets, les ordonnances, la jurisprudence et, au Cameroun, le droit coutumier. Il est \textbf{posé} (lat. \emph{ponere}) par une autorité compétente.
\end{definition}

\textbf{Caractères :}
\begin{itemize}
    \item \textbf{Relativité spatiale :} Le Code pénal camerounais ne s'applique qu'au Cameroun.
    \item \textbf{Relativité temporelle :} La loi nouvelle abroge l'ancienne (ex. Loi n°2016/007 du 12 juil. 2016 portant Code pénal abroge l'ordonnance de 1965).
    \item \textbf{Coercibilité :} Sanction étatique organisée (police, prisons, amendes).
\end{itemize}

\begin{attention}[Piège : « Naturel » $\neq$ « Instinctif »]
Ne confondez pas \textbf{droit naturel} (découvert par la \emph{raison}, norme éthique universelle : « ne tue pas ») et \textbf{loi naturelle au sens physique/biologique} (loi de la gravité, instinct de conservation, loi du plus fort biologique). Le droit naturel est une exigence \textbf{rationnelle et morale}, non un déterminisme biologique. L'homme peut violer le droit naturel (commettre un meurtre) ; il ne peut violer la loi de la gravité.
\end{attention}

\courssub{La thèse conventionnaliste : le contrat social}

Comment passer de l'état de nature à la société civile ? Par un pacte, une convention.

\begin{propriete}[Le conventionnalisme (Contractualisme)]
Le droit et l'État ne sont pas naturels mais artificiels : ils naissent d'un \textbf{contrat social} par lequel des individus libres et égaux s'engagent mutuellement à respecter des règles communes pour sortir de l'insécurité de l'état de nature.
\end{propriete}

\begin{center}
\begin{tabularx}{\linewidth}{|l|X|X|}
\hline
\textbf{Auteur} & \textbf{État de nature} & \textbf{Fondement du droit / But du contrat} \\ \hline
\textbf{Hobbes} (\emph{Léviathan}, 1651) & Guerre de tous contre tous (\emph{bellum omnium contra omnes}), vie « misérable, brutale et courte ». & \textbf{Sécurité.} Aliénation totale de la liberté au Souverain (Léviathan). Le droit = ordre du Souverain. \\ \hline
\textbf{Rousseau} (\emph{Du contrat social}, 1762) & Homme bon mais corrompu par la société naissante ; indépendance naturelle. & \textbf{Liberté civile \& Légitimité.} Aliénation totale à la \textbf{Volonté Générale}. « Obéir à la loi qu'on s'est prescrite, c'est la liberté. » Le droit = expression de la Volonté Générale. \\ \hline
\end{tabularx}
\end{center}

\begin{exoresolu}[Analyse d'un extrait : Rousseau, \emph{Du Contrat Social}, Livre I, Chap. 6]
\textbf{Énoncé :} « Trouver une forme d'association qui défende et protège de toute la force commune la personne et les biens de chaque associé, et par laquelle chacun, s'unissant à tous, n'obéisse pourtant qu'à lui-même et reste aussi libre qu'auparavant. »
\textbf{Solution :}
\begin{enumerate}
    \item \textbf{Problème :} Comment concilier obligation légale et liberté individuelle ?
    \item \textbf{Thèse :} Le contrat social fonde la légitimité du droit. La loi n'est pas une contrainte extérieure (hétéronomie) mais l'expression de ma propre volonté en tant que citoyen (autonomie).
    \item \textbf{Argument :} En votant la loi (ou en acceptant le pacte), je m'oblige moi-même. Le droit positif tire sa force obligatoire du consentement des gouvernés.
    \item \textbf{Portée :} Fondement démocratique de la légalité. Au Cameroun, le Préambule de 1996 pose : « La souveraineté nationale appartient au peuple qui l'exerce soit par référendum, soit par l'intermédiaire de ses représentants élus. »
\end{enumerate}
\end{exoresolu}

\courssub{La thèse de la force : critique du « droit du plus fort »}

\emph{Intuition :} « La raison du plus fort est toujours la meilleure » (La Fontaine, \emph{Le Loup et l'Agneau}). Calliclès (dans le \emph{Gorgias} de Platon) affirme que la justice naturelle, c'est que le fort domine le faible ; la loi est une invention des faibles pour se protéger.

\begin{attention}[Erreur fatale : Confondre Force et Droit]
\begin{itemize}
    \item \textbf{La force (Factum) :} Relation de causalité physique (contrainte, violence). Elle produit de l'obéissance par la peur.
    \item \textbf{Le droit (Jus) :} Relation de normativité (obligation morale/juridique). Il produit de l'obéissance par la reconnaissance de la légitimité.
    \item \textbf{Critique (Kant, Kelsen) :} Un bandit qui me vole sous la menace d'une arme exerce une force, il ne me fait pas un droit. « \emph{Ex injuria jus non oritur} » : le droit ne naît pas de l'injustice/violence. Un régime totalitaire a la force, mais lui manque la \textbf{légitimité} (le droit).
\end{itemize}
\end{attention}

\courssub{Cas camerounais : pluralisme juridique et coexistence des ordres}

Le Cameroun offre un laboratoire unique de \textbf{pluralisme juridique} (dualisme légal) hérité de l'histoire coloniale (française et britannique) et de la réalité sociologique.

\begin{definition}[Droit coutumier au Cameroun]
Ensemble de règles non écrites, issues de la tradition des groupes ethniques, régissant le statut des personnes (mariage, filiation, succession) et les biens (terres communautaires), reconnues et appliquées par l'ordre judiciaire (Loi n°72/5 du 26 août 1972, Art. 1er al. 2 : « La coutume est source de droit »).
\end{definition}

\begin{popfigure}
\begin{tikzpicture}[scale=0.85, every node/.style={font=\small}]
    % Contour très simplifié du Cameroun (polygone approximatif)
    \draw[popdark, thick, fill=popcream] 
        (0,0) -- (1.5,0.2) -- (3.5,0.5) -- (5.5,0.2) -- (7,0) -- % Bas (Sud)
        (7.2,1.5) -- (7.5,3) -- (7.2,4.5) -- (6.5,5.5) -- (5.5,6) -- % Est
        (4.5,6.5) -- (3.5,6.2) -- (2.5,6.5) -- (1.5,6) -- (0.5,5.5) -- % Nord
        (0,4) -- (0.2,2.5) -- (0,1) -- cycle;
    
    % Aires coutumières dominantes (zones colorées)
    % Nord (Soudano-sahélien) : Droit coutumier islamisé / chefferies lamidats
    \fill[popblue!20] (1.5,4.5) -- (3.5,4.8) -- (5.5,5) -- (5.5,6) -- (3.5,6.2) -- (1.5,6) -- cycle;
    \node[popblue, align=center, font=\small\bfseries] at (3.5,5.3) {Zone Nord\\ (Lamidats, Lawane,\\ Droit coutumier islamisé)};
    
    % Ouest (Hauts-Plateaux) : Droit coutumier bantou / chefferies supérieures
    \fill[popgreen!20] (0.5,2.5) -- (3.5,2.8) -- (3.5,4.8) -- (1.5,4.5) -- cycle;
    \node[popgreen, align=center, font=\small\bfseries] at (2,3.5) {Zone Ouest\\ (Chefferies 1er/2nd degré,\\ Droit coutumier bantou)};
    
    % Sud/Centre/Est : Droit écrit dominant (Code civil), coutume résiduelle
    \fill[poporange!15] (0,0) -- (5.5,0.2) -- (5.5,2.5) -- (3.5,2.8) -- (0.5,2.5) -- cycle;
    \node[poporange, align=center, font=\small\bfseries] at (2.5,1.2) {Zone Sud/Centre/Est\\ (Droit écrit dominant,\\ Coutume : successions/mariage)};
    
    % Villes repères
    \node[circle, fill=popdark, inner sep=1.5pt] (Ya) at (2.2, 1.8) {};
    \node[above right] at (Ya) {Yaoundé};
    \node[circle, fill=popdark, inner sep=1.5pt] (Do) at (4.5, 0.8) {};
    \node[above right] at (Do) {Douala};
    \node[circle, fill=popdark, inner sep=1.5pt] (Ga) at (3.5, 5.2) {};
    \node[above right] at (Ga) {Garoua};
    \node[circle, fill=popdark, inner sep=1.5pt] (Ba) at (1.5, 3.8) {};
    \node[above right] at (Ba) {Bafoussam};
    \node[circle, fill=popdark, inner sep=1.5pt] (Ma) at (6.2, 3.0) {};
    \node[above right] at (Ma) {Maroua};
    
    % Flèche Nord
    \draw[popdark, thick, ->] (6.5, 1.5) -- (6.5, 2.2) node[right] {N};
\end{tikzpicture}
\legende{Carte schématique du Cameroun : aires de dominance du droit coutumier (Nord : chefferies lamidats ; Ouest : chefferies bantoues) et zone de prédominance du droit écrit (Sud/Centre). La loi positive (Code civil, Code pénal) s'applique sur l'ensemble du territoire, la coutume intervient en matière personnelle/foncière selon l'article 1er de la loi de 1972.}
\end{popfigure}

\begin{exemplebox}[Concret : Mariage et Successions]
\begin{itemize}
    \item \textbf{Mariage :} Un Camerounais peut choisir le mariage civil (Code civil, monogame, devant l'officier d'état civil) \textbf{OU} le mariage coutumier (polygame possible, dot, devant le chef de famille/chef traditionnel). Les deux produisent des effets juridiques (Art. 144 Code civil + Coutume).
    \item \textbf{Successions :} Si le défunt n'a pas fait de testament, la dévolution successorale suit souvent la coutume (ex. primogéniture masculine, exclusion des filles/femme chez certains groupes) sauf opposition expresse ou saisine du tribunal pour application de la loi moderne (Code civil : égalité hommes/femmes). C'est un contentieux majeur.
    \item \textbf{Terres :} Dans les zones rurales (Nord, Ouest), l'accès à la terre relève souvent du chef de terre / lamido (coutume), alors que la loi foncière (Ordonnance 74-1, 74-2) impose l'immatriculation pour la propriété privée. Conflit fréquent : coutume vs titre foncier.
\end{itemize}
\end{exemplebox}

\courssub{Tableau comparatif : Droit naturel vs Droit positif}

\begin{popfigure}
\begin{center}
\begin{tabularx}{\linewidth}{|>{\bfseries\centering\arraybackslash}p{2.8cm}|X|X|}
\hline
\rowcolor{popblue}
\textcolor{white}{\textbf{Critère}} & \textcolor{white}{\textbf{Droit Naturel (\emph{Jus Naturale})}} & \textcolor{white}{\textbf{Droit Positif (\emph{Jus Positum})}} \\ \hline
\textbf{Fondement} & La nature raisonnable de l'homme / Dignité intrinsèque & La volonté du souverain / Parlement / Coutume / Contrat social \\ \hline
\textbf{Caractères} & Universel (valable partout) / Immuable (éternel) / Imprescriptible & Relatif (spatial \& temporel) / Révocable / Coercible par l'État \\ \hline
\textbf{Source} & La Raison / La Loi divine (pour certains) / La Nature & La Loi écrite (Constitution, Codes) / Jurisprudence / Coutume reconnue \\ \hline
\textbf{Contenu} & Droits de l'homme (Vie, Liberté, Égalité, Propriété) / Justice idéale & Règles techniques (Code route, Fiscalité) / Organisation de l'État / Sanctions pénales \\ \hline
\textbf{Exemples} & DUDH 1948 / Déclaration 1789 / Préambule Const. Cameroun 1996 & Code civil camerounais / Code pénal 2016 / Loi électorale / Coutume (mariage, terre) \\ \hline
\end{tabularx}
\end{center}
\legende{Tableau comparatif : Le droit naturel fournit la \textbf{légitimité morale} (critère de justice), le droit positif assure l'\textbf{effectivité juridique} (sanction étatique). Le droit positif \textbf{légitime} doit s'inspirer du droit naturel.}
\end{popfigure}

\courssub{Conséquence majeure : La désobéissance civile}

Si le droit positif peut s'écarter du droit naturel (lois injustes), le citoyen se trouve face à un conflit de conscience : \emph{obéir à la loi injuste ou obéir à sa conscience (droit naturel) ?}

\begin{propriete}[Légitimité de la désobéissance civile]
Une loi positive manifestement contraire aux droits fondamentaux (droit naturel) perd son caractère obligatoire en conscience. La \textbf{désobéissance civile} (désobéissance publique, non-violente, assumée, acceptant la sanction légale) devient un devoir moral pour alerter la communauté politique sur l'injustice.
\end{propriete}

\begin{exemplebox}[Exemples historiques et camerounais]
\begin{itemize}
    \item \textbf{Lois de l'apartheid (Afrique du Sud) :} Lois positives (Population Registration Act, Group Areas Act) légalisant la ségrégation raciale. Nelson Mandela, Martin Luther King (USA, ségrégation), Gandhi (Inde, lois coloniales britanniques) ont pratiqué la désobéissance civile au nom du droit naturel (égalité, dignité).
    \item \textbf{Au Cameroun :} La lutte pour l'indépendance (UPC, années 50-60) a été une désobéissance massive à la loi coloniale (code de l'indigénat, travail forcé) au nom du droit des peuples à disposer d'eux-mêmes (droit naturel / DUDH Art. 1 \& 21).
    \item \textbf{Contexte actuel :} Manifestations pacifiques (syndicats, société civile) contre des lois perçues comme liberticides (ex. loi antiterroriste 2014, code pénal 2016 sur la diffamation/propagation de fausses nouvelles) invoquent la supériorité des libertés fondamentales (Constitution, DUDH) sur la loi ordinaire.
\end{itemize}
\end{exemplebox}

\begin{methode}[Comment analyser un conflit Droit Naturel / Droit Positif (Dissertation/Commentaire)]
\begin{enumerate}
    \item \textbf{Identifier la loi positive en cause} (texte, article, contexte d'application).
    \item \textbf{Identifier le principe de droit naturel invoqué} (dignité, liberté, égalité, vie, propriété - se référer à la Constitution, DUDH, Charte africaine).
    \item \textbf{Montrer la contradiction} : En quoi la loi positive viole-t-elle le principe naturel ? (Ex. : Loi autorisant la torture vs Droit à l'intégrité physique Art. 5 DUDH / Art. 1er Const. Cameroun).
    \item \textbf{Qualifier le régime} : Est-ce une simple imperfection (loi perfectible) ou une loi \textbf{radicalement injuste} (atteinte au noyau dur des droits de l'homme) ?
    \item \textbf{Conclure sur la réponse juridique et morale} : Voie de droit (recours constitutionnel, Cour suprême, Cour africaine DH) vs Désobéissance civile (si voies de droit bloquées/ineffectives).
\end{enumerate}
\end{methode}

\begin{aretenir}[Synthèse : Hiérarchie des normes et fondements]
\begin{itemize}
    \item Le \textbf{Droit Naturel} = Fondement \textbf{éthique} (Justice idéale, Universalité). Critère de validité morale.
    \item Le \textbf{Droit Positif} = Fondement \textbf{juridique/formel} (Légalité, Effectivité, Coercibilité). Outil de régulation sociale.
    \item \textbf{Relation idéale :} Tout droit positif \textbf{légitime} s'inspire du droit naturel (Préambule Const. 1996).
    \item \textbf{Relation réelle :} Le droit positif \textbf{peut} s'écarter du droit naturel (lois iniques).
    \item \textbf{Recours :} Contrôle de constitutionnalité (Conseil constitutionnel), Contrôle de conventionalité (Traités internationaux > Loi, Art. 45 Const.), Désobéissance civile (dernier ressort moral).
\end{itemize}
\end{aretenir}

\textbf{Transition :} Maintenant que les fondements du droit sont posés (naturels, positifs, conventionnels), nous devons examiner comment la justice, comme valeur suprême, se réalise concrètement à travers l'égalité et les différentes formes de justice (commutative, distributive, réparatrice), pour tendre vers l'idée absolue de justice. \par
\courssub{Justice et égalité : vers les formes de justice (Annonce de la suite)}

\coursec{Justice et égalité}

Après avoir examiné les fondements qui donnent au droit sa force obligatoire — la volonté de l'État, la raison naturelle ou le contrat social —, nous devons maintenant interroger le critère même de la justice. Une loi peut être légitime dans son origine, mais injuste dans ses effets si elle ne respecte pas l'égalité réelle entre les citoyens. La question centrale de cette section est donc : \emph{être juste, est-ce traiter tout le monde de la même manière ?} L'intuition spontanée répond souvent par l'affirmative : la justice serait l'égalité pure et simple. Pourtant, l'expérience montre qu'appliquer la même règle à des situations différentes peut produire des résultats profondément injustes. C'est toute la distinction entre \emph{égalité arithmétique} et \emph{égalité proportionnelle} qu'il faut maîtriser.

\courssub{Le problème : être juste, est-ce traiter tout le monde de la même manière ?}

Imaginez un professeur qui donne la même note à tous ses élèves, qu'ils aient travaillé ou non, au nom de l'égalité. Ou un juge qui inflige la même peine de prison à un voleur de pomme et à un assassin, au nom de l'égalité devant la loi. Dans les deux cas, l'égalité formelle (même traitement) conduit à l'injustice matérielle (mépris du mérite ou de la gravité). Le problème de la justice est donc de déterminer \emph{quel critère} de répartition ou de traitement est pertinent selon les situations : le nombre de personnes (arithmétique), le mérite, le besoin, la contribution, la gravité de la faute ?

\begin{definition}[Égalité arithmétique et égalité proportionnelle]
\textbf{Égalité arithmétique} : mode de répartition qui attribue la même part à chaque individu, sans tenir compte de leurs différences (besoins, mérites, capacités). C'est l'égalité \emph{quantitative}, « un pour un » (ex. : un citoyen = une voix).

\textbf{Égalité proportionnelle} (ou géométrique, chez Aristote) : mode de répartition qui attribue à chacun une part proportionnée à un critère pertinent (mérite, besoin, travail, contribution, faute). C'est l'égalité \emph{qualitative}, « à chacun selon... » (ex. : impôt proportionnel au revenu).
\end{definition}

\courssub{L'égalité arithmétique : même part pour tous}

L'égalité arithmétique s'applique là où les individus sont considérés comme \emph{strictement égaux} en dignité ou en droit fondamental, sans critère de différenciation pertinent. C'est le principe « une personne, une voix » dans le suffrage universel : le vote d'un professeur d'université compte autant que celui d'un cultivateur, d'un étudiant ou d'un chef d'entreprise. De même, devant la loi pénale, le principe \emph{nulla poena sine lege} implique que la même infraction entraîne la même peine prévue par le texte, quel que soit l'auteur (égalité formelle des justiciables).

\begin{exemplebox}[Le droit de vote au Cameroun]
Aux élections présidentielles, législatives ou municipales, chaque électeur inscrit dispose d'une seule voix. Que l'on soit milliardaire à Douala ou paysan dans le Mayo-Danay, le bulletin de vote a le même poids. C'est l'application pure de l'égalité arithmétique : la souveraineté populaire s'exprime par le nombre, non par la fortune ou le diplôme.
\end{exemplebox}

\courssub{L'égalité proportionnelle (Aristote) : à chacun selon son mérite, son travail ou ses besoins}

Dans l'\emph{Éthique à Nicomaque} (Livre V), Aristote distingue la justice \emph{distributive} (répartition des honneurs, richesses, charges) de la justice \emph{commutative} (échanges, contrats). Pour la justice distributive, il pose le principe : « Le juste est le proportionnel. » Il ne s'agit pas de donner la même chose à tous, mais de donner \emph{à chacun selon sa valeur} (mérite, vertu), \emph{selon son travail} (contribution) ou \emph{selon ses besoins} (nécessité vitale). La proportionnalité s'écrit : $\frac{\text{Part}_A}{\text{Critère}_A} = \frac{\text{Part}_B}{\text{Critère}_B}$.

\begin{exemplebox}[Trois critères de proportionnalité]
\begin{itemize}
    \item \textbf{Mérite} : bourses d'excellence aux meilleurs notes.
    \item \textbf{Travail / Contribution} : salaire proportionnel aux heures ou à la productivité ; cotisation retraite proportionnelle aux versements.
    \item \textbf{Besoin} : aide sociale, allocation familiale, gratuité des soins pour les indigents.
\end{itemize}
Le choix du critère est lui-même un acte de justice politique : une société qui ne retient que le mérite néglige les plus fragiles ; une qui ne retient que le besoin peut décourager l'effort.
\end{exemplebox}

\courssub{Applications concrètes au Cameroun}

\begin{exemplebox}[L'impôt progressif sur le revenu (IRPP) au Cameroun]
L'impôt sur le revenu des personnes physiques (IRPP) applique une \emph{progressivité} par tranches : plus le revenu net imposable est élevé, plus le taux marginal augmente. C'est l'égalité proportionnelle au service de la solidarité nationale : ceux qui ont une plus grande capacité contributive paient non seulement plus en valeur absolue, mais proportionnellement plus.

\begin{center}
\begin{tabularx}{\linewidth}{|l|X|X|}\hline
\textbf{Tranche de revenu net imposable (FCFA)} & \textbf{Taux marginal} & \textbf{Observation} \\ \hline
Jusqu'à 2 000 000 & 11\,\% & Taux d'entrée \\ \hline
2 000 001 -- 3 000 000 & 16,5\,\% & \\ \hline
3 000 001 -- 5 000 000 & 22\,\% & \\ \hline
5 000 001 -- 10 000 000 & 27,5\,\% & \\ \hline
10 000 001 -- 20 000 000 & 33\,\% & \\ \hline
Au-delà de 20 000 000 & 38,5\,\% & Taux maximal \\ \hline
\end{tabularx}
\end{center}

\emph{Exemple chiffré (salariés, après abattement de 20\,\% pour frais professionnels)} : Un célibataire sans enfant gagnant 1 500 000 FCFA/mois (18 000 000/an, soit 14 400 000 net imposable) paiera environ 3 650 000 FCFA d'IRPP (taux effectif sur le brut ~20\,\%). Un cadre gagnant 5 000 000 FCFA/mois (60 000 000/an, soit 48 000 000 net imposable) paiera environ 16 300 000 FCFA (taux effectif sur le brut ~27\,\%). Le second gagne 3,3 fois plus, mais paie environ 4,5 fois plus d'impôt. C'est la \textbf{progressivité} : l'égalité proportionnelle à la capacité contributive.
\end{exemplebox}

\begin{exemplebox}[Bourses d'études et discrimination positive (quotas régionaux)]
L'attribution des bourses d'excellence (mérite) repose sur l'égalité proportionnelle aux résultats scolaires. Mais l'accès aux grandes écoles (ENSET, ENS, Mines, etc.) intègre souvent une \emph{discrimination positive} via des quotas régionaux : chaque région a un nombre minimal de places réservées. Le critère n'est plus seulement le mérite individuel, mais l'\emph{équilibre territorial} et la correction des inégalités structurelles (écoles sous-équipées, éloignement). C'est une application de l'égalité proportionnelle aux \emph{besoins collectifs} de représentation et de développement équilibré.
\end{exemplebox}

\courssub{L'équité : corriger la rigueur de la loi (Aristote)}

L'\textbf{équité} (grec \emph{epikeia}, latin \emph{aequitas}) est la vertu qui permet de corriger la rigueur de la loi générale quand son application aveugle à un cas particulier aboutirait à l'injustice. Aristote déjà : « L'équitable est juste, mais non pas le juste légal ; c'est une correction du juste légal. » Le juge équitable (ou le bon administrateur) regarde la \emph{situation concrète} : l'intention, les circonstances, la vulnérabilité, l'histoire personnelle. C'est le passage de la \emph{justice légale} (égalité formelle) à la \emph{justice équitable} (égalité réelle).

\begin{definition}[Équité]
L'équité est la disposition à adapter la règle générale aux particularités du cas d'espèce pour atteindre la justice matérielle. Elle ne nie pas la loi, mais la tempère là où sa lettre contredit son esprit. Ex. : ne pas sanctionner un retard justifié par un cas de force majeure, bien que le règlement dise « tout retard est sanctionné ».
\end{definition}

\begin{popfigure}
\begin{tikzpicture}[scale=0.9, every node/.style={font=\small}]
    % Base
    \draw[popdark, thick] (0,0) -- (10,0);
    % Piedestal
    \draw[popdark, thick] (4.5,0) -- (4.5,-0.5) -- (5.5,-0.5) -- (5.5,0);
    % Colonne centrale
    \draw[popdark, thick] (5,-0.5) -- (5,3);
    % Bras de la balance (horizontal = équilibre)
    \draw[popdark, thick] (1,3) -- (9,3);
    % Pivot
    \draw[popdark, fill=popgold] (5,3) circle (0.15);
    % Plateaux (horizontaux)
    \draw[popblue, thick] (1.5,3) -- (1.5,3.8) -- (3.5,3.8) -- (3.5,3);
    \draw[poporange, thick] (6.5,3) -- (6.5,3.8) -- (8.5,3.8) -- (8.5,3);
    % Ficelles
    \draw[popmuted] (1.5,3.8) -- (1.5,4.3) -- (2.5,4.3);
    \draw[popmuted] (3.5,3.8) -- (3.5,4.3) -- (2.5,4.3);
    \draw[popmuted] (6.5,3.8) -- (6.5,4.3) -- (7.5,4.3);
    \draw[popmuted] (8.5,3.8) -- (8.5,4.3) -- (7.5,4.3);
    % Poids sur plateau gauche (besoin élevé -> gros poids)
    \draw[popblue, fill=popblue!20, rounded corners] (2,4.4) rectangle (3,5.6);
    \node[popblue, align=center, font=\small\bfseries] at (2.5,5.0) {Part importante\\ (Besoin élevé)};
    % Poids sur plateau droit (besoin faible -> petit poids)
    \draw[poporange, fill=poporange!20, rounded corners] (7,4.4) rectangle (8,5.0);
    \node[poporange, align=center, font=\small\bfseries] at (7.5,4.7) {Part faible\\ (Besoin faible)};
    % Flèche d'équilibre
    \draw[popgreen, thick, <->] (3.5,3.4) -- (6.5,3.4) node[midway, above] {\textbf{Équilibre (Justice)}};
    % Légendes
    \node[popdark, align=center] at (2.5,2.2) {\textbf{Égalité proportionnelle}\\\emph{Parts inégales, balance juste}};
    \node[popdark, align=center] at (7.5,2.2) {\textbf{Égalité arithmétique}\\\emph{(Parts égales, déséquilibre si besoins inégaux)}};
    % Titre
    \node[popdark, font=\bfseries\large] at (5,6.2) {La balance de la justice : l'égalité proportionnelle rétablit l'équilibre};
\end{tikzpicture}
\legende{Schéma : l'égalité proportionnelle place des poids différents (parts adaptées aux besoins) sur les plateaux. La balance reste horizontale : c'est la justice effective. L'égalité arithmétique (mêmes poids pour besoins inégaux) ferait pencher la balance du côté de celui qui a moins besoin (injustice matérielle).}
\end{popfigure}

\courssub{Les limites de l'égalité stricte : quand l'égalité produit l'injustice}

\begin{attention}[Piège fréquent : confondre justice et égalité absolue]
\textbf{Erreur} : « La justice, c'est traiter tout le monde pareil. » \textbf{Faux.} La justice formelle (même loi pour tous) est une condition nécessaire, mais non suffisante. Appliquer aveuglément la même règle à des situations inégales \emph{crée} de l'injustice.
\begin{itemize}
    \item \textbf{Même punition pour fautes inégales} : un élève qui oublie son livre et un élève qui insulte le professeur reçoivent la même exclusion ? Injuste.
    \item \textbf{Même impôt pour revenus inégaux} : un impôt forfaitaire (même somme pour tous) écrase les pauvres et épargne les riches. C'est l'injustice de la \emph{capitation} (impôt par tête) que la Révolution française a abolie.
    \item \textbf{Mêmes droits sans moyens d'exercice} : le droit à la santé « pour tous » reste vide si les hôpitaux sont absents dans les zones rurales. L'égalité réelle exige des moyens différenciés (plus de ressources là où le besoin est plus grand).
\end{itemize}
\end{attention}

\courssub{Illustration comparative : répartition arithmétique vs proportionnelle}

\begin{popfigure}
\begin{tikzpicture}
\begin{axis}[
    ybar,
    bar width=18pt,
    width=\linewidth,
    height=9cm,
    enlarge x limits=0.25,
    ylabel={Montant reçu (FCFA)},
    xtick={1,2,3},
    xticklabels={A (faible), B (moyen), C (élevé)},
    nodes near coords,
    nodes near coords align={vertical},
    every node near coord/.append style={font=\small, rotate=90, anchor=west, yshift=2pt},
    ymin=0,
    ymax=120000,
    legend style={at={(0.5,-0.22)}, anchor=north, legend columns=2, font=\small},
    ymajorgrids=true,
    grid style=dashed,
]
\addplot[popblue!70, fill=popblue!30] coordinates {
    (1, 50000)
    (2, 50000)
    (3, 50000)
};
\addplot[poporange!70, fill=poporange!30] coordinates {
    (1, 20000)
    (2, 50000)
    (3, 100000)
};
\legend{Égalité arithmétique (150 000 / 3), Égalité proportionnelle (selon le besoin)}
\end{axis}
\end{tikzpicture}
\legende{Répartition de 150 000 FCFA entre trois personnes. \textbf{Bleu} : égalité arithmétique (50 000 chacun) — la personne C (besoin élevé) est sous-dotée, la personne A (besoin faible) sur-dotée. \textbf{Orange} : égalité proportionnelle aux besoins — la justice matérielle est respectée.}
\end{popfigure}

\courssub{Méthode : analyser une situation de justice}

\begin{methode}[Analyser une situation de répartition ou de sanction]
Pour juger si une décision est juste, suivez ces étapes :
\begin{enumerate}
    \item \textbf{Identifier la règle appliquée} : quelle est la norme, la loi, le règlement, l'usage ?
    \item \textbf{Identifier le critère de répartition implicite ou explicite} : est-ce le nombre (arithmétique) ? Le mérite ? Le besoin ? Le travail ? La faute ?
    \item \textbf{Comparer les situations réelles des personnes concernées} : sont-elles égales ou inégales au regard du critère pertinent ?
    \item \textbf{Évaluer l'adéquation} : le critère choisi est-il pertinent par rapport à l'objet de la décision (éducation, fiscalité, sanction, santé) ?
    \item \textbf{Conclure} : si le critère est inadapté (ex. arithmétique pour des besoins inégaux), la décision est formellement égale mais matériellement injuste. Proposer le critère proportionnel pertinent.
\end{enumerate}
\end{methode}

\begin{exoresolu}[Analyse d'une sanction scolaire]
\textbf{Énoncé} : Dans un lycée, deux élèves arrivent en retard de 15 minutes. L'un habitait à 500 m de l'école, l'autre à 15 km et a eu une panne de bus. Le règlement stipule : « Tout retard > 10 min = 1 heure de retenue. » Les deux reçoivent la même retenue. Est-ce juste ?

\textbf{Solution détaillée} :
\begin{enumerate}
    \item \textbf{Règle} : Règlement intérieur, sanction unique pour tout retard > 10 min (critère arithmétique : 1 retard = 1 retenue).
    \item \textbf{Critère} : La gravité de la faute (volonté, négligence) — mais le règlement ne la regarde pas.
    \item \textbf{Situations réelles} : Élève A (proche, pas d'excuse valable) = négligence. Élève B (loin, panne transport) = cas de force majeure, non responsable.
    \item \textbf{Adéquation} : Le critère « retard brut » est inadapté : il ne distingue pas la faute volontaire de l'empêchement extérieur. L'égalité formelle produit une injustice : l'élève B est puni pour un événement qu'il ne contrôle pas.
    \item \textbf{Conclusion} : La décision est \emph{légalement} conforme, mais \emph{équitablement} injuste. L'équité exigerait d'exonérer l'élève B (ou de sanctionner différemment). Le règlement devrait prévoir une clause de force majeure.
\end{enumerate}
\end{exoresolu}

\begin{exercice}[À vous de jouer]
Une entreprise verse une prime de fin d'année de 1 000 000 FCFA à partager entre trois employés :
\begin{itemize}
    \item Employé X : ancienneté 15 ans, salaire de base 300 000 FCFA, performance moyenne.
    \item Employé Y : ancienneté 5 ans, salaire de base 500 000 FCFA, performance excellente.
    \item Employé Z : ancienneté 1 an, salaire de base 200 000 FCFA, performance bonne.
\end{itemize}
Le directeur propose trois modes de répartition :
\begin{enumerate}
    \item Égalité arithmétique (un tiers chacun).
    \item Proportionnelle au salaire de base.
    \item Proportionnelle à un score composite : 40\,\% ancienneté, 40\,\% performance, 20\,\% salaire.
\end{enumerate}
\textbf{Question} : Calculez les trois répartitions. Laquelle vous semble la plus juste ? Justifiez en identifiant le critère dominant de chaque mode et son adéquation à la finalité d'une « prime de fin d'année ».
\end{exercice}

\begin{corrige}[Exercice n°1]
\begin{enumerate}
    \item \textbf{Arithmétique} : 333 333 FCFA chacun. Critère : l'appartenance à l'entreprise (statut d'employé). Ignore tout effort ou ancienneté.
    \item \textbf{Proportionnelle au salaire} : Total salaires = 1 000 000. X : 300 000 ; Y : 500 000 ; Z : 200 000. Critère : la valeur marchande du poste / la responsabilité hiérarchique. Avantage les cadres, pénalise les anciens aux salaires modestes.
    \item \textbf{Score composite} : Calculons les points (base 100).
    \begin{itemize}
        \item Ancienneté (max 15 ans) : X=100, Y=33, Z=7.
        \item Performance (échelle : excellente=100, bonne=70, moyenne=50) : X=50, Y=100, Z=70.
        \item Salaire (max 500 000) : X=60, Y=100, Z=40.
    \end{itemize}
    Score X = 0,4$\times$100 + 0,4$\times$50 + 0,2$\times$60 = 40 + 20 + 12 = 72.
    Score Y = 0,4$\times$33 + 0,4$\times$100 + 0,2$\times$100 = 13,2 + 40 + 20 = 73,2.
    Score Z = 0,4$\times$7 + 0,4$\times$70 + 0,2$\times$40 = 2,8 + 28 + 8 = 38,8.
    Total = 184. Parts : X = 72/184 $\times$ 1 000 000 $\approx$ 391 304 ; Y $\approx$ 397 826 ; Z $\approx$ 210 870.
    \item \textbf{Analyse} : La prime de fin d'année vise généralement à \emph{récompenser la contribution globale} (fidélité + effort + responsabilité). Le mode 3 est le plus juste car il intègre les trois dimensions pertinentes. Le mode 1 est trop plat (démotivant). Le mode 2 réduit la justice à la hiérarchie salariale, oubliant l'ancienneté et le mérite actuel.
\end{enumerate}
\end{corrige}

\begin{aretenir}[Synthèse : Justice et égalité]
\begin{itemize}
    \item \textbf{Justice $\neq$ Égalité arithmétique}. Traiter égaux des inégaux est une injustice.
    \item \textbf{Égalité arithmétique} : même part pour tous. Pertinente pour les droits fondamentaux (vote, dignité, peine légale pour même infraction).
    \item \textbf{Égalité proportionnelle (Aristote)} : à chacun selon un critère pertinent (mérite, travail, besoin, faute). Pertinente pour les répartitions (richesses, charges, honneurs, sanctions).
    \item \textbf{Équité} : correction de la règle générale par l'appréciation du cas concret. Passage de la légalité à la justice réelle.
    \item \textbf{Méthode} : identifier le critère $\rightarrow$ comparer les situations réelles $\rightarrow$ juger la pertinence du critère $\rightarrow$ conclure.
\end{itemize}
\end{aretenir}

\coursec{Les formes de justice}

\courssub{Transition : de l'égalité aux formes concrètes de la justice}

Dans la section précédente, nous avons vu que la justice suppose l'égalité, mais que cette égalité ne saurait être arithmétique (donner la même chose à tous) : elle doit être \emph{proportionnelle} (donner à chacun selon un critère pertinent). C'est précisément ici que la philosophie distingue les \textbf{formes de la justice} : selon qu'il s'agit de régler un échange entre deux personnes, de répartir des biens communs dans la société ou de sanctionner une faute, le critère de proportionnalité change. Aristote, dans l'\emph{Éthique à Nicomaque} (Livre V), est le premier à opérer cette distinction fondamentale entre \textbf{justice commutative}, \textbf{justice distributive} et \textbf{justice rétributive} (ou corrective/pénale). Comprendre ces formes, c'est saisir comment le droit tente de réaliser concrètement l'idéal de justice.

\begin{definition}[Les trois formes classiques de la justice (Aristote)]
\emph{Justice commutative} : régit les échanges entre particuliers (contrats, ventes, réparations de dommages). Elle vise l'\textbf{égalité arithmétique} : chose pour chose, valeur pour valeur (ex. : rembourser exactement le préjudice subi).

\emph{Justice distributive} : régit la répartition des biens, honneurs, charges et devoirs au sein de la communauté politique. Elle vise l'\textbf{égalité géométrique (proportionnelle)} : à chacun selon son mérite, son besoin, sa contribution ou son rang (ex. : impôt progressif, attribution des bourses sur critères sociaux).

\emph{Justice rétributive (ou pénale)} : régit la sanction des fautes et des crimes. Elle vise la \textbf{proportionnalité entre la peine et la gravité de la faute} (lex talionis tempérée : « œil pour œil » devient peine proportionnée). Elle remplit trois fonctions : \emph{réparation} (pour la victime), \emph{dissuasion} (empêcher la récidive), \emph{rééducation} (réinsérer le coupable).
\end{definition}

\begin{propriete}[Exhaustivité des formes]
Toute décision de justice concrète, tout litige porté devant un tribunal relève d'au moins une de ces trois formes. Un même procès peut les combiner : le voleur est condamné à une peine de prison (rétributive), à rembourser l'objet volé (commutative) et l'État peut confisquer ses biens illicites pour les redistribuer (distributive).
\end{propriete}

\courssub{La justice commutative : l'égalité dans l'échange}

\textbf{Intuition et structure.}\par
Imaginez un marché à Douala : un acheteur paie 10 000 FCFA pour un poulet. Si le vendeur donne un poulet maigre valant 5 000 FCFA, l'échange est inégal, donc injuste. La justice commutative exige que la valeur donnée = valeur reçue. C'est la justice du \emph{do ut des} (« je donne pour que tu donnes »).

\textbf{Mécanisme : la restitution de l'égalité.}\par
En cas de violation (vol, dol, inexécution contractuelle, accident), le juge ramène les parties à l'égalité initiale par la \textbf{réparation intégrale du préjudice} (dommages-intérêts). Le critère est purement \emph{arithmétique} : le montant du dommage ne dépend ni du mérite de la victime, ni de la richesse du responsable, mais de la perte réelle subie.

\begin{exemplebox}[Réparation civile d'un accident de la route à Douala]
Monsieur \textsc{Ngono} traverse légalement au passage piéton. Un taxi, excès de vitesse, le renverse. Fracture de la jambe, 3 mois d'arrêt de travail (salaire : 200 000 FCFA/mois), frais médicaux 150 000 FCFA, préjudice moral estimé 100 000 FCFA.
\begin{itemize}
    \item \textbf{Justice commutative appliquée :} Le tribunal condamne l'assureur du taxi à verser : $(3 \times 200\,000) + 150\,000 + 100\,000 = 850\,000$ FCFA.
    \item \textbf{Point clé :} On ne demande pas si M. Ngono est « méritant » ou « pauvre » (critère distributif), ni si le chauffeur est « méchant » (critère rétributif). On calcule la \emph{perte} pour la rétablir. L'égalité arithmétique est restaurée.
\end{itemize}
\end{exemplebox}

\begin{attention}[Piège fréquent : confondre réparation et punition]
En justice commutative, le but n'est \textbf{pas} de punir le responsable (c'est le rôle du pénal), mais de \textbf{réparer la victime}. Si le chauffeur paie 850 000 FCFA, la justice commutative est satisfaite, même s'il n'a pas de peine de prison. À l'inverse, une peine de prison sans indemnisation laisse la victime non réparée : injustice commutative persistante.
\end{attention}

\courssub{La justice distributive : la répartition proportionnelle dans la cité}

\textbf{Intuition et structure.}\par
L'État perçoit des impôts et distribue des services (éducation, santé, bourses, marchés publics). Faut-il donner la même chose à tous ? Non : un étudiant brillant sans argent a \emph{besoin} d'une bourse ; un contribuable riche a la \emph{capacité} de payer plus d'impôts. La justice distributive cherche le \textbf{critère pertinent de proportionnalité}.

\textbf{Les grands critères historiques et actuels.}\par
\begin{itemize}
    \item \textbf{Le mérite (Aristote, méritocratie républicaine) :} À chacun selon ses capacités et ses efforts (concours, notes, décorations).
    \item \textbf{Le besoin (socialisme, droit social, Rawls) :} À chacun selon ses besoins vitaux (allocations familiales, CMU, bourses sur critères sociaux).
    \item \textbf{La contribution / l'investissement (libéralisme) :} À chacun selon ce qu'il a apporté au gâteau commun (dividendes, retraites par capitalisation).
    \item \textbf{Le rang / la naissance (Ancien Régime, castes) :} Critère aujourd'hui rejeté par les droits de l'homme (égalité devant la loi).
\end{itemize}

\begin{exemplebox}[Attribution des bourses d'excellence au Cameroun]
Le MINESUP publie un arrêté : bourses sur \textbf{critères de mérite} (moyenne $\ge$ 14/20) ET \textbf{critères sociaux} (revenus parents $<$ seuil).
\begin{itemize}
    \item Un étudiant moy. 15/20, parents hauts revenus $\to$ bourse mérite seulement.
    \item Un étudiant moy. 12/20, parents très modestes $\to$ bourse sociale seulement.
    \item Un étudiant moy. 15/20, parents modestes $\to$ bourse complète (mérite + besoin).
\end{itemize}
C'est une \textbf{justice distributive mixte} : le critère n'est pas unique, la loi combine mérite et besoin pour approcher l'équité.
\end{exemplebox}

\courssub{La justice rétributive (pénale) : punir proportionnellement}

\textbf{Intuition et structure.}\par
Quelqu'un commet un vol, une agression, un meurtre. La société ne peut pas laisser faire : l'ordre public est rompu. La justice rétributive inflige une \textbf{peine} (prison, amende, TIG, peine de mort là où elle existe). Le principe cardinal : \textbf{proportionnalité de la peine à la gravité de la faute} (art. 8 DHUD, art. 7 PIDCP, Constitution camerounaise).

\textbf{Les trois fonctions de la peine (doctrine pénale moderne).}\par
\begin{enumerate}
    \item \textbf{Réparation (expiation) :} La peine « paye la dette » envers la société et la victime (réparation morale, parfois financière via partie civile).
    \item \textbf{Dissuasion (prévention générale) :} La sévérité annoncée décourage les candidats au crime.
    \item \textbf{Rééducation / Réinsertion (prévention spéciale) :} La peine doit transformer le détenu pour qu'il ne récidive pas (formation, travail, suivi psychologique). \emph{C'est l'idéal du Code pénal camerounais (art. 1 : « La peine a pour but... l'amendement du condamné »)}.
\end{enumerate}

\begin{exemplebox}[Procès pour escroquerie au Tribunal de Première Instance de Yaoundé]
\textbf{Faits :} Mme \textsc{Mvondo} a détourné 50 millions FCFA de son employeur (fausses factures) sur 2 ans. Victime : entreprise, licenciement de 3 employés par ricochet.
\textbf{Jugement :}
\begin{itemize}
    \item \textbf{Voisinage rétributif :} 3 ans de prison ferme (proportionné à la gravité : abus de confiance, montant élevé, récidive légale).
    \item \textbf{Voisinage commutatif (partie civile) :} Restitution des 50 millions + 5 millions dommages-intérêts pour préjudice moral/licenciements.
    \item \textbf{Voisinage distributif (confiscation) :} Biens acquis par le produit du délit saisis et vendus au profit du Trésor public (redistribution collective).
\end{itemize}
\textbf{Analyse :} Une seule affaire mobilise les trois formes. La prison ne « répare » pas l'argent volé (commutatif), l'argent ne « punit » pas la faute morale (rétributif), la confiscation ne « compense » pas la victime individuelle (distributif).
\end{exemplebox}

\begin{exoresolu}[Analyse des formes de justice dans un jugement mixte]
\textbf{Énoncé :} Un jugement condamne un conducteur ivre ayant tué un piéton à : 2 ans de prison ferme, 10 millions FCFA à la famille de la victime, retrait du permis 5 ans, confiscation du véhicule au profit de l'État. Identifiez pour chaque mesure la forme de justice mobilisée et justifiez.
\tcblower
\textbf{Solution détaillée :}
\begin{enumerate}
    \item \textbf{2 ans prison ferme} $\to$ \textbf{Justice rétributive}. Sanction de la faute (homicide involontaire aggravé par l'ivresse). Fonction : dissuasion, rééducation, expiation.
    \item \textbf{10 millions à la famille} $\to$ \textbf{Justice commutative}. Réparation du préjudice subi (perte du soutien financier, préjudice moral). Égalité arithmétique : compensation du dommage.
    \item \textbf{Retrait permis 5 ans} $\to$ \textbf{Justice rétributive (mesure de sûreté) / distributive (protection collective)}. Mesure préventive : retire un droit (conduire) pour protéger la société (distribution de la sécurité).
    \item \textbf{Confiscation véhicule} $\to$ \textbf{Justice distributive (confiscation au profit Trésor)}. Privation d'un bien acquis/utilié pour la faute, redistribué à la collectivité. Pas de lien direct avec le dommage de la victime.
\end{enumerate}
\end{exoresolu}

\courssub{La question de la peine de mort et de la prison : la punition rend-elle justice ?}

\textbf{Le débat philosophique.}\par
\begin{itemize}
    \item \textbf{Pour la peine de mort (rétentionnistes) :} Seule peine proportionnée au crime suprême (meurtre). Dissuasion ultime. « Qui verse le sang de l'homme, par l'homme son sang sera versé » (Genèse 9,6 ; Kant, \emph{Métaphysique des mœurs} : « Il faut qu'il meure »).
    \item \textbf{Contre (abolitionnistes) :} Droit à la vie absolu (art. 3 DHUD). Irréversibilité de l'erreur judiciaire. Aucune preuve de dissuasion supérieure à la prison à vie. Peine inhumaine et dégradante (Protocole n°13 CEDH, Résolutions ONU). \emph{Cameroun : moratoire \emph{de facto} depuis 1997 (dernière exécution), mais peine maintenue dans le Code pénal (art. 231-232) pour meurtre, trahison, espionnage.}
    \item \textbf{La prison :} Michel Foucault (\emph{Surveiller et punir}) montre que la prison a remplacé le supplice public mais produit souvent \emph{plus de délinquance} (école du crime, désocialisation, stigmatisation). Surpopulation, conditions indignes $\to$ la peine devient \emph{injuste} par ses conditions d'exécution.
\end{itemize}

\begin{savaistu}[Surpopulation carcérale au Cameroun : justice réelle vs justice idéale]
Selon les rapports du MINJUSTICE et d'ONG (ACAT, Commission nationale des droits de l'homme) :
\begin{itemize}
    \item Capacité officielle des 79 prisons : \textbf{environ 17 000 places}.
    \item Population carcérale réelle (2023-2024) : \textbf{30 000 à 35 000 détenus} (dont \textbf{60 à 70 \% de prévenus} en attente de jugement, parfois des années).
    \item Taux d'occupation moyen : \textbf{180 \% à 200 \%} (pointe à 300 \% à la prison centrale de Douala, Yaoundé-Kondengui).
    \item Conséquences : promiscuité extrême, malnutrition, maladies (tuberculose, beri-béri), violences entre détenus, absence de programmes de rééducation.
\end{itemize}
\textbf{Problème philosophique :} Si la peine a pour but l'amendement (Code pénal), une prison qui détruit physiquement et psychiquement le détenu \textbf{contredit sa propre finalité}. La justice rétributive réelle (la prison telle qu'elle est) trahit la justice rétributive idéale (la peine juste, proportionnée, rééducative). C'est l'expérience vécue de l'\textbf{injustice institutionnelle}.
\end{savaistu}

\begin{popfigure}
\begin{tikzpicture}[
    scale=0.85, transform shape,
    every node/.style={font=\small},
    level 1/.style={sibling distance=5.5cm},
    level 2/.style={sibling distance=3.2cm},
    level 3/.style={sibling distance=2.5cm},
    edge from parent/.style={draw=popdark, thick, ->},
    concept/.style={draw=popblue, fill=popblueL, thick, rounded corners, minimum width=2.8cm, minimum height=1.1cm, text centered, text width=2.6cm},
    subconcept/.style={draw=poporange, fill=poporangeL, thick, rounded corners, minimum width=2.6cm, minimum height=1.0cm, text centered, text width=2.4cm},
    example/.style={draw=popgreen, fill=popgreenL, thick, rounded corners, minimum width=3.0cm, minimum height=0.9cm, text centered, text width=2.8cm, font=\footnotesize}
]
\node[concept, align=center]{\textbf{JUSTICE}\\(Aristote, Éthique V)}
    child { node[subconcept, align=center]{\textbf{Commutative}\\(Échanges particuliers)}
        child { node[example, align=center]{Ex. camerounais :\\Accident taxi Douala\\$\to$ Réparation 850k FCFA\\ (Égalité arithmétique)} }
    }
    child { node[subconcept, align=center]{\textbf{Distributive}\\(Répartition sociale)}
        child { node[example, align=center]{Ex. camerounais :\\Bourses MINESUP\\ Mérite + Besoin social\\ (Proportionnalité géométrique)} }
    }
    child { node[subconcept, align=center]{\textbf{Rétributive (Pénale)}\\(Sanction des fautes)}
        child { node[example, align=center]{Ex. camerounais :\\Escroquerie 50M TPI Yaoundé\\ Prison + Restitution + Confiscation\\ (Proportionnalité peine/faute)} }
    };
\end{tikzpicture}
\legende{Classification des trois formes de la justice (Aristote) avec un exemple concret camerounais par branche. La justice commutative restaure l'égalité arithmétique ; la distributive applique une proportionnalité géométrique selon un critère (mérite, besoin) ; la rétributive proportionne la peine à la faute.}
\end{popfigure}

\courssub{Vers l'idée absolue de justice : au-delà des lois positives}

\textbf{Le problème : la loi peut être injuste.}\par
Les trois formes ci-dessus opèrent \emph{dans} un cadre légal (Code civil, Code pénal, Code du travail). Mais que faire si la loi elle-même est injuste ? (Ex. : lois d'apartheid, lois vichystes, législation coloniale du travail forcé). La justice positive (celle des tribunaux) ne coïncide pas toujours avec la \textbf{justice idéale}.

\textbf{Platon : la justice comme harmonie de l'âme et de la Cité.}\par
Dans \emph{La République}, Platon refuse que la justice soit « l'intérêt du plus fort » (Thrasymaque) ou un simple contrat (Glaucon). La justice vraie est un \textbf{ordre ontologique} : chaque partie de l'âme (raison, courage, appétit) et chaque classe de la Cité (gardiens, auxiliaires, producteurs) accomplit \emph{sa propre fonction} sans empiéter sur l'autre. La justice absolue est cet \emph{idéal régulateur} vers lequel les lois humaines doivent tendre, sans jamais l'atteindre pleinement.

\textbf{Kant : le droit comme condition de la liberté externe.}\par
Pour Kant (\emph{Métaphysique des mœurs}, 1797), la justice n'est pas le bonheur mais la \textbf{conformité de la volonté extérieure à la loi universelle}. Le principe : « Agis extérieurement de telle sorte que le libre usage de ton arbitre puisse coexister avec la liberté de chacun selon une loi universelle. » L'État de droit (\emph{Rechtsstaat}) garantit cette coexistence. La justice absolue est l'\textbf{Idée régulatrice} (au sens kantien) : elle nous commande de réformer sans cesse le droit positif pour qu'il approche le droit pur.

\begin{aretenir}[Justice positive vs Justice idéale]
\begin{itemize}
    \item \textbf{Justice positive (droit en vigueur) :} Ce que disent les codes, ce que jugent les tribunaux aujourd'hui. Elle est \emph{relative} (variante selon les époques, les pays) et \emph{imparfaite} (erreurs, lenteurs, inégalités d'accès).
    \item \textbf{Justice idéale (idée absolue) :} Critère critique qui permet de dire « cette loi est injuste », « ce jugement est inique ». Elle est \emph{universelle} (droits de l'homme), \emph{exigeante} (exige la réforme), \emph{jamais pleinement réalisée}.
    \item \textbf{Lien :} Le juriste applique le droit positif ; le philosophe (et le citoyen) le juge à l'aune de l'idée absolue. Le progrès juridique = réduction de l'écart entre les deux.
\end{itemize}
\end{aretenir}

\courssub{L'injustice comme expérience vécue : point de départ de la réflexion}

\textbf{Phénoménologie de l'injustice (Ricoeur, Levinas).}\par
On ne découvre pas la justice par abstraction, mais par le \textbf{sentiment d'injustice} : « Cela ne va pas », « C'est pas juste ». C'est une \emph{expérience vécue} (pathos) avant d'être un concept. La victime crie justice ; le philosophe conceptualise ce cri.

\textbf{Cas camerounais emblématique : les détentions prolongées sans jugement.}\par
\begin{itemize}
    \item \textbf{Fait :} Des milliers de prévenus (60-70 \% de la population carcérale) attendent leur procès 2, 5, parfois 10 ans (affaires « dossier vide », perte de dossiers, grèves des magistrats, absence d'avocat).
    \item \textbf{Violation flagrante :} Code de procédure pénale (délais de détention provisoire), Constitution (présomption d'innocence), PIDCP (art. 9 : « toute personne arrêtée sera jugée dans un délai raisonnable ou libérée »).
    \item \textbf{Expérience vécue :} Le détenu innocent (ou non encore jugé) subit une peine (prison) \emph{sans} jugement $\to$ confusion totale des formes : ni rétributive (pas de condamnation), ni commutative (pas de réparation), ni distributive (charge pour l'État sans utilité). C'est l'\textbf{arbitraire pur}.
    \item \textbf{Réflexion philosophique :} Cette injustice crie le besoin d'une \textbf{justice transitionnelle / restauratrice} (audience foraine, libération conditionnelle massive, réforme de l'aide juridictionnelle) pour rapprocher le réel de l'idéal.
\end{itemize}

\begin{methode}[Analyser une situation concrète à l'aide des trois formes (Bac / Probatoire)]
\textbf{Étape 1 :} Identifier le \textbf{type de relation} : Échange entre deux ? $\to$ Commutative. Répartition collective ? $\to$ Distributive. Faute sanctionnée ? $\to$ Rétributive.
\textbf{Étape 2 :} Identifier le \textbf{critère de proportionnalité} exigé : Valeur pour valeur (arithmétique) ? Mérite/Besoin (géométrique) ? Gravité de la faute (proportionnalité peine/faute) ?
\textbf{Étape 3 :} Vérifier si la \textbf{décision/loi} respecte ce critère. Si non $\to$ \textbf{injustice} (préciser laquelle : commutative, distributive, rétributive).
\textbf{Étape 4 :} Situer par rapport à l'\textbf{idée absolue} : La loi elle-même est-elle juste ? (Ex. : loi autorisant détention 10 ans sans procès = loi injuste).
\textbf{Étape 5 :} Conclure en articulant \textbf{droit positif} et \textbf{idéal de justice}.
\end{methode}

\begin{exercice}[Application : Grève des agents de l'État pour arriérés de salaires]
Des agents contractuels de l'État (enseignants, infirmiers) n'ont pas été payés depuis 6 mois. Ils font grève. Le gouvernement les menace de remplacement. La presse parle d'« injustice ».
\begin{enumerate}
    \item Quelle(s) forme(s) de justice sont en jeu ? Justifiez.
    \item Le non-paiement viole-t-il la justice commutative, distributive, ou les deux ?
    \item La menace de remplacement est-elle juste au regard de la justice rétributive ?
    \item En quoi cette situation révèle-t-elle l'écart entre justice positive (contrats signés, budget voté) et justice idéale ?
\end{enumerate}
\end{exercice}

\begin{corrige}[Exercice : Grève des agents pour arriérés]
\begin{enumerate}
    \item \textbf{Commutative :} Contrat de travail = échange travail $\leftrightarrow$ salaire. Non-paiement = rupture de l'égalité arithmétique (travail donné, salaire non reçu). \textbf{Distributive :} L'État a l'obligation de répartir le budget national pour payer ses agents (critère : besoin/service rendu). L'arbitraire budgétaire viole la proportionnalité géométrique. \textbf{Rétributive (secondaire) :} Menace de sanction (remplacement) pour exercice droit de grève (liberté fondamentale) $\to$ disproportion peine/faute.
    \item \textbf{Les deux.} Commutative : dette contractuelle précise (montant dû). Distributive : obligation sociale de l'État employeur de garantir la rémunération du service public.
    \item \textbf{Non.} La grève est un droit (Constitution, OIT). Remplacer des grévistes pour réclamer un droit (salaire dû) est une sanction disproportionnée $\to$ injustice rétributive. De plus, cela aggrave l'injustice commutative (nouveaux agents payés, anciens non).
    \item \textbf{Justice positive :} Contrats existent, budget voté $\to$ obligation légale de payer. \textbf{Réalité :} Trésorerie tendue, choix politiques autres $\to$ non-paiement. \textbf{Idée absolue :} Le travail mérite salaire (justice commutative) ; l'État doit l'exemplarité (justice distributive). L'écart = \textbf{injustice structurelle} qui mine la légitimité de l'État de droit.
\end{enumerate}
\end{corrige}

\begin{popfigure}
\begin{tikzpicture}
\begin{axis}[
    width=\linewidth,
    height=9cm,
    ybar stacked,
    bar width=18pt,
    symbolic x coords={2018,2019,2020,2021,2022,2023},
    xtick=data,
    xlabel={Année},
    ylabel={Nombre de détenus (en milliers)},
    ymin=0, ymax=40,
    legend style={at={(0.5,-0.25)}, anchor=north, legend columns=2, font=\footnotesize},
    ymajorgrids=true,
    grid style=dashed,
    title style={font=\normalsize\bfseries, yshift=0.3cm},
    title={Évolution de la population carcérale au Cameroun (estimations MINJUSTICE/ONG)},
    enlarge x limits=0.15,
]
\addplot[fill=popblue!70] coordinates {(2018,22) (2019,24) (2020,26) (2021,28) (2022,30) (2023,33)};
\addplot[fill=poporange!70] coordinates {(2018,5) (2019,5) (2020,5) (2021,5) (2022,5) (2023,5)};
\legend{Prévenus (non jugés), Condamnés}
\end{axis}
% Ligne de capacité
\draw[popred, thick, dashed] ({rel axis cs:0,17/40}) -- ({rel axis cs:1,17/40}) node[right, pos=1, font=\footnotesize, popred] {Capacité théorique $\approx$ 17 000};
\end{tikzpicture}
\legende{Statistiques du système pénitentiaire camerounais : surpopulation chronique (dépassement capacitaire 180-200 \%) et part massive des prévenus (non jugés) illustrant l'injustice rétributive réelle (peine sans jugement) et l'échec de la fonction rééducative. Sources : Rapports MINJUSTICE, ACAT, CNDHL 2018-2023.}
\end{popfigure}

\begin{attention}[Erreur méthodologique majeure en dissertation]
Ne pas réduire la justice à \textbf{une seule forme} (souvent la rétributive/pénale). Un sujet sur « La justice » exige d'analyser \textbf{les trois dimensions} : commutative (réparation), distributive (partage), rétributive (sanction). Une dissertation qui ne parle que de prison ou de peine de mort rate le cœur du programme. De même, ne pas confondre \textbf{égalité arithmétique} (commutative) et \textbf{égalité proportionnelle} (distributive) : c'est la distinction conceptuelle clé attendue.
\end{attention}

\coursec{De la justice relative à l'idée absolue de justice}

Après avoir analysé les \emph{formes concrètes de justice} (commutative, distributive, réparatrice, pénale) qui structurent la vie sociale, nous devons maintenant nous demander : ces formes suffisent-elles à épuiser l'idée de justice ? L'histoire et l'expérience juridique montrent que ce qui est « juste » selon la loi positive d'un temps ou d'un lieu peut apparaître profondément injuste sous un autre angle. Cette tension entre le \textbf{droit positif} (les lois en vigueur) et l'\textbf{idéal de justice} est le moteur de l'évolution du droit. Cette section explore comment la justice relative, incarnée dans les lois historiques, tend vers une idée absolue qui la dépasse et la juge.

\courssub{La relativité des lois : l'histoire comme preuve de la contingence du droit}

\par
L'intuition première est simple : ouvrons un code pénal du XIXe siècle ou une coutume médiévale, nous y trouverons des peines de mort pour des vols mineurs, l'esclavage légalisé, l'interdiction du vote aux femmes, la criminalisation de l'homosexualité. Ce qui était \emph{légal} hier est \emph{illégal} aujourd'hui, et vice-versa. La légalité n'est pas la justice ; elle n'en est que la forme historique, contingente, révisable.

\begin{experience}[Comparaison de législations : l'évolution du statut de la femme]
\textbf{Objectif :} Mesurer l'écart entre droit positif et justice à travers le temps.
\textbf{Corpus documentaire :}
\begin{itemize}
    \item \textbf{Code Napoléon (1804), Art. 213 :} « La femme doit obéissance à son mari. »
    \item \textbf{Code du travail camerounais (1992, révisé 2023) :} Interdiction des discriminations fondées sur le sexe (Art. 2) ; égalité de rémunération pour travail de valeur égale.
    \item \textbf{Constitution du Cameroun (1996), Préambule :} « La femme et l'homme sont égaux en droits et en devoirs. »
\end{itemize}
\textbf{Protocole d'analyse :}
\begin{enumerate}
    \item Identifier la norme juridique (droit positif) à chaque époque.
    \item Évaluer sa conformité au principe d'égalité (justice distributive).
    \item Conclure sur le caractère historique et perfectible de la loi.
\end{enumerate}
\textbf{Observation :} La loi de 1804 organisait l'infériorité juridique de la femme comme une « justice » naturelle. La loi actuelle la proscrit. La justice n'a pas changé (l'égalité est un idéal constant) ; c'est la loi qui a rattrapé la justice.
\textbf{Interprétation :} La relativité des lois prouve que le droit positif n'est pas la source de la justice, mais son instrument imparfait. Si la loi \emph{faisait} la justice par sa seule existence, toute réforme serait injustifiée.
\legende{Expérience de pensée : la loi change, l'exigence de justice demeure.}
\end{experience}

\begin{popfigure}
\begin{tikzpicture}[scale=0.9, every node/.style={font=\small}]
    % Axes
    \draw[->, thick, popdark] (0,0) -- (12,0) node[right] {\textbf{Temps / Espace (Contextes historiques)}};
    \draw[->, thick, popdark] (0,0) -- (0,6) node[above] {\textbf{Degré de conformité à la Justice}};
    % Courbe idéale (Justice absolue)
    \draw[popblue, very thick, dashed, domain=0:11, samples=100] plot (\x, {5}) node[right, pos=0.95, popblue] {\textbf{Idée absolue de Justice (Idéal régulateur)}};
    % Courbes successives du Droit Positif (escalier montant)
    % Étape 1 : Code Napoléon / Esclavage
    \draw[poporange, very thick] (0.5,1) -- (2.5,1) node[midway, above=2pt, poporange] {Code Napoléon 1804};
    \draw[poporange, thick] (2.5,1) -- (2.5,2.5);
    % Étape 2 : Abolition / Suffrage universel masculin
    \draw[poporange, very thick] (2.5,2.5) -- (5,2.5) node[midway, above=2pt, poporange] {Abolition esclavage 1848 / Suffrage masculin};
    \draw[poporange, thick] (5,2.5) -- (5,3.5);
    % Étape 3 : Suffrage féminin / DUDH
    \draw[poporange, very thick] (5,3.5) -- (7.5,3.5) node[midway, above=2pt, poporange] {DUDH 1948 / Suffrage féminin};
    \draw[poporange, thick] (7.5,3.5) -- (7.5,4.2);
    % Étape 4 : Constitution Cameroun / Droits LGBTQ+ (partiel)
    \draw[poporange, very thick] (7.5,4.2) -- (10,4.2) node[midway, above=2pt, poporange] {Const. Cameroun 1996 / CEDAW};
    \draw[poporange, thick] (10,4.2) -- (10,4.7);
    % Flèche future
    \draw[poporange, very thick, ->] (10,4.7) -- (11,4.7) node[midway, above=2pt, poporange] {Progrès à venir...};
    % Zone d'écart (hachurée)
    \foreach \x in {0.5,1,...,10.5} {
        \draw[popmuted, pattern=north east lines, pattern color=popmuted!50] (\x, {5}) -- (\x+0.5, {5}) -- (\x+0.5, {1+(\x/11)*3.7}) -- (\x, {1+(\x/11)*3.7}) -- cycle;
    }
    % Légende visuelle
    \node[align=left, anchor=north west, fill=popcream, draw=popline, rounded corners, inner sep=4pt] at (0.5, 5.5) {
        \textbf{Légende :}\\
        \textcolor{popblue}{--- --- ---} Idée absolue de Justice (Inatteignable, immuable)\\
        \textcolor{poporange}{———} Droit positif historique (Relatif, perfectible)\\
        \textcolor{popmuted}{\rule{5pt}{5pt}} Écart = Injustice légale / Marge de progression
    };
\end{tikzpicture}
\legende{Schéma en escalier : le droit positif (marche par marche) s'approche asymptotiquement de l'idée absolue de justice (ligne d'horizon). L'écart hachuré représente l'injustice légale à un instant T.}
\end{popfigure}

\begin{aretenir}[Relativité du droit positif]
Le \cle{droit positif} (l'ensemble des lois en vigueur dans un État à un moment donné) est \textbf{relatif} : il varie selon les époques (diachronie) et les pays (synchronie). 
\begin{itemize}
    \item \textbf{Preuve historique :} L'esclavage, la torture, la peine de mort pour vol, l'interdiction du divorce étaient « légaux ».
    \item \textbf{Preuve géographique :} La peine de mort, l'avortement, le blasphème, l'homosexualité sont légaux ici, illégaux là-bas.
    \item \textbf{Conséquence :} La légalité ne saurait être le critère ultime de la justice. Une loi peut être \textbf{légale mais injuste}.
\end{itemize}
\end{aretenir}

\courssub{Le conflit entre loi positive et justice : Antigone, figure fondatrice}

\par
La tragédie grecque \emph{Antigone} (Sophocle, V\textsuperscript{e} s. av. J.-C.) met en scène le conflit archétypal entre la \cle{loi positive} (l'édit du tyran Créon, loi de la cité, \emph{nomos}) et la \cle{loi naturelle / non écrite} (les lois des dieux, la piété filiale, la justice immuable, \emph{dikaï}).

\begin{exemplebox}[Antigone : le droit naturel contre la loi positive]
\textbf{Situation :} Créon, roi de Thèbes, interdit sous peine de mort d'enterrer Polynice (traître à la patrie). Antigone, sa sœur, lui désobéit au nom des « lois non écrites et immuables des dieux ».
\textbf{Le dialogue clé (v. 450-457) :}
\begin{quote}
\textsc{Créon.} — Tu as osé transgresser cette loi ?\\
\textsc{Antigone.} — Car ce n'est pas Zeus qui me l'a proclamée, ni la Justice qui siège avec les dieux d'en bas ; [...] ce ne sont pas d'aujourd'hui ni d'hier, c'est de toujours qu'elles vivent, et nul ne sait d'où elles sont apparues.
\end{quote}
\textbf{Analyse philosophique :}
\begin{itemize}
    \item \textbf{Créon incarne le légalisme pur :} « Nul ne doit désobéir à la loi de la cité. » L'État est la source unique du droit. La justice = obéissance.
    \item \textbf{Antigone incarne le jusnaturalisme :} Il existe une justice \textbf{antérieure et supérieure} à la loi humaine. La loi positive n'est légitime que si elle est conforme à cette justice supérieure.
    \item \textbf{Le tragique :} Antigone meurt, mais l'histoire (et la philosophie) lui donne raison : Créon perd son fils, sa femme, son pouvoir. La loi positive qui viole la justice naturelle se détruit elle-même.
\end{itemize}
\legende{Antigone enterre son frère : l'acte de désobéissance civile fondatrice.}
\end{exemplebox}

\begin{propriete}[L'idée absolue de justice comme idéal régulateur (Kant)]
L'idée absolue de justice n'est pas un état de fait réalisable parfaitement dans l'histoire (ce n'est pas une \emph{constitution} empirique), mais un \textbf{idéal régulateur} (\emph{Regulativbegriff}).
\begin{itemize}
    \item \textbf{Fonction :} Elle ne se \emph{réalise} jamais pleinement (pas de société parfaitement juste), mais elle \emph{oriente} : elle sert de critère pour juger les lois existantes et guider les réformes.
    \item \textbf{Analogie :} Comme l'horizon guide le marcheur sans jamais être atteint, l'idée absolue de justice guide le législateur et le juge.
    \item \textbf{Conséquence :} Toute loi positive est \textbf{provisoire} ; elle est un essai d'approximation de la justice. La critique du droit est donc permanente.
\end{itemize}
\legende{Propriété fondamentale : l'idéal régule le réel sans s'identifier à lui. Notion clé pour l'épreuve du Probatoire : expliquer la distinction constitutive/régulateur chez Kant.}
\end{propriete}

\courssub{Le rôle du juge : dire le droit tout en visant l'équité}

\par
Si la loi est le texte, le juge est l'interprète. Face à la généralité abstraite de la loi et à la singularité concrète du cas, le juge ne peut être une « bouche de la loi » (Montesquieu) sans risque d'injustice.

\begin{definition}[Équité (\emph{aequitas})]
L'\cle{équité} est la correction de la loi là où elle est défectueuse par sa généralité. Elle permet au juge d'adapter la règle universelle au cas particulier pour atteindre la justice matérielle, sans nier la loi formelle.
\begin{itemize}
    \item \textbf{Aristote, \emph{Éthique à Nicomaque} (V, 10) :} « L'équitable est juste, mais non pas le juste légal, plutôt la correction du juste légal. »
    \item \textbf{Application :} Circonstances atténuantes, interprétation extensive/restrictive, principe de proportionnalité de la peine.
\end{itemize}
\end{definition}

\begin{attention}[Piège : Confondre équité et arbitraire judiciaire]
\begin{itemize}
    \item \textbf{Équité :} Respect de l'esprit de la loi, motivation juridique, prévisibilité, contrôle en appel/cassation. C'est une \textbf{justice supérieure} dans le cadre du droit.
    \item \textbf{Arbitraire :} Décision sans fondement légal, selon le bon plaisir du juge, imprévisible. C'est la \textbf{négation du droit}.
    \item \textbf{Au Probatoire :} Ne pas écrire « le juge fait ce qu'il veut ». Écrire : « Le juge interprète la loi à la lumière des principes supérieurs (équité, droits de l'homme) pour combler les lacunes du texte. »
\end{itemize}
\end{attention}

\courssub{La désobéissance civile : résister à une loi injuste}

\par
Lorsque l'écart entre la loi positive et l'idée de justice devient insupportable (lois racistes, totalitaires, attentatoires aux droits fondamentaux), l'obéissance passive devient complicité. La \cle{désobéissance civile} est la réponse politique et morale moderne au dilemme d'Antigone.

\begin{definition}[Désobéissance civile]
La \cle{désobéissance civile} est une \textbf{transgression publique, non violente, consciente et assumée} d'une loi jugée injuste, dans le but de provoquer son abrogation ou sa réforme, en acceptant les sanctions légales.
\begin{itemize}
    \item \textbf{Publique :} Pas de clandestinité ; acte politique adressé à l'opinion et au pouvoir.
    \item \textbf{Non violente :} Respect de l'intégrité physique des personnes (distinction cruciale avec l'insurrection/révolte armée).
    \item \textbf{Assumée :} Le désobéissant accepte la prison pour manifester la supériorité de sa conscience sur la loi inique.
    \item \textbf{Finalité :} Réformer la loi, pas détruire l'État de droit.
\end{itemize}
\end{definition}

\begin{exemplebox}[Gandhi et Martin Luther King : l'éthique de la résistance]
\begin{center}
\begin{tabularx}{\linewidth}{|l|X|X|}
\hline
\rowcolor{popblueL} \textbf{Critère} & \textbf{Gandhi (Inde, \emph{Satyagraha})} & \textbf{Martin Luther King (USA, Droits civiques)} \\
\hline
\textbf{Loi injuste ciblée} & Lois coloniales britanniques (sel, coton, vote) & Lois de ségrégation raciale (Jim Crow) \\
\hline
\textbf{Fondement moral} & \emph{Ahimsa} (non-violence) + Loi de la conscience supérieure à la loi de l'État & Loi naturelle / Loi de Dieu > Loi positive ségrégationniste (Lettre de Birmingham) \\
\hline
\textbf{Moyens} & Marche du sel, boycott, jeûne, non-coopération & Sit-ins, marches, boycott bus Montgomery \\
\hline
\textbf{Acceptation sanction} & Prison (plusieurs fois) : « Aller en prison est un honneur » & Prison : « On a une responsabilité morale de désobéir aux lois injustes » \\
\hline
\textbf{Résultat} & Indépendance de l'Inde (1947) & Civil Rights Act (1964), Voting Rights Act (1965) \\
\hline
\end{tabularx}
\end{center}
\legende{Deux modèles historiques de désobéissance civile réussie : la force morale contre la force légale.}
\end{exemplebox}

\begin{attention}[Erreur fréquente : Désobéissance civile $\neq$ Révolte violente / Émeute]
\begin{itemize}
    \item \textbf{Désobéissance civile :} Cadre démocratique (ou visé), non-violence stricte, visée \textbf{juridique} (changer la loi), acceptation du procès. \textbf{Exemple :} Manifestants climat bloquant une route, acceptant l'amende.
    \item \textbf{Révolte / Insurrection :} Violence physique, visée \textbf{politique} (renverser le pouvoir), clandestinité ou affrontement. \textbf{Exemple :} Prise d'armes contre un régime dictatorial (droit de résistance à l'oppression, Préambule Constitution Cameroun).
    \item \textbf{Nuance Probatoire :} La désobéissance civile \emph{présuppose} un minimum d'État de droit (liberté d'expression, procès équitable) pour s'exprimer. En dictature totale, elle bascule en résistance/insurrection.
\end{itemize}
\end{attention}

\courssub{Synthèse : Droit sans justice = arbitraire ; Justice sans droit = vengeance privée}

\par
La leçon toute entière converge vers cette double vérité dialectique :

\begin{propriete}[L'indissociabilité du Droit et de la Justice]
\begin{center}
\begin{tabularx}{\linewidth}{|>{\bfseries}X|X|X|}
\hline
\rowcolor{popblueL} \textbf{Configuration} & \textbf{Analyse} & \textbf{Risque / Dérive} \\
\hline
\textbf{Droit \emph{sans} Justice} & Loi positive déconnectée de l'équité, de l'humanité, des droits fondamentaux. & \textbf{Arbitraire légal / Totalitarisme légal.} « Tout est permis à l'État car c'est la loi. » (Ex. : lois nazies, apartheid, esclavage légal). \\
\hline
\textbf{Justice \emph{sans} Droit} & Revendication subjective de « ma justice » sans procédure, sans juge tiers, sans règle commune. & \textbf{Vengeance privée / Lynchage / État de nature.} « Je me fais justice moi-même. » Fin de la sécurité juridique, guerre de tous contre tous. \\
\hline
\textbf{Droit \textbf{et} Justice} & Droit positif constamment critiqué, réformé, interprété à la lumière de l'idéal de justice (Droits de l'homme, équité). & \textbf{État de droit démocratique.} Le droit est \textbf{instrument} de la justice ; la justice est \textbf{fin} du droit. \\
\hline
\end{tabularx}
\end{center}
\legende{Tableau de synthèse : le droit est la forme, la justice est l'esprit. Sans esprit, la forme est vide (arbitraire) ; sans forme, l'esprit est destructeur (vengeance).}
\end{propriete}

\courssub{Bilan de la leçon et ouverture vers la question de l'État}

\begin{methode}[Conclure une dissertation sur le Droit et la Justice]
\textbf{Sujet type :} « La loi suffit-elle à définir la justice ? » ou « Peut-on désobéir à la loi au nom de la justice ? »
\begin{enumerate}
    \item \textbf{Rappeler la tension :} Reprendre la distinction fondamentale : Droit positif (fait, relatif, coercitif) vs Idée de Justice (valeur, absolue, idéale).
    \item \textbf{Nuire au simplisme :} Refuser les deux écueils : légalisme béat (loi = justice) et subjectivisme anarchique (ma conscience = justice).
    \item \textbf{Proposer la synthèse dialectique :} Le droit \textbf{a besoin} de la justice pour être légitime (sinon arbitraire) ; la justice \textbf{a besoin} du droit pour être effective et paisible (sinon vengeance).
    \item \textbf{Ouvrir sur l'institution garante :} C'est l'\textbf{État de droit} (séparation des pouvoirs, justice indépendante, constitution, contrôle de constitutionnalité) qui institutionnalise cette tension pour la rendre féconde et non explosive.
    \item \textbf{Conclusion ouverte :} La justice n'est pas un trésor qu'on possède, mais un chantier permanent. La démocratie, c'est l'institutionnalisation du conflit droit/justice par le débat public et le juge constitutionnel.
\end{enumerate}
\end{methode}

\begin{savaistu}[Le Cameroun face à l'idéal de justice]
\begin{itemize}
    \item \textbf{Préambule Constitution 1996 :} Rattache le Cameroun aux « idéaux de la Déclaration universelle des droits de l'homme » (référence à l'absolu).
    \item \textbf{Conseil Constitutionnel (créé 2018, opérationnel 2020) :} Juge de la conformité des lois à la Constitution (garde-fou contre l'arbitraire légal).
    \item \textbf{Problématique actuelle :} Peine de mort (maintenue dans le Code pénal mais moratoire \emph{de facto} depuis 1997) ; Droit coutumier vs Droit moderne (ex. : succession, mariage) ; Indépendance réelle du pouvoir judiciaire vs exécutif.
    \item \textbf{Ouverture :} La question « Qu'est-ce qu'un État juste ? » (Programme de Terminale, Philosophie politique) prend ici tout son sens : l'État n'est pas seulement l'appareil qui fait la loi (légalité), mais celui qui garantit que la loi tend vers la justice (légitimité).
\end{itemize}
\end{savaistu}

\begin{popfigure}
\begin{tikzpicture}[scale=0.85, every node/.style={font=\small}]
    % Frise chronologique horizontale
    \draw[thick, popdark] (0,0) -- (14,0);
    % Graduations et événements
    \def\events{
        -0.5/Ve s. av. J.-C./Antigone (Sophocle)/Naissance du conflit Loi vs Justice/Conflit fondateur,
        3.5/1789/Déclaration DDHC/Droits naturels et résistance à l'oppression (Art. 2)/Positivisme révolutionnaire,
        7/1948/DUDH (ONU)/Droits universels et supranationaux/Justice comme norme internationale,
        10.5/1996/Constitution Cameroun/Préambule DUDH et Conseil Constitutionnel/Ancrage national de l'idéal,
        14/Aujourd'hui/Défis/Peine de mort et justice transitionnelle et État de droit/Horizon ouvert}
    \foreach \x/\date/\event/\desc/\nature in \events {
        % Point sur la ligne
        \fill[popblue] (\x,0) circle (3pt);
        % Date au-dessus
        \node[above=4pt, anchor=south, popdark, font=\bfseries\small] at (\x,0.6) {\date};
        % Événement
        \node[above=14pt, anchor=south, text width=2.5cm, align=center, popblue] at (\x,1.2) {\event};
        % Description (alternance haut/bas pour lisibilité)
        \ifdim\x pt<7pt
            \node[below=14pt, anchor=north, text width=2.5cm, align=center, popmuted] at (\x,-1.2) {\desc};
            \node[below=38pt, anchor=north, text width=2.5cm, align=center, poporange!80!popdark, font=\itshape\small] at (\x,-2.0) {\nature};
        \else
            \node[above=38pt, anchor=south, text width=2.5cm, align=center, popmuted] at (\x,2.6) {\desc};
            \node[above=62pt, anchor=south, text width=2.5cm, align=center, poporange!80!popdark, font=\itshape\small] at (\x,3.4) {\nature};
        \fi
        % Trait de repère
        \draw[popline, thin] (\x,0) -- (\x,0.4);
    }
    % Flèches de progression
    \draw[popgreen, thick, ->] (0.5, -3.2) -- (13.5, -3.2) node[midway, below=2pt] {Progrès de la reconnaissance juridique de la justice (Relatif $\to$ Absolu)};
    % Titre
    \node[above=8pt, anchor=south, popdark, font=\bfseries\large] at (7, 4.5) {Frise : De la tragédie antique à l'État de droit camerounais};
\end{tikzpicture}
\legende{Frise chronologique : l'histoire du droit est l'histoire de la réduction progressive de l'écart entre le légal et le juste. Le Cameroun (1996) s'inscrit dans cette dynamique universaliste post-1948.}
\end{popfigure}

\begin{exoresolu}[Sujet de dissertation type Probatoire]Énoncé : « Respecter la loi, est-ce toujours agir justement ? »\tcblower
\textbf{Analyse du sujet :}
\begin{itemize}
    \item \textbf{Termes :} « Respecter la loi » = obéissance légale, conformité formelle. « Agir justement » = conformité à l'équité, à la morale, à l'idée absolue de justice.
    \item \textbf{Problématique :} La légalité est-elle condition suffisante de la légitimité morale ? La coïncidence droit/justice est-elle structurelle ou accidentelle ?
\end{itemize}

\textbf{Plan dialectique :}

\textbf{I. Thèse : La loi comme condition nécessaire de la justice (Légalisme).}
\begin{itemize}
    \item A. La loi assure l'égalité formelle : même règle pour tous (Kelsen, pure théorie du droit). Sans loi, arbitraire du plus fort.
    \item B. La sécurité juridique : prévisibilité, stabilité, paix civile. Respecter la loi, c'est respecter le pacte social (Rousseau, Contrat social).
    \item C. Le juge comme « bouche de la loi » (Montesquieu) : garantit l'impartialité contre le subjectivisme.
\end{itemize}

\textbf{II. Antithèse : La loi peut être injuste (Relativité / Jusnaturalisme).}
\begin{itemize}
    \item A. Argument historique : lois esclavagistes, ségrégationnistes, totalitaires étaient « légales » (Antigone, Nuremberg).
    \item B. Argument structurel : la loi est générale, le cas est singulier $\to$ risque d'injustice formelle (Aristote, équité).
    \item C. La distinction légalité/légitimité : une loi inique (contraire aux droits de l'homme) perd son caractère obligatoire en conscience (Thoreau, Gandhi, King).
\end{itemize}

\textbf{III. Synthèse : L'État de droit, institution de la tension droit/justice.}
\begin{itemize}
    \item A. Le droit n'est pas la justice, mais son \textbf{instrument}. Il doit être \textbf{justiciable} (contrôle de constitutionnalité, conventionnalité).
    \item B. Le citoyen a le \textbf{devoir d'obéissance} aux lois justes, et le \textbf{droit/devoir de résistance} (désobéissance civile) aux lois manifestement contraires aux droits fondamentaux (Préambule Const. Cameroun, Art. 2 DDHC 1789).
    \item C. La justice n'est pas un état, mais un \textbf{processus} : le débat démocratique, le juge indépendant, la révision constitutionnelle sont les mécanismes qui font tendre le droit vers la justice.
\end{itemize}

\textbf{Conclusion :} Respecter la loi, c'est agir \emph{légalement} ; agir justement, c'est agir \emph{légitimement}. La coïncidence n'est pas donnée, elle est \textbf{à conquérir sans cesse} par la vigilance citoyenne et l'exigence constitutionnelle.
\end{exoresolu}

\begin{exercice}[Entraînement Probatoire - Commentaire de texte]
\textbf{Texte :} « Il n'y a pas de justice sans loi, mais il y a des lois sans justice. » (Inspiré de Montesquieu / Kant)
\textbf{Consignes :}
\begin{enumerate}
    \item Expliquez la distinction opérée entre « loi » et « justice ».
    \item Montrez pourquoi « il y a des lois sans justice » (donnez deux exemples historiques ou actuels).
    \item En quoi l'État de droit est-il la réponse institutionnelle à cette tension ?
\end{enumerate}
\end{exercice}

\begin{corrige}[Exercice n°1 - Éléments de correction]
\begin{enumerate}
    \item \textbf{Distinction :} La \textbf{loi} (droit positif) est une norme formelle, écrite, coercitive, édictée par le souverain/législateur, relative à un temps/lieu. La \textbf{justice} est une valeur morale idéale (égalité, équité, droit naturel), universelle, inconditionnelle, qui sert de critère de validité morale à la loi.
    \item \textbf{Exemples de lois sans justice :}
        \begin{itemize}
            \item \textbf{Lois de Nuremberg (1935, Allemagne) :} Légalisent l'antisémitisme, privent les Juifs de citoyenneté. Légales, radicalement injustes.
            \item \textbf{Code de l'indigénat (Colonies françaises, jusqu'en 1946) :} Justice sommaire, travail forcé, pas d'égalité devant la loi pour les « sujets » vs « citoyens ». Légal colonial, injuste humainement.
            \item \textbf{Exemple camerounais (débattu) :} Peine de mort maintenue dans le Code pénal (Art. 222-1) alors que le Préambule 1996 renvoie à la DUDH (droit à la vie) et que le moratoire est respecté. Incohérence droit/justice.
        \end{itemize}
    \item \textbf{L'État de droit comme réponse :}
        \begin{itemize}
            \item \textbf{Séparation des pouvoirs :} Le législateur fait la loi, mais le juge (judiciaire + constitutionnel) la contrôle au regard de la Constitution (bloc de constitutionnalité = justice formelle).
            \item \textbf{Hiérarchie des normes :} Loi < Traité international (DUDH) < Constitution. Le juge peut écarter la loi injuste (inconstitutionnalité / non-conventionnalité).
            \item \textbf{Procédure :} Le justiciable a accès à un juge indépendant pour faire valoir ses droits (accès au droit, aide juridictionnelle).
            \item \textbf{Démocratie :} Le peuple peut changer la loi injuste par le vote, le référendum, la pétition, la manifestation (liberté d'expression).
        \end{itemize}
\end{enumerate}
\end{corrige}

\coursec{Méthode et application aux épreuves}

Après avoir exploré le chemin qui mène des formes concrètes de justice à l'idée absolue de justice, il nous faut maintenant transformer ces savoirs en compétences éprouvées à l'examen. Le Probatoire et le Baccalauréat technique ne demandent pas de réciter le cours, mais de \emph{mobiliser} les concepts (droit, justice, égalité, équité, légalité, légitimité) pour analyser un texte ou construire une argumentation rigoureuse. Cette section vous donne les clés méthodologiques et un entraînement ciblé.

\courssub{Méthode de l'explication de texte philosophique}

L'explication de texte n'est pas un résumé ni une paraphrase. C'est un exercice d'\emph{analyse} qui vise à faire apparaître la pensée de l'auteur dans sa structure et sa portée.

\begin{methode}[Les quatre temps de l'explication de texte]
\begin{enumerate}
\item \textbf{Introduction (environ 10 lignes) :} Présenter l'auteur, l'œuvre, le contexte général, \textbf{la thèse principale} (en une phrase), le \textbf{problème} traité et l'\textbf{annonce du plan} de l'explication (qui suit le fil du texte).
\item \textbf{Analyse détaillée (corps) :} Suivre le texte paragraphe par paragraphe (ou idée par idée). Pour chaque étape : \emph{citer} (mots-clés ou courte phrase), \emph{expliquer} (définir les concepts, montrer le raisonnement), \emph{commenter} (porter un regard critique : la thèse est-elle solide ? Y a-t-il des présupposés ? Des distinctions utiles à faire ?).
\item \textbf{Synthèse (environ 5-8 lignes) :} Reprendre la thèse, montrer comment l'argumentation s'articule, indiquer la portée philosophique du texte (à quel débat classique répond-il ?).
\item \textbf{Ouverture (2-3 lignes) :} Une question ou une référence brève qui prolonge la réflexion (autre auteur, situation contemporaine).
\end{enumerate}
\end{methode}

\begin{attention}[Piège fréquent : le « commentaire de texte » hors sujet]
Ne pas confondre explication et dissertation. L'explication \textbf{reste collée au texte} : on ne développe pas ses propres arguments, on expose et évalue ceux de l'auteur. L'erreur majeure au Probatoire et au Baccalauréat est de réciter son cours sur la justice sans montrer comment \emph{ce texte précis} la traite. Chaque paragraphe de votre analyse doit renvoyer explicitement à une citation ou une idée du texte.
\end{attention}

\begin{exemplebox}[Repérer les concepts opératoires dans un texte]
Dans un extrait sur la justice, soulignez systématiquement : \cle{droit}, \cle{justice}, \cle{égalité}, \cle{équité}, \cle{loi}, \cle{légitimité}, \cle{distributive}, \cle{corrective}, \cle{proportionnalité}. Ce sont les « charnières » du raisonnement. Si l'auteur oppose « justice légale » et « justice naturelle », c'est l'opposition \cle{légal / légitime} qui structure le texte.
\end{exemplebox}

\courssub{Méthode de la dissertation philosophique}

La dissertation répond à une \textbf{question problématique} (ex. : « La loi est-elle toujours juste ? »). Elle exige une \emph{problématisation} (montrer pourquoi la question n'a pas de réponse évidente) et une \emph{argumentation dialectique} (thèse, antithèse, synthèse).

\begin{methode}[Problématiser à partir d'une opposition conceptuelle]
\begin{enumerate}
\item \textbf{Repérer les concepts en tension :} Dans « La loi est-elle toujours juste ? », les concepts sont \cle{loi} (ordre positif, contrainte, généralité) et \cle{justice} (idéal moral, équité, proportionnalité).
\item \textbf{Identifier l'opposition :} \emph{Légal} (ce qui est conforme à la loi) $\neq$ \emph{Légitime} (ce qui est fondé moralement). \emph{Égalité arithmétique} (même traitement pour tous) $\neq$ \emph{Équité / proportionnalité géométrique} (traitement proportionné au mérite ou au besoin).
\item \textbf{Transformer en question philosophique :} « La conformité à la loi suffit-elle à garantir la justice ? » ou « L'obéissance à la loi peut-elle entrer en conflit avec les exigences de la justice ? »
\item \textbf{Construire le plan dialectique :}
\begin{itemize}
    \item \textbf{Thèse :} La loi \emph{est} la justice (positivisme juridique : Kelsen, Hart). La loi assure l'égalité formelle, la sécurité juridique, la prévisibilité.
    \item \textbf{Antithèse :} La loi \emph{peut être injuste} (jusnaturalisme : Antigone, Luther King, Rawls). Loi inique, discrimination légale, « légalité » $\neq$ « légitimité ». L'équité corrige la rigueur de la loi (Aristote, \emph{Éthique à Nicomaque} V).
    \item \textbf{Synthèse :} La loi \emph{vise} la justice mais ne l'\emph{épuise} pas. Nécessité d'un contrôle de constitutionnalité, de droits fondamentaux supérieurs à la loi, de l'équité comme correctif (équité du juge, esprit des lois). La justice est l'horizon régulateur de la loi.
\end{itemize}
\end{enumerate}
\end{methode}

\begin{popfigure}
\begin{tikzpicture}[
    node distance=1.2cm and 2.5cm,
    box/.style={draw=popblue, thick, rounded corners, fill=popblue!5, minimum width=3.2cm, minimum height=1.1cm, align=center, font=\small},
    arrow/.style={-Stealth, thick, popdark},
    labelbox/.style={font=\tiny, text=popmuted, align=center}
]
\node[box, align=center] (intro){Introduction\\ \textit{Accroche, définitions,}\\\textit{ problématique, annonce du plan}};
\node[box, below=of intro, align=center] (these){Thèse\\ \textit{La loi = justice}\\\textit{ (sécurité, égalité formelle)}};
\node[box, below=of these, align=center] (antithese){Antithèse\\ \textit{La loi $\neq$ justice}\\\textit{ (lois iniques, équité)}};
\node[box, below=of antithese, align=center] (synthese){Synthèse\\ \textit{Loi $\to$ horizon de justice}\\\textit{ contrôle, droits fondamentaux}};
\node[box, below=of synthese, align=center] (conclu){Conclusion\\ \textit{Réponse nuancée,}\\\textit{ ouverture}};
\draw[arrow] (intro) -- (these) node[midway, right, labelbox, align=center]{Problématique :\\ « La loi suffit-elle ? »};
\draw[arrow] (these) -- (antithese) node[midway, right, labelbox, align=center]{Objection :\\ lois injustes};
\draw[arrow] (antithese) -- (synthese) node[midway, right, labelbox, align=center]{Dépassement :\\ équité, contrôle};
\draw[arrow] (synthese) -- (conclu) node[midway, right, labelbox] {Bilan + ouverture};
\end{tikzpicture}
\legende{Schéma de la structure dialectique d'une dissertation (thèse -- antithèse -- synthèse). Chaque flèche représente un mouvement de la pensée : affirmation, objection, dépassement.}
\end{popfigure}

\begin{methode}[Rédiger et justifier une conclusion argumentée]
\begin{enumerate}
\item \textbf{Reprendre la problématique} (ne pas la copier-coller, la reformuler).
\item \textbf{Donner la réponse synthétique} : « La loi est une condition nécessaire mais non suffisante de la justice. »
\item \textbf{Justifier en une phrase le chemin parcouru} : « Si la loi garantit l'égalité formelle et la paix civile (thèse), elle peut devenir instrument d'oppression sans le contrôle de principes supérieurs et le correctif de l'équité (antithèse) ; c'est pourquoi tout État de droit institue des garde-fous (constitution, cours suprêmes, droits de l'homme) pour rapprocher sans cesse le légal du légitime (synthèse). »
\item \textbf{Ouvrir} : une question sur l'effectivité (accès au droit, corruption), une référence (Antigone, Déclaration 1789, Constitution camerounaise), ou un lien avec l'actualité.
\end{enumerate}
\end{methode}

\courssub{Entraînement guidé : analyse d'un extrait d'Aristote (\emph{Éthique à Nicomaque}, Livre V)}

\begin{exoresolu}[Analyse d'un extrait d'Aristote sur la justice distributive]
\textbf{Extrait (adapté, Éthique à Nicomaque, V, 3, 1131a-1131b) :}
\begin{quote}
« La justice distributive consiste dans la répartition des honneurs, de l'argent ou des autres biens qui se partagent entre les membres de la cité ; car c'est dans ce partage que l'inégalité ou l'égalité se manifeste. Or le juste dans la distribution doit être selon une proportion géométrique ; car c'est ainsi que le tout se rapporte au tout, et la partie à la partie. Si la distribution est faite selon une proportion arithmétique, il y a injustice ; car on donne une part égale à des inégaux. Le juste est donc une proportion, et l'injuste est ce qui viole la proportion. »
\end{quote}

\textbf{Éléments d'analyse (structure type Bac) :}

\textbf{Introduction} : Aristote (384-322 av. J.-C.), philosophe grec, fonde la philosophie politique sur l'étude de la cité. Dans l'\emph{Éthique à Nicomaque} (Livre V), il distingue justice générale (légalité) et justice particulière (équité). Le texte porte sur la \textbf{justice distributive} : comment répartir les biens communs ? \textbf{Thèse} : La justice distributive exige une \textbf{proportionnalité géométrique} (selon le mérite), non une égalité arithmétique (même part pour tous). \textbf{Problème} : Quel critère de répartition assure la justice dans le partage des biens sociaux ? \textbf{Plan} : 1) Définition de la justice distributive ; 2) La proportion géométrique comme critère ; 3) La critique de l'égalité arithmétique ; 4) Portée : la justice comme proportion.

\textbf{Analyse} :
\begin{itemize}
    \item \textbf{§1 (Définition) :} « La justice distributive consiste dans la répartition... » $\to$ Aristote circonscrit le domaine : honneurs, argent, biens partageables. \cle{Justice distributive} $\neq$ justice corrective (échanges, contrats). Ici, la cité distribue selon un critère de \cle{mérite} (valeur civique, vertu, richesse selon les régimes).
    \item \textbf{§2 (Proportion géométrique) :} « Le juste... doit être selon une proportion géométrique. » $\to$ $A/B = C/D$ (ex. : mérite 1 / part 1 = mérite 2 / part 2). Si le citoyen A a 2 fois le mérite de B, il reçoit 2 fois sa part. C'est l'\cle{équité} (proportionnalité) opposée à l'égalité arithmétique ($A-B = C-D$).
    \item \textbf{§3 (Critique de l'arithmétique) :} « Si la distribution est faite selon une proportion arithmétique, il y a injustice. » $\to$ Donner part égale à des inégaux (même salaire pour travail inégal, même honneur pour vertu inégale) \emph{est} une injustice. L'égalité formelle devient inéquitable si elle ignore les différences pertinentes.
    \item \textbf{§4 (Définition formelle) :} « Le juste est donc une proportion. » $\to$ La justice \emph{est} structurellement proportionnalité. L'injustice = violation de la proportion (donner trop à qui mérite peu, trop peu à qui mérite beaucoup).
\end{itemize}

\textbf{Synthèse} : Aristote fonde la justice distributive sur la \textbf{proportionnalité géométrique} (mérite $\leftrightarrow$ part). L'égalité arithmétique (même traitement) n'est juste que si les sujets sont égaux en mérite. La justice n'est pas l'uniformité, mais l'adéquation proportionnée.

\textbf{Ouverture} : Rawls (\emph{Théorie de la justice}, 1971) reprend l'idée de proportionnalité mais remplace le « mérite aristocratique » par les « avantages sociaux » à répartir au profit des plus défavorisés (principe de différence). Au Cameroun, la répartition des postes publics ou des bourses pose la question du critère : mérite scolaire, quota régional, parité genre ? Quelle proportion géométrique choisir ?
\tcblower
\textbf{Solution détaillée (attendus Bac) :}
\begin{itemize}
    \item \textbf{Introduction complète} (auteur, œuvre, thèse, problème, plan) : 4/4 pts.
    \item \textbf{Analyse} : 4 étapes clairement séparées, chaque étape cite + explique + commente. Concepts clés définis (justice distributive, proportion géométrique/arithmétique, mérite, équité). 10/10 pts.
    \item \textbf{Synthèse} : reprend thèse + articulation + définition formelle. 3/3 pts.
    \item \textbf{Ouverture} : pertinente (Rawls, actualité camerounaise). 3/3 pts.
    \item \textbf{Qualité de la langue} : termes philosophiques précis, connecteurs logiques (donc, or, car, en revanche). 5/5 pts.
\end{itemize}
\textbf{Total : 25/25.}
\end{exoresolu}

\courssub{TP : Construction d'un tableau synoptique « Formes de justice × Exemples camerounais × Philosophes »}

Ce travail pratique vise à \emph{ancrer} les concepts dans le réel camerounais et à préparer des exemples « prêts à l'emploi » pour la dissertation.

\begin{methode}[Comment remplir le tableau synoptique]
\begin{enumerate}
\item \textbf{Colonne 1 : Formes de justice} (cours §4) : Justice légale / Justice naturelle -- Justice distributive / Justice corrective (commutative) -- Justice rétributive / Justice réparatrice -- Équité (correctif du juge).
\item \textbf{Colonne 2 : Exemple concret camerounais} : Situation réelle, loi, jurisprudence, fait social, coutume.
\item \textbf{Colonne 3 : Philosophe / Référence} : Auteur qui théorise ou illustre cette forme.
\item \textbf{Colonne 4 : Concept-clé / Critère} : Égalité arithmétique, proportion géométrique, responsabilité, restauration, esprit des lois.
\end{enumerate}
\end{methode}

\begin{popfigure}
\begin{tikzpicture}
\begin{scope}[
    cell/.style={draw=popline, minimum height=0.9cm, align=center, font=\small, inner sep=2pt},
    header/.style={cell, fill=popblue, font=\small\bfseries, text=white},
    rowA/.style={cell, fill=popblue!5},
    rowB/.style={cell, fill=popcream},
]
% En-têtes
\node[header, minimum width=3.2cm] (h1) at (0,0) {Forme de justice};
\node[header, minimum width=4.5cm] (h2) at (3.2,0) {Exemple camerounais (concret)};
\node[header, minimum width=3.5cm] (h3) at (7.7,0) {Philosophe / Référence};
\node[header, minimum width=3.0cm] (h4) at (11.2,0) {Concept-clé / Critère};
% Lignes vides à compléter (6 lignes)
\foreach \i/\rowstyle in {1/rowA, 2/rowB, 3/rowA, 4/rowB, 5/rowA, 6/rowB} {
    \node[\rowstyle, minimum width=3.2cm] at (0,-\i*0.9) {};
    \node[\rowstyle, minimum width=4.5cm] at (3.2,-\i*0.9) {};
    \node[\rowstyle, minimum width=3.5cm] at (7.7,-\i*0.9) {};
    \node[\rowstyle, minimum width=3.0cm] at (11.2,-\i*0.9) {};
}
% Grille
\draw[popline] (0,0) rectangle (14.2,-5.4);
\draw[popline] (3.2,0) -- (3.2,-5.4);
\draw[popline] (7.7,0) -- (7.7,-5.4);
\draw[popline] (11.2,0) -- (11.2,-5.4);
\foreach \i in {1,...,5} {
    \draw[popline] (0,-\i*0.9) -- (14.2,-\i*0.9);
}
\end{scope}
\end{tikzpicture}
\legende{Tableau synoptique vierge à compléter (TP). Imprimez-le ou recopiez-le sur une feuille A4 paysage. Remplissez au moins 6 lignes avec des exemples variés (droit coutumier, code du travail, jurisprudence Cour suprême, allocation bourses, réparation préjudice, médiation traditionnelle).}
\end{popfigure}

\begin{exemplebox}[Lignes types pour démarrer le tableau (ne pas copier, adapter)]
\begin{center}
\begin{tabularx}{\linewidth}{|l|X|X|l|}\hline
\textbf{Forme} & \textbf{Exemple camerounais} & \textbf{Philosophe} & \textbf{Critère} \\\hline
Justice distributive & Attribution des bourses universitaires (mérite + quota régional) & Aristote, Rawls & Proportion géométrique (mérite / besoin) \\\hline
Justice corrective & Code du travail : indemnité de licenciement abusif (Art. 37) & Aristote, Code civil & Égalité arithmétique (réparation = dommage) \\\hline
Équité (correctif) & Juge des référés suspendant une expulsion locative pour raison humanitaire (esprit de la loi) & Aristote, Montesquieu & Discrétion du juge, proportionnalité \\\hline
Justice réparatrice & Médiation traditionnelle (chef de village) : auteur + victime + communauté $\to$ réconciliation & Paul Ricoeur, Droit coutumier & Restauration du lien social, pardon \\\hline
Justice transitionnelle & Commission Vérité, Justice et Réconciliation (CVJR, 2019-2021) & Hannah Arendt, J. Habermas & Reconnaissance, mémoire, non-répétition \\\hline
Légalité vs Légitimité & Loi sur la décentralisation (2019) : texte voté mais transfert de compétences incomplet $\to$ légalité sans effectivité & Kelsen (légalité), Locke (légitimité) & Effectivité, consentement des gouvernés \\\hline
\end{tabularx}
\end{center}
\end{exemplebox}

\courssub{TP : Rédaction d'une introduction de dissertation sur « La loi est-elle toujours juste ? »}

\begin{methode}[Gabarit d'introduction type Bac (4 étapes, ~15 lignes)]
\begin{enumerate}
\item \textbf{Accroche (1-2 lignes) :} Citation, fait d'actualité, paradoxe. \emph{Exemple :} « "La loi est la même pour tous, qu'elle protège ou qu'elle opprime", écrit Anatole France. Pourtant, l'histoire montre des lois racistes, esclavagistes, totalitaires. »
\item \textbf{Définition des concepts (2-3 lignes) :} \cle{Loi} = règle générale, abstraite, obligatoire, sanctionnée par la force publique (positivisme). \cle{Justice} = idéal moral d'équité, de proportionnalité, de respect des droits (jusnaturalisme).
\item \textbf{Problématisation (3-4 lignes) :} Montrer la tension. « Si la loi \emph{incarne} la justice en assurant l'égalité formelle et la paix civile (légalité), elle peut aussi \emph{la trahir} lorsqu'elle consacre des privilèges, discrimine ou viole les droits fondamentaux (légitimité). Dès lors, \textbf{la conformité à la loi suffit-elle à garantir la justice, ou la justice exige-t-elle des principes supérieurs à la loi ?} »
\item \textbf{Annonce du plan (2 lignes) :} « Nous montrerons d'abord que la loi est le fondement nécessaire de la justice sociale (I), avant de démontrer qu'elle peut devenir injuste sans contrôle par des normes supérieures (II), pour conclure que la loi n'est juste que si elle se soumet sans cesse à l'exigence d'équité et de droits fondamentaux (III). »
\end{enumerate}
\end{methode}

\begin{exoresolu}[Introduction rédigée — Sujet : « La loi est-elle toujours juste ? »]
« "La loi est l'expression de la volonté générale", affirme Rousseau dans \emph{Du contrat social}. Pourtant, les lois de Nuremberg (1935), le Code de l'indigénat en Afrique coloniale, ou l'apartheid légal en Afrique du Sud étaient des \emph{lois} : générales, obligatoires, sanctionnées. Etaient-elles \emph{justes} ? La \cle{loi} est la règle positive, écrite, imposée par l'État (légalité formelle). La \cle{justice} est l'idéal moral d'équité, de proportionnalité, de respect de la dignité (légitimité matérielle). Si la loi assure l'égalité formelle et la sécurité juridique — conditions \emph{nécessaires} de la justice —, elle n'en est pas la condition \emph{suffisante} : une loi peut être injuste par son contenu (discrimination) ou par son application (inéquité). \textbf{La loi est-elle toujours juste, ou la justice exige-t-elle de juger la loi elle-même à l'aune de principes supérieurs ?} Nous analyserons d'abord pourquoi la loi est le rempart indispensable de la justice (I), puis pourquoi la légalité ne saurait se confondre avec la légitimité (II), pour enfin montrer comment l'État de droit moderne tente d'articuler loi et justice par le contrôle de constitutionnalité et l'équité (III). »
\tcblower
\textbf{Analyse de la copie (grille Bac) :}
\begin{itemize}
    \item Accroche : citation + contre-exemples historiques précis (3/3).
    \item Définitions : loi (positivisme), justice (jusnaturalisme), distinction légalité/légitimité (4/4).
    \item Problématisation : tension réelle, question claire, concepts opposés (5/5).
    \item Plan : dialectique (nécessité / limite / articulation), annoncée sans lourdeur (3/3).
    \item Langue : termes techniques, connecteurs, pas de familier (5/5).
\end{itemize}
\textbf{Note : 20/20.}
\end{exoresolu}

\courssub{Critères d'évaluation au Probatoire et au Baccalauréat technique}

\begin{aretenir}[Grille de correction type (sur 20 pts) — Dissertation / Explication]
\begin{itemize}
    \item \textbf{Exactitude des définitions (4 pts) :} Droit, justice, égalité, équité, légalité, légitimité, justice distributive/corrective définis \emph{avec leurs nuances} (pas de définitions approximatives).
    \item \textbf{Pertinence des exemples (4 pts) :} Au moins 2 exemples \textbf{concrets, datés, camerounais ou universels} (pas d'exemples vagues : « dans la vie quotidienne... »). Ex. : « Loi n°2019/024 du 24 déc. 2019 sur la décentralisation », « Jurisprudence Cour Suprême n°... », « Coutume Bamiléké de règlement des conflits fonciers ».
    \item \textbf{Cohérence de l'argumentation (6 pts) :} Problématique réelle, plan dialectique (thèse/antithèse/synthèse), transitions logiques, pas de « liste de cours ». Chaque partie \emph{répond} à la problématique.
    \item \textbf{Maîtrise des références philosophiques (3 pts) :} Noms d'auteurs (Aristote, Kant, Rawls, Montesquieu, Kelsen, Locke, Rousseau, Ricoeur...) associés à \emph{leur thèse précise} sur le point discuté. Pas de name-dropping vide.
    \item \textbf{Qualité de l'expression (3 pts) :} Français correct, vocabulaire philosophique, connecteurs logiques (donc, toutefois, en revanche, par conséquent, c'est pourquoi), introduction/conclusion soignées.
\end{itemize}
\end{aretenir}

\begin{attention}[Erreur fatale au Bac : réciter le cours sans répondre à la problématique]
Le correcteur ne note pas « ce que vous savez » mais « comment vous \emph{répondez} à la question ». Un plan qui énumère « 1) Définition du droit 2) Fondements 3) Formes de justice » est \textbf{hors sujet} sur un sujet comme « La loi est-elle toujours juste ? » ou « Faut-il toujours obéir aux lois ? ». \textbf{Chaque partie doit être une étape de la réponse.} Vérifiez : si je remplace le titre de la partie par « Argument 1 », « Argument 2 », « Argument 3 », le fil logique tient-il ?
\end{attention}

\begin{exemplebox}[Sujet type Bac : « Faut-il toujours obéir aux lois ? » — Plan détaillé]
\textbf{Problématique :} L'obéissance à la loi est-elle une obligation absolue (ordre, pacte social) ou une obligation conditionnée (conscience, justice, droits de l'homme) ?
\begin{itemize}
    \item \textbf{Thèse (Obligation absolue / quasi-absolue) :} 1) Pacte social (Hobbes, Locke, Rousseau) : on a consenti à la loi pour sortir de l'état de nature ; désobéir = retour au chaos. 2) Positivisme juridique (Kelsen) : la validité de la loi ne dépend pas de sa moralité ; seul le juge constitutionnel peut l'invalider. 3) Utilité sociale : prévisibilité, sécurité, égalité formelle.
    \item \textbf{Antithèse (Droit/Dévoir de désobéissance) :} 1) Jusnaturalisme (Antigone, Thomas d'Aquin, Locke) : loi injuste $\neq$ loi ; « \emph{lex injusta non est lex} ». 2) Désobéissance civile (Thoreau, Gandhi, Martin Luther King) : publique, non-violente, acceptation de la peine, pour changer la loi. 3) Droit de résistance (DDHC 1789, Art. 2 ; Constitution camerounaise Préambule) : l'obéissance cesse si la loi viole les droits naturels.
    \item \textbf{Synthèse (Obéissance critique dans l'État de droit) :} 1) Principe : on obéit aux lois justes \emph{et} aux lois imparfaites mais démocratiques (stabilité). 2) Exception : désobéissance \emph{encadrée} (contrôle constitutionnel, objection de conscience, désobéissance civile) quand la loi porte atteinte au noyau dur des droits fondamentaux. 3) La loi n'est pas un maître absolu, elle est l'instrument d'une justice qui la dépasse.
\end{itemize}
\textbf{Exemples camerounais à mobiliser :} Grèves syndicales (droit de grève vs réquisition), refus d'appliquer une circulaire illégale (fonctionnaire), mouvements citoyens (ex. « Villes mortes » — analyse juridique vs légitimité politique), jurisprudence Cour Suprême sur l'inconstitutionnalité.
\end{exemplebox}

\begin{savaistu}[Le savais-tu ? Les sujets de philosophie au Probatoire et au Baccalauréat technique camerounais]
Depuis 2015, les sujets du Probatoire et du Baccalauréat technique reviennent \emph{régulièrement} sur le triptyque \textbf{Droit -- État -- Liberté} :
\begin{itemize}
    \item « La loi est-elle la condition de la liberté ? » (Session 2018, 2022)
    \item « Faut-il toujours obéir à l'État ? » (2019, 2023)
    \item « La justice sociale est-elle possible ? » (2020)
    \item « Le droit peut-il fonder la morale ? » (2021)
    \item « L'égalité devant la loi suffit-elle à garantir la justice ? » (2024)
\end{itemize}
\emph{Conseil stratégique :} Maîtrisez \textbf{par cœur} 3 exemples camerounais datés (loi, jurisprudence, coutume, fait social) pour \emph{chacune} de ces grandes questions. C'est ce qui fait la différence entre 12/20 et 17/20.
\end{savaistu}

\begin{exercice}[Entraînement final — À faire en autonomie (2h)]
\begin{enumerate}
    \item \textbf{Explication de texte :} Analysez l'extrait suivant (Kant, \emph{Métaphysique des mœurs}, § 49) : « La liberté extérieure qui peut coexister avec la liberté de chacun selon une loi universelle, c'est le droit. » (Introduction, analyse en 3 étapes, synthèse, ouverture).
    \item \textbf{Dissertation :} Sujet : « L'égalité devant la loi suffit-elle à réaliser la justice ? » Rédigez l'introduction complète + le plan détaillé (thèse/antithèse/synthèse avec arguments et exemples camerounais pour chaque sous-partie).
    \item \textbf{Tableau synoptique :} Complétez le tableau de la figure (p. 3) avec 6 lignes originales (différentes de l'exemple donné). Faites valider par votre professeur.
\end{enumerate}
\end{exercice}

\begin{corrige}[Exercice n°2 — Dissertation : « L'égalité devant la loi suffit-elle à réaliser la justice ? »]
\textbf{Introduction type :} « "La loi, dans sa majestueuse égalité, interdit aux riches comme aux pauvres de dormir sous les ponts, de mendier dans les rues et de voler du pain" (Anatole France, \emph{Le Lys rouge}). Cette ironie révèle le fossé possible entre \cle{égalité formelle} (même texte pour tous) et \cle{justice réelle} (équité des situations). L'égalité devant la loi (Art. 6 DDHC 1789, Art. 1er Constitution camerounaise) est un principe cardinal de l'État de droit. La justice, elle, vise la proportionnalité, la correction des inégalités de fait, le respect des besoins et des mérites. \textbf{L'uniformité juridique est-elle une justice, ou masque-t-elle des inégalités structurelles que seule l'équité peut corriger ?} Nous verrons que l'égalité formelle est le socle indispensable (I), qu'elle devient injuste si elle ignore les inégalités réelles (II), pour conclure que la justice exige de passer de l'égalité arithmétique à l'équité proportionnelle (III). »

\textbf{Plan détaillé :}
\begin{itemize}
    \item \textbf{I. L'égalité devant la loi : fondement nécessaire de la justice (Thèse)}
    \begin{itemize}
        \item A. Principe d'universalité et d'abstraction : la loi ne connaît pas les personnes (Kelsen, Art. 1er Const. CM). Ex. : Code pénal s'applique au ministre comme au paysan.
        \item B. Sécurité juridique et prévisibilité : même cause $\to$ même effet. Évite l'arbitraire (Montesquieu, \emph{Esprit des lois}).
        \item C. Condition de la citoyenneté : égalité des droits civiques (vote, éligibilité, fonction publique). Ex. : Loi n°2012/001 du 19 avr. 2012 (Code électoral) : un citoyen = une voix.
    \end{itemize}
    \item \textbf{II. Les limites de l'égalité formelle : quand la loi égale produit de l'injustice (Antithèse)}
    \begin{itemize}
        \item A. Égalité arithmétique $\neq$ équité : même impôt proportionnel (TVA) pèse plus sur le pauvre (Rawls, \emph{Théorie de la justice}). Ex. : Fiscalité camerounaise — TVA 19,25\% régressive.
        \item B. Inégalités de fait non corrigées : même droit à la santé, mais hôpital loin du village ; même droit à l'école, mais frais de scolarité exclusifs. Ex. : Disparités régionales (Extrême-Nord vs Littoral) en taux de scolarisation.
        \item C. Lois neutres en apparence, discriminantes en effet : condition de taille pour concours (femmes exclues), exigence de diplôme non pertinent (jeunes ruraux). Jurisprudence Cour Suprême : discrimination indirecte.
    \end{itemize}
    \item \textbf{III. Vers une justice équitable : corriger l'égalité par la proportionnalité (Synthèse)}
    \begin{itemize}
        \item A. Justice distributive (Aristote) : proportion géométrique selon mérite/besoin. Ex. : Bourses sur critères sociaux + mérite ; quotas régionaux (décret 2011/...).
        \item B. Discrimination positive / action affirmative : traitement \emph{inégal} pour rétablir l'égalité \emph{réelle}. Ex. : Parité genre dans les listes électorales (Loi 2012/001), recrutement personnes handicapées (Loi 2010/002).
        \item C. Équité du juge / contrôle de constitutionnalité : le juge adapte la sanction (circonstances atténuantes), le Conseil constitutionnel censure les lois créant des inégalités injustifiées. Ex. : Décision n°... du Conseil constitutionnel sur la loi de finances.
    \end{itemize}
\end{itemize}
\textbf{Conclusion :} L'égalité devant la loi est la \emph{condition de possibilité} de la justice, non sa \emph{réalisation achevée}. La justice commande de regarder non seulement le texte de la loi, mais ses \emph{effets} sur des sujets inégaux. L'État de droit camerounais, par sa Constitution et ses lois récentes, tente ce passage de l'égalité formelle à l'équité effective — mais le chemin reste long (accès au droit, corruption, effectivité des jugements).
\end{corrige}

\newpage

\coursec{Activité d'intégration}

\coursec{Le Droit et la Justice}

\courssub{Objectifs de l'activité}
Cette activité de type \emph{tribunal simulé} vise à mettre en œuvre les savoirs et savoir-faire de l'unité « Le Droit et la Justice » à travers une situation concrète inspirée du contentieux commercial camerounais. Elle permet aux élèves de :
\begin{itemize}
    \item \textbf{Identifier} et \textbf{distinguer} les \cle{droits subjectifs} (créance, propriété, responsabilité) qui s'opposent dans un litige ;
    \item \textbf{Qualifier} la nature de la norme applicable : \cle{droit positif} (Code civil, OHADA) vs \cle{droit coutumier} (accord verbal, usage du transport) ;
    \item \textbf{Classifier} la décision judiciaire attendue selon les trois \cle{formes de la justice} aristotéliciennes : \cle{commutative}, \cle{distributive}, \cle{rétributive} ;
    \item \textbf{Rédiger} un \cle{verdict motivé} en opérant la distinction critique entre \cle{légalité} (conformité à la règle de droit) et \cle{équité} (justice concrète, adaptation au cas particulier) ;
    \item \textbf{Construire} un \cle{tableau synoptique} de synthèse juridico-philosophique ;
    \item \textbf{Argumenter} à l'oral lors de la restitution : « Le verdict légal est-il forcément juste ? ».
\end{itemize}

\courssub{Situation}
\textbf{Contexte :} Douala, quartier Bépanda, mars 2025. La route nationale 3 (axe Douala-Yaoundé) est dans un état critique suite aux fortes pluies : nids-de-poule, éboulements, ralentissements obligatoires. Les transporteurs routiers subissent des pertes sèches (carburant, temps, entretien véhicules).

\textbf{Les parties :}
\begin{itemize}
    \item \textbf{M. Étienne Ngando (Plaignant) :} Commerçant en quincaillerie à Akwa. Il a confié un chargement de \textbf{500 kg de tubes PVC} (valeur marchande : \textbf{2 500 000 FCFA}) à la société de transport « \textit{Express Bassa} » pour livraison à Bertoua (Est-Cameroun). Facture de transport : \textbf{150 000 FCFA} (payée d'avance). Contrat écrit : bon de commande \textit{Express Bassa} n° 4482, mention « livraison sous 48h, responsabilité du transporteur en cas de perte/avarie ».
    \item \textbf{M. Jean-Baptiste Owona (Défendeur) :} Gérant de « \textit{Express Bassa} ». Le camion (immatriculé LT 458 AX) a eu un accident non responsable (glissement sur boue, sortie de route) au PK 120 (vers Édéa). La marchandise a été \textbf{volée par des riverains} pendant que le chauffeur (M. Pierre) cherchait de l'aide. Le camion est immobilisé (réparations : 800 000 FCFA).
    \item \textbf{Élément coutumier :} M. Owona invoque un \textbf{accord verbal} conclu avec M. Ngando en 2023 : « \textit{Si la route fait défaut, on partage la perte à moitié. C'est la coutume entre nous, commerçants de Bassa.} » Aucun écrit ne constate cet accord. La Chambre de Commerce de Douala confirme l'existence d'un \textbf{usage local} dit « \textit{partage du risque route} » non codifié.
\end{itemize}

\textbf{Données chiffrées clés :}
\begin{center}
\begin{tabularx}{\linewidth}{|l|X|}\hline
\textbf{Poste} & \textbf{Montant (FCFA)} \\ \hline
Valeur marchandise (tubes PVC) & 2 500 000 \\
Fret payé d'avance & 150 000 \\
Réparation camion (charge transporteur) & 800 000 \\
Perte d'exploitation transporteur (15 jours) & 450 000 (estimé) \\
Offre transactionnelle transporteur (amiable) & 1 000 000 (refusée) \\ \hline
\end{tabularx}
\end{center}

\textbf{Procédure :} M. Ngando a saisi le \textbf{Tribunal de Première Instance de Douala-Bonanjo} (droit positif OHADA, Acte uniforme sur le transport de marchandises). M. Owona soulève l'\textbf{exception d'incompétence} au profit de la \textbf{chefferie traditionnelle Bassa} (droit coutumier) pour appliquer l'accord verbal.

\courssub{Consignes}
Les élèves sont répartis en trois collèges : \textbf{Plaignants} (avocats de M. Ngando), \textbf{Défenseurs} (avocats de M. Owona), \textbf{Juges} (3 à 5 élèves). Chaque collège travaille sur les 5 étapes ci-dessous, puis les Juges rendent un verdict unique après délibération.

\begin{enumerate}
    \item \textbf{Identification des droits subjectifs :} Listez les \cle{droits subjectifs} (droits de créance, droits réels, droits de la personnalité) dont se prévalent chacune des parties. Distinguez le \cle{droit subjectif} (pouvoir juridique reconnu à un sujet) du \cle{droit objectif} (ensemble des règles).
    \item \textbf{Qualification de la norme applicable :} Le litige relève-t-il du \cle{droit positif} (OHADA, Code civil camerounais) ou du \cle{droit coutumier} ? Justifiez en citant l'article 1\textsuperscript{er} du Code civil (hiérarchie des normes) et la condition de \cle{non-contrariété à l'ordre public} de la coutume.
    \item \textbf{Classification selon les formes de justice :} La décision à rendre (condamnation à payer des dommages-intérêts, partage de la perte, relaxe) relève-t-elle de la \cle{justice commutative} (égalité arithmétique contractuelle), de la \cle{justice distributive} (égalité géométrique, proportion au mérite/besoin) ou de la \cle{justice rétributive} (sanction proportionnée à la faute) ? Peut-elle en relever plusieurs simultanément ?
    \item \textbf{Rédaction du verdict motivé :} Rédigez le dispositif du jugement (300 à 500 mots). Distinguez explicitement : \textit{« Attendu que... »} (motifs de \cle{légalité} : application stricte de la loi) et \textit{« Considérant que... »} (motifs d'\cle{équité} : circonstances atténuantes, usage local, bonne foi).
    \item \textbf{Tableau synoptique récapitulatif :} Construisez un tableau à 4 colonnes : \textit{Notion mobilisée} / \textit{Définition} / \textit{Application au cas Ngando vs Owona} / \textit{Portée philosophique (légalité vs justice)}.
\end{enumerate}

\courssub{Ressources à mobiliser}
\begin{definition}[Droit objectif et Droit subjectif]
Le \cle{droit objectif} est l'ensemble des règles juridiques obligatoires régissant la société (ex. : Acte uniforme OHADA transport). Le \cle{droit subjectif} est la prérogative, le pouvoir reconnu à un sujet de droit par la règle objective (ex. : droit de propriété de Ngando sur les tubes, droit à réparation du préjudice).
\end{definition}

\begin{definition}[Justice commutative, distributive, rétributive (Aristote, \emph{Éthique à Nicomaque}, V)]
\begin{itemize}
    \item \textbf{Commutative :} Régit les échanges entre particuliers (contrats, délits). Vise l'\textbf{égalité arithmétique} : restituer l'équivalent exact (valeur pour valeur). \textit{Ex. : payer le prix convenu, réparer le dommage causé.}
    \item \textbf{Distributive :} Régit la répartition des biens/honneurs/charges dans la cité. Vise l'\textbf{égalité géométrique} (proportionnelle) : à chacun selon son mérite, son besoin, sa contribution. \textit{Ex. : impôt progressif, allocation des marchés publics.}
    \item \textbf{Rétributive (ou corrective) :} Sanctionne la violation de la loi. Vise la \textbf{proportionnalité de la peine} à la gravité de la faute. \textit{Ex. : peine de prison, amende.}
\end{itemize}
\end{definition}

\begin{propriete}[Hiérarchie des normes en droit camerounais (Art. 1\textsuperscript{er} Code civil + Constitution)]
La \cle{loi} (droit positif écrit : Constitution, traités OHADA, lois parlementaires) prime sur la \cle{coutume}. La coutume ne s'applique qu'\textbf{à défaut de loi} et \textbf{si elle n'est pas contraire à l'ordre public} ni aux bonnes mœurs. Un accord verbal ne peut déroger à une obligation contractuelle écrite (principe de la force obligatoire des contrats, art. 1101 Code civil / Art. 4 Acte uniforme OHADA transport).
\end{propriete}

\begin{aretenir}[Légalité vs Équité]
La \cle{légalité} est la conformité stricte à la règle de droit positive (le juge \textit{bouche de la loi}). L'\cle{équité} (au sens d'\textit{aequitas}) est le pouvoir du juge d'adapter la règle au cas concret pour éviter l'injustice rigide (\textit{summum jus, summa injuria}). En droit OHADA, le juge peut statuer \textit{ex aequo et bono} seulement si les parties le lui demandent expressément (Art. 31 Règlement de procédure CCJA), sinon il juge en droit.
\end{aretenir}

\courssub{Démarche et corrigé commenté}
\begin{exoresolu}[Corrigé guidé du Tribunal de la classe]
\textbf{Énoncé :} Se reporter à la situation, aux consignes et aux ressources mobilisées ci-dessus.
\tcblower
\textbf{Solution détaillée :}

\textbf{Étape 1 : Identification des droits subjectifs}
\begin{itemize}
    \item \textbf{M. Ngando (Plaignant) :}
    \begin{itemize}
        \item \cle{Droit de propriété} sur les tubes PVC (art. 544 Code civil : droit absolu, exclusif). La perte du bien = atteinte à ce droit.
        \item \cle{Droit de créance} (droit personnel) né du contrat de transport : exiger l'exécution (livraison) ou, à défaut, la \cle{réparation intégrale du préjudice} (dommages-intérêts = valeur marchandise + frais accessoires). Fondement : responsabilité contractuelle du transporteur (Art. 22 Acte uniforme OHADA transport : « Le transporteur est responsable de la perte... »).
        \item \cle{Droit à l'exécution de bonne foi} du contrat (Art. 1104 Code civil).
    \end{itemize}
    \item \textbf{M. Owona (Défendeur) :}
    \begin{itemize}
        \item \cle{Droit de propriété} sur son camion (atteint par l'accident).
        \item \cle{Droit subjectif à ne pas répondre d'un cas de force majeure} (Art. 1218 Code civil / Art. 23 Acte uniforme : exonération si \textit{fait étranger, imprévisible, irrésistible}). Il invoque l'état de la route (fait de la nature/autorité publique) et le vol par tiers (fait d'un tiers).
        \item \cle{Droit à l'indemnisation} de son propre préjudice (camion, perte d'exploitation) si la responsabilité de l'État (entretien route) ou des voleurs est établie (action récursoire).
    \end{itemize}
\end{itemize}
\textbf{Distinction clé :} Le droit objectif (la loi OHADA) \textbf{crée} le droit subjectif de créance de Ngando. L'accord verbal invoqué par Owona prétend créer un \textbf{droit subjectif nouveau} (droit à partage de perte) dérivé de la coutume.

\textbf{Étape 2 : Qualification de la norme applicable (Droit positif vs Coutume)}
\begin{itemize}
    \item \textbf{Principe :} L'Acte uniforme OHADA sur le transport de marchandises (ratifié par le Cameroun) est une \cle{loi} (droit positif). Il s'applique \textbf{impérativement} aux contrats de transport routier de marchandises.
    \item \textbf{Application :} L'art. 22 pose la responsabilité de plein droit du transporteur pour perte/avarie. L'art. 23 limite les exonérations à la \cle{force majeure} (événement extérieur, imprévisible, irrésistible). \textbf{L'état délabré de la route au Cameroun est-il imprévisible ?} Jurisprudence constante (CCJA, arrêts 2018, 2021) : \textbf{Non}. L'état des routes est un risque normal, prévisible de l'activité de transporteur. Donc \textbf{pas de force majeure}.
    \item \textbf{L'accord verbal / Coutume :} L'art. 1\textsuperscript{er} Code civil : « La coutume ne s'applique qu'à défaut de dispositions législatives... pourvu qu'elle ne soit pas contraire à l'ordre public. » Ici, la loi (OHADA) existe \textbf{et} dispose de l'exonération (force majeure). La coutume « partage du risque route » ne peut s'appliquer \textbf{à défaut de loi}. De plus, elle contredit l'ordre public économique (protection du chargeur, sécurité juridique du transport). L'accord verbal ne peut déroger au contrat écrit (principe \textit{sollemnité} relative pour le transport, mais surtout force obligatoire du contrat écrit).
    \item \textbf{Conclusion :} Le tribunal est \textbf{compétent} (droit positif). La coutume/accord verbal est \textbf{inopposable} à M. Ngando pour réduire la responsabilité du transporteur.
\end{itemize}

\textbf{Étape 3 : Classification selon les formes de justice}
La décision à rendre (condamnation du transporteur à payer 2 500 000 FCFA + intérêts + dépens) relève principalement de :
\begin{itemize}
    \item \textbf{Justice commutative (principale) :} Il s'agit d'un \textbf{contrat} (échange volontaire : fret contre obligation de livrer). La perte de la marchandise rompt l'égalité arithmétique du contrat (Ngando a payé le fret, n'a rien reçu). La condamnation à \textbf{réparer intégralement le préjudice} (valeur exacte de la chose perdue) restaure l'égalité arithmétique : \textit{valeur reçue = valeur due}. C'est la \cle{justice corrective/commutative} par excellence (Aristote, \textit{Éthique à Nicomaque}, V, 4).
    \item \textbf{Justice rétributive (accessoire) :} Si le juge condamne le transporteur pour \textbf{faute de gestion} (non adaptation du véhicule à la saison, négligence du chauffeur), les dommages-intérêts majorés ou une astreinte relèvent de la sanction (proportionnée à la faute). Mais en responsabilité contractuelle de droit commun, la faute n'est pas requise (responsabilité de plein droit), ce qui renforce le caractère \textbf{commutatif} (réparation du déséquilibre contractuel) et non rétributif (punition).
    \item \textbf{Justice distributive (absente) :} Il ne s'agit pas de répartir un bien commun entre citoyens selon un critère de mérite/besoin (impôts, emplois publics), mais de corriger une inégalité née d'un lien contractuel bilatéral.
\end{itemize}
\textbf{Nuance :} Si les Juges décident un \textbf{partage de la perte} (ex. 50/50) au nom de l'équité/coutume, ils font glisser la décision vers une \textbf{justice distributive} (répartition proportionnelle du fardeau « route » entre deux acteurs économiques) ou une \textbf{transaction} (accord des parties), mais sortent du cadre strict de la justice commutative légale.

\textbf{Étape 4 : Rédaction du verdict motivé (Modèle de rédaction attendue)}

\begin{center}
\textbf{TRIBUNAL DE PREMIÈRE INSTANCE DE DOUALA-BONANJO}\\
\textbf{AUDIENCE PUBLIQUE DU 15 MAI 2025}\\
\textbf{JUGEMENT N\textsuperscript{o} 2025/0458}
\end{center}

\textbf{ENTRE :}
\textbf{M. Étienne NGANDO}, demeurant Douala-Akwa, \textbf{Plaignant} / \textbf{Demandeur}.
\textbf{ET :}
\textbf{Société EXPRESS BASSA}, prise en la personne de son gérant M. Jean-Baptiste OWONA, sise Douala-Bépanda, \textbf{Défenderesse} / \textbf{Défenderesse}.

\textbf{LE TRIBUNAL,}
\textbf{Vu} la requête introductive d'instance du 10 avril 2025 ;
\textbf{Vu} les conclusions de la défenderesse du 28 avril 2025 soulevant l'exception d'incompétence et le fond ;
\textbf{Vu} l'Acte uniforme OHADA relatif au transport de marchandises (juillet 2003, révisé 2010) ;
\textbf{Vu} le Code civil camerounais (notamment articles 1\textsuperscript{er}, 1101, 1104, 1218, 1240) ;
\textbf{Vu} la Constitution du 18 janvier 1996 (prééminence du droit écrit) ;

\textbf{ATTENDU QUE} (Motifs de \textbf{LÉGALITÉ}) :
\begin{enumerate}
    \item \textbf{Sur la compétence :} L'Acte uniforme OHADA est une loi applicable au Cameroun. L'exception d'incompétence au profit de la chefferie traditionnelle est \textbf{rejetée}. La coutume invoquée (« partage du risque route ») ne s'applique qu'à défaut de loi (Art. 1\textsuperscript{er} Code civil). Or la loi (OHADA) régit exhaustivement la responsabilité du transporteur. De plus, cet usage contredit l'ordre public économique en vidant de sa substance l'obligation de sécurité du transporteur. L'accord verbal, non écrit, ne peut prévaloir sur le bon de commande signé (Art. 1101 Code civil : force obligatoire des contrats).
    \item \textbf{Sur le fond (Responsabilité) :} Aux termes de l'Art. 22 de l'Acte uniforme, le transporteur est responsable de la perte des marchandises. L'exonération (Art. 23) exige la preuve d'un cas de force majeure (extérieur, imprévisible, irrésistible). L'état de la route Nationale 3, bien que dégradé, est un risque \textbf{prévisible} et \textbf{permanent} inhérent à l'activité de transport au Cameroun (jurisprudence CCJA constante). Le vol par des tiers, survenu pendant l'immobilisation du véhicule due à un accident non force majeure, ne rompt pas le lien de causalité. La société \textit{Express Bassa} est donc \textbf{pleinement responsable} de la perte.
    \item \textbf{Sur le quantum :} Le préjudice certain et direct est la valeur de la marchandise au lieu et moment de la livraison convenue (Bertoua) : \textbf{2 500 000 FCFA}. S'y ajoutent le remboursement du fret payé pour une prestation non exécutée : \textbf{150 000 FCFA}. Les intérêts au taux légal courent du jour de la mise en demeure (requête). Les dépens sont mis à la charge de la défenderesse.
\end{enumerate}

\textbf{CONSIDÉRANT QUE} (Motifs d'\textbf{ÉQUITÉ} - \textit{Aequitas}) :
\begin{enumerate}
    \item Le Tribunal note la gravité exceptionnelle de l'état de la route (pluies diluviennes, absence de signalisation, défaillance de l'entretien public) qui a contribué à l'accident. Si la loi ne permet pas d'exonérer le transporteur (pas de force majeure), l'équité commande de ne pas alourdir la condamnation par des dommages-intérêts punitifs (intérêts moratoires majorés, astreintes) demandés par le demandeur.
    \item Le Tribunal relève la bonne foi apparente de M. Owona (entreprise locale, difficultés réelles du secteur) et l'existence d'un usage professionnel connu, bien que juridiquement inopérant. Il invite les parties à saisir la \textbf{Commission de Conciliation de la Chambre de Commerce} pour tenter un accord transactionnel sur les modalités de paiement (échelonnement), sans préjudice du principal.
\end{enumerate}

\textbf{PAR CES MOTIFS,}
\textbf{REJETTE} l'exception d'incompétence ;
\textbf{CONDAMNE} la société \textit{Express Bassa} à payer à M. Étienne Ngando la somme de \textbf{2 650 000 FCFA} (deux millions six cent cinquante mille francs CFA) en principal, avec intérêts au taux légal dès signification, et aux dépens ;
\textbf{DIT} que le présent jugement est exécutoire par provision nonobstant appel ;
\textbf{ORDONNE} la mention du présent jugement en marge du registre du commerce de la défenderesse.

\textbf{Étape 5 : Tableau synoptique récapitulatif}
\begin{center}
\begin{tabularx}{\linewidth}{|l|X|X|X|}\hline
\rowcolor{popblueL}
\textbf{Notion mobilisée} & \textbf{Définition synthétique} & \textbf{Application au cas Ngando vs Owona} & \textbf{Portée : Légalité vs Justice} \\ \hline
\textbf{Droit objectif (OHADA Transport)} & Ensemble des règles écrites, générales, obligatoires. & Source unique de la responsabilité du transporteur (Art. 22). Exclut la coutume. & \textbf{Légalité pure} : Le juge applique la règle, sans pouvoir la modifier. \\ \hline
\textbf{Droit subjectif (Créance de Ngando)} & Pouvoir d'exiger un comportement (payer, livrer, réparer). & Ngando exige réparation intégrale (2,5 M + fret). Fondé sur le contrat + loi. & \textbf{Légalité} : Le droit subjectif naît du droit objectif. \\ \hline
\textbf{Coutume / Accord verbal} & Règle non écrite, constante, crue obligatoire par un groupe. & Invoquée par Owona pour partager la perte (50/50). & \textbf{Inopérante en droit (légalité)} car contraire à la loi et à l'ordre public. \textbf{Pertinente en équité} (contexte réel). \\ \hline
\textbf{Justice Commutative} & Égalité arithmétique dans les échanges bilatéraux (contrats, délits). & Condamnation à valeur exacte perdue (2,5 M) = restauration de l'égalité contractuelle. & \textbf{Justice légale} : La loi OHADA réalise la justice commutative par la réparation intégrale. \\ \hline
\textbf{Justice Distributive} & Égalité géométrique (proportionnelle) dans la répartition sociale. & Absente du verdict légal. Aurait été mobilisée si partage 50/50 (répartition du risque « route »). & \textbf{Justice équitable potentielle} : Partager le fardeau commun (route) selon la capacité de chacun. \\ \hline
\textbf{Justice Rétributive} & Sanction proportionnée à la faute (peine). & Non mobilisée (responsabilité sans faute). Serait pertinente si faute prouvée (survitesse, camion défectueux). & \textbf{Légalité pénale/administrative} distincte de la réparation civile. \\ \hline
\textbf{Équité (Aequitas)} & Adaptation de la règle au cas concret pour éviter l'injustice rigide. & Motifs « Considérant que » : pas de majoration punitive, invitation à transaction amiable pour paiement. & \textbf{Justice concrète} : Tempère la légalité stricte par l'humanité du jugement. \\ \hline
\end{tabularx}
\end{center}

\textbf{Restitution orale \& Débat : « Le verdict légal est-il forcément juste ? »}
\begin{itemize}
    \item \textbf{Thèse 1 (Oui - Légalisme/Kelsen) :} Le verdict est juste \textit{parce que} légal. La loi OHADA incarne la volonté démocratique (État de droit), protège le faible (chargeur) contre le fort (transporteur organisé). La sécurité juridique exige l'application stricte. \textit{« Dura lex, sed lex »}.
    \item \textbf{Thèse 2 (Non - Jusnaturalisme/Équité/Aristote) :} La loi est imparfaite. Elle ignore la réalité camerounaise (routes défaillantes = responsabilité partagée État/transporteur). Condamner seul le transporteur à 2,65 M FCFA (risque de faillite) pour un accident dû à la route est une \textit{injustice matérielle} (\textit{summum jus, summa injuria}). L'équité (partage, transaction) est plus juste que la loi rigide.
    \item \textbf{Synthèse (Ricœur / Paul Ricoeur, \textit{Le Juste}) :} La justice vise l'\textbf{équivalence} (commutative) mais passe par la \textbf{loi} (universelle). Le juge \textbf{doit} juger en droit (légalité) mais \textbf{peut/doit} user de l'équité dans l'interprétation (quantum, délais, exécution) pour approcher le Juste. Le verdict \textbf{est légal} ; il \textbf{tend vers la justice} par l'équité procédurale (pas de ruine, invitation transaction), mais \textbf{n'épuise pas l'idée absolue de justice} (qui exigerait que l'État paye sa part de la route).
\end{itemize}
\end{exoresolu}

\courssub{Bilan de l'activité}
Cette activité a permis de mettre en évidence, par la pratique simulée, les articulations centrales du programme :
\begin{enumerate}
    \item \textbf{La distinction opératoire Droit objectif / Droit subjectif :} Les élèves ont vu que le droit objectif (OHADA) \textbf{genère} les droits subjectifs (créance de Ngando) et que la prétention coutumière d'Owona ne créait pas de droit subjectif opposable face à la loi.
    \item \textbf{La hiérarchie des normes comme garde-fou de la légalité :} L'application de l'Art. 1\textsuperscript{er} Code civil a montré que la \cle{coutume} est une source \textbf{subsidiaire} et \textbf{conditionnée} (non contrariété à l'ordre public), ce qui ancre le droit positif comme fondement premier de la décision judiciaire.
    \item \textbf{La pertinence de la typologie aristotélicienne pour qualifier le jugement :} Le contentieux contractuel de transport est l'archétype de la \cle{justice commutative} (restauration de l'égalité arithmétique rompue). La confusion fréquente avec la justice distributive (partage) ou rétributive (punition) a été levée par l'analyse du fondement juridique (responsabilité sans faute = commutative, pas rétributive).
    \item \textbf{La tension constitutive Légalité / Équité :} Le verdict type « Attendu que / Considérant que » incarne la dualité du métier de juge : dire le droit (légalité formelle) tout en cherchant le juste (équité concrète). Les élèves ont expérimenté que \textbf{la légalité est nécessaire mais non suffisante à la justice} (risque de \textit{summum jus, summa injuria}).
    \item \textbf{L'ancrage camerounais comme révélateur philosophique :} Le contexte (routes, OHADA, coutume Bassa, CCJA) a montré que les concepts philosophiques (droit, justice, équité, force majeure, ordre public) ne sont pas des abstractions vides mais des \textbf{outils de résolution de conflits réels} dans un État de droit pluraliste.
\end{enumerate}
\textbf{Compétence évaluée :} Capacité à \textbf{mobiliser des concepts philosophiques} (droit subjectif, justice commutative, équité) pour \textbf{analyser une situation juridique concrète}, \textbf{construire un raisonnement argumenté} (verdict) et \textbf{problématiser le rapport loi/justice} à l'oral et par écrit. Cette activité constitue un entraînement direct aux épreuves de \textbf{dissertation} (sujet type : « La loi assure-t-elle la justice ? ») et de \textbf{commentaire de texte} (extrait d'Aristote, Ricœur, Kelsen, ou décision de justice).

\newpage

\coursec{Méthodes et exercices résolus}

\courssub{Méthodes et exercices résolus}

\begin{methode}[Définir rigoureusement les concepts de Droit et de Justice]
    \textbf{Objectif :} Produire une définition philosophique précise, distincte du sens commun ou juridique pur.
    \begin{enumerate}
        \item \textbf{Identifier le genre prochain :} Le Droit est un \cle{système de normes} ; la Justice est une \cle{vertu} ou un \cle{principe d'ordre}.
        \item \textbf{Préciser la différence spécifique :}
            \begin{itemize}
                \item \textbf{Droit :} Ensemble de règles obligatoires (juridiques) sanctionnées par la force publique (État), régissant les rapports sociaux. Distinction \cle{Droit naturel} (universel, antérieur à l'État) / \cle{Droit positif} (écrit, promulgué par l'État).
                \item \textbf{Justice :} Disposition constante à rendre à chacun ce qui lui est dû (\textit{suum cuique tribuere}, Ulpien). Distinction \cle{Justice légale} (conformité à la loi) / \cle{Justice équitable} (conformité à l'équité, correction de la loi).
            \end{itemize}
        \item \textbf{Articuler le lien :} Le Droit est l'instrument ; la Justice est la fin/valeur. Une loi peut être juridique (valide) mais injuste (injuste en équité).
        \item \textbf{Formatter :} « Concept = Genre + Différence spécifique + Exemple/Contrastes ».
    \end{enumerate}
\end{methode}

\begin{exoresolu}[Définition et distinction : Droit positif vs Justice naturelle]
\textbf{Énoncé :} « Une loi votée à l'unanimité par le Parlement peut-elle être injuste ? » Définissez les termes, opposez Droit positif et Justice naturelle, et répondez en justifiant.

\tcblower
\textbf{Solution détaillée :}
\begin{enumerate}
    \item \textbf{Analyse des concepts :}
    \begin{itemize}
        \item \textbf{Droit positif :} Ensemble des règles juridiques écrites, promulguées par l'autorité légitime (Parlement), applicables dans un État donné à un moment donné (ex: Code pénal camerounais). Caractères : \cle{positivity} (écrit), \cle{coercitivité} (sanction étatique), \cle{relativité} (variable selon temps/lieu).
        \item \textbf{Justice naturelle / Justice idéale :} Principe moral universel selon lequel il faut « rendre à chacun son dû ». Elle repose sur la \cle{dignité humaine}, les \cle{droits de l'homme}, l'\cle{équité}. Elle est \cle{universelle} et \cle{intemporelle}.
    \end{itemize}
    \item \textbf{Problématisation :} La validité formelle (procédure parlementaire, unanimité) garantit-elle la conformité à la justice matérielle ?
    \item \textbf{Argumentation :}
    \begin{itemize}
        \item \textbf{Thèse (Légalisme / Positivisme - Kelsen, Hart) :} Le Droit est ce qui est posé. Si la procédure est respectée, la loi est valide. « Justice = Légalité ». L'unanimité renforce la légitimité démocratique.
        \item \textbf{Antithèse (Jusnaturalisme - Aristote, Thomas d'Aquin, DDHC 1789, Dworkin) :} \cle{Lex injusta non est lex} (Une loi injuste n'est pas une loi). L'histoire montre des lois unanimement votées mais criminelles (lois raciales, lois d'apartheid, lois totalitaires). La justice positive (loi) $\neq$ Justice naturelle (droit moral).
        \item \textbf{Synthèse :} La distinction \cle{légalité / légitimité} (Weber, Rawls). Une loi est \emph{légale} (forme) mais peut manquer de \emph{légitimité} (fond moral). Le contrôle de constitutionnalité (Conseil Constitutionnel) et le droit international des droits de l'homme sont des garde-fous institutionnels.
    \end{itemize}
    \item \textbf{Conclusion :} \textbf{Oui, une loi votée à l'unanimité peut être injuste.} La validité juridique (Droit positif) ne coïncide pas nécessairement avec la validité morale (Justice naturelle). L'unanimité politique ne saurait valider la violation des droits fondamentaux (ex: lois nazies, apartheid).
\end{enumerate}
\textbf{Phrase de conclusion type :} « La distinction entre légalité et légitimité impose de juger le Droit positif à l'aune de la Justice naturelle, critère suprême de la validité morale des normes. »
\end{exoresolu}

\begin{methode}[Analyser les fondements du Droit : Nature vs Convention, Force vs Consentement]
    \textbf{Objectif :} Expliquer \emph{pourquoi} on obéit au Droit et \emph{d'où} il tire son autorité.
    \begin{enumerate}
        \item \textbf{Lister les théories classiques :}
            \begin{itemize}
                \item \textbf{Fondement naturel (Jusnaturalisme) :} Droit issu de la nature humaine/raison/Dieu (Aristote, Stoïciens, Thomas d'Aquin, Locke, DDHC 1789). $\rightarrow$ Droit antérieur à l'État, universel.
                \item \textbf{Fondement conventionnel / Contractuel :} Droit né du \cle{contrat social} (Hobbes, Locke, Rousseau). Consentement des gouvernés = source de légitimité. $\rightarrow$ Souveraineté populaire.
                \item \textbf{Fondement historique / sociologique :} Droit comme produit de l'histoire, des mœurs (\cle{coutume}), de l'esprit du peuple (Savigny, École historique). $\rightarrow$ Droit vivant, non décrété.
                \item \textbf{Fondement formel / Positivisme (Kelsen) :} Fondement = \cle{Grundnorm} (norme fondamentale hypothétique). Validité purement formelle, sans référence morale.
                \item \textbf{Fondement réaliste / Force (Thrasymaque, Marx, Weber) :} « La force prime le droit » / Droit = volonté de la classe dominante / Monopole de la violence physique légitime (État).
            \end{itemize}
        \item \textbf{Structurer la réponse :} Thèse (ex: Nature) $\rightarrow$ Antithèse (ex: Convention/Force) $\rightarrow$ Synthèse (Articulation : La force assure l'effectivité, la raison/morale assure la légitimité, l'histoire assure l'ancrage).
        \item \textbf{Vocabulaire clé :} \cle{Légitimité}, \cle{Légalité}, \cle{Souveraineté}, \cle{Contractualisme}, \cle{Positivisme}, \cle{Jusnaturalisme}.
    \end{enumerate}
\end{methode}

\begin{exoresolu}[Fondements : Le contrat social comme source de l'obligation politique]
\textbf{Énoncé :} « Obéit-on à la loi par peur du gendarme ou par devoir moral ? » Analysez les fondements de l'obligation juridique en distinguant contrainte et obligation.

\tcblower
\textbf{Solution détaillée :}
\begin{enumerate}
    \item \textbf{Définitions :}
    \begin{itemize}
        \item \textbf{Contrainte :} Soumission à une force extérieure (physique, sanction). Hétéronomie. Motif : \cle{peur} (Thrasymaque, Hobbes \textit{Léviathan} : pacte de soumission par peur de la mort violente).
        \item \textbf{Obligation :} Soumission à une loi qu'on se donne soi-même (autonomie) ou qu'on reconnaît comme juste. Motif : \cle{devoir moral}, \cle{respect de la loi} (Kant, Rousseau).
    \end{itemize}
    \item \textbf{Analyse des fondements :}
    \begin{itemize}
        \item \textbf{Thèse 1 : La Force (Réalisme politique).} Hobbes : État de nature = guerre de tous contre tous. On obéit par \cle{peur} du pouvoir souverain (Gendarme/État). Le Droit = ordre du plus fort. $\rightarrow$ Obéissance = prudence, non moralité.
        \item \textbf{Thèse 2 : Le Consentement (Contractualisme).} Locke / Rousseau : On obéit car on a \cle{consenti} au pacte social. La loi est l'expression de la \cle{Volonté générale} (Rousseau). Obéir à la loi, c'est obéir à soi-même (\cle{autonomie}). $\rightarrow$ Fondement moral : la parole donnée, le contrat.
        \item \textbf{Thèse 3 : La Raison morale (Kant).} Impératif catégorique : « Agis uniquement d'après une maxime telle que tu puisses vouloir qu'elle devienne une loi universelle. » Le fondement est la \cle{raison pratique pure}. La contrainte étatique n'est que le moyen de garantir la coexistence des libertés (Droit = condition de la liberté).
    \end{itemize}
    \item \textbf{Synthèse :} La \cle{contrainte} (force publique, police, tribunaux) est la \textbf{condition d'effectivité} du Droit (sans elle, le droit reste vœu pieux). L'\cle{obligation morale} (consentement, raison, justice) est la \textbf{condition de légitimité}. Un droit sans force est impuissant ; une force sans droit est tyrannie.
    \item \textbf{Exemple camerounais :} Le respect du Code de la route. Je m'arrête au feu rouge : 1) Peur de l'amende/contrainte (gendarme). 2) Devoir moral / contrat social : je veux que les autres s'arrêtent pour ma sécurité (réciprocité, volonté générale).
\end{enumerate}
\textbf{Phrase de conclusion type :} « Le Droit repose sur la double assise de la force qui l'effectue et de la raison (ou du consentement) qui le légitime ; réduire l'obéissance à la peur, c'est confondre assujettissement et obligation. »
\end{exoresolu}

\begin{methode}[Articuler Justice et Égalité : Proportionnalité arithmétique vs géométrique]
    \textbf{Objectif :} Maîtriser la distinction aristotélicienne fondamentale pour traiter les sujets sur l'égalité, les discriminations positives, la justice sociale.
    \begin{enumerate}
        \item \textbf{Rappeler la définition aristotélicienne (\textit{Éthique à Nicomaque}, V) :} La Justice = \cle{Égalité proportionnelle}. « Le juste est proportionnel. »
        \item \textbf{Distinguuer les deux figures :}
            \begin{itemize}
                \item \textbf{Justice commutative (Arithmétique) :} Échanges entre particuliers (contrats, dommages). \textbf{Égalité stricte :} $A = B$. « À travail égal, salaire égal ». Réparation : restitution exacte (\textit{restitutio in integrum}). Aveugle aux mérites/besoins.
                \item \textbf{Justice distributive (Géométrique) :} Distribution des honneurs, richesses, charges par la Cité. \textbf{Égalité proportionnelle :} $\frac{A}{B} = \frac{Mérite_A}{Mérite_B}$. « À chacun selon son mérite / son besoin / son rang ». Critère de répartition = \cle{mérite} (Aristote), \cle{besoin} (Marx/Socialisme), \cle{capacité} (Rawls).
            \end{itemize}
        \item \textbf{Appliquer aux débats actuels :}
            \begin{itemize}
                \item Égalité formelle (droit identique pour tous) vs Égalité réelle / équité (moyens différenciés pour résultats égaux).
                \item Discriminations positives = Justice distributive géométrique (corriger inégalités structurelles).
            \end{itemize}
        \item \textbf{Piège à éviter :} Confondre « Inégalité » (différence de traitement justifiée par proportionnalité) et « Injustice » (différence non justifiée).
    \end{enumerate}
\end{methode}

\begin{exoresolu}[Justice distributive et discrimination positive au Cameroun]
\textbf{Énoncé :} « Réserver 30\% des places en Grandes Écoles aux élèves des zones d'éducation prioritaires (ZEP) est-il une injustice envers les autres candidats ? » Mobilisez la distinction justice commutative / distributive.

\tcblower
\textbf{Solution détaillée :}
\begin{enumerate}
    \item \textbf{Qualification du litige :} Il s'agit d'une \textbf{distribution de biens rares} (places en Grandes Écoles) par la puissance publique. $\rightarrow$ Relève de la \cle{Justice distributive} (géométrique), non commutative.
    \item \textbf{Analyse commutative (argument de l'injustice apparente) :} Si on raisonne en \cle{justice commutative} (échange marchand, concours anonyme) : même épreuve $\rightarrow$ même note $\rightarrow$ même droit. Réserver des places brise l'égalité arithmétique ($Note_A = Note_B \Rightarrow Place_A = Place_B$). Apparence d'injustice / « passe-droit ».
    \item \textbf{Analyse distributive (justification de l'équité) :} La Justice distributive exige une \cle{proportionnalité} entre la part accordée et le \cle{critère de mérite/besoin}.
    \begin{itemize}
        \item \textbf{Critère mérite pur (Aristote) :} Notes brutes. Favore les favorisés (capital culturel, écoles privées). Reproduit les inégalités sociales.
        \item \textbf{Critère mérite corrigé / potentiel (Rawls, Justice sociale) :} Le mérite réel = performance \emph{au regard des moyens}. Un 14/20 en ZEP $\approx$ 16/20 en lycée d'élite. La mesure vise l'\cle{égalité des chances} (équité), pas l'égalité formelle.
        \item \textbf{Finalité :} Corriger une \cle{injustice structurelle} (inégalité de départ) pour tendre vers une justice absolue.
    \end{itemize}
    \item \textbf{Synthèse :} Ce n'est pas une \cle{injustice} (violation de l'égalité arithmétique), c'est une \cle{justice correctrice} (application de l'égalité géométrique). Elle traite \emph{inégalement} des situations \emph{inégales} pour rétablir l'égalité réelle des chances.
    \item \textbf{Nuance :} Condition de validité : critère transparent, temporaire, proportionné (pas de quota rigide excluant le mérite total), révisable.
\end{enumerate}
\textbf{Phrase de conclusion type :} « La discrimination positive n'est pas une violation de la justice, mais l'application rigoureuse de la justice distributive géométrique : rendre à chacun selon ses conditions réelles de concurrence, pour que l'égalité formelle ne masque pas l'inégalité réelle. »
\end{exoresolu}

\begin{methode}[Classifier les formes de la Justice : Commutative, Distributive, Corrective (Pénale), Sociale]
    \textbf{Objectif :} Identifier le type de justice en jeu dans un énoncé pour mobiliser les bons arguments.
    \begin{enumerate}
        \item \textbf{Justice Commutative (Échanges privés) :} Relation \textbf{A $\leftrightarrow$ B}. Égalité arithmétique. Contrats, ventes, responsabilité civile (réparation). Mots-clés : \cle{Échange}, \cle{Contrat}, \cle{Restitution}, \cle{Égalité stricte}.
        \item \textbf{Justice Distributive (Allocation publique) :} Relation \textbf{État/Collectivité $\rightarrow$ Individus}. Proportionnalité géométrique. Salaires fonction publique, bourses, impôts, logements sociaux. Mots-clés : \cle{Partage}, \cle{Mérite}, \cle{Besoin}, \cle{Contribution}, \cle{Proportionnalité}.
        \item \textbf{Justice Corrective / Pénale (Sanction) :} Relation \textbf{Société $\rightarrow$ Coupable}. But : \cle{Rétablir l'ordre} rompu par le crime/délit.
            \begin{itemize}
                \item \textbf{Rétributive (Kant, Droit classique) :} Peine = \cle{Expiation}. « Œil pour œil » (proportionnalité peine/gravité). Regard vers le passé.
                \item \textbf{Utilitaire / Préventive (Beccaria, Bentham, Droit moderne) :} Peine = \cle{Prévention} (spéciale : récidive ; générale : exemple) + \cle{Réinsertion}. Regard vers l'avenir.
            \end{itemize}
        \item \textbf{Justice Sociale (Synthèse moderne - Rawls) :} Structure les institutions de base. Deux principes : 1) Libertés égales maximales. 2) Inégalités sociales acceptées seulement si elles profitent aux plus défavorisés (\cle{Principe de différence}) et positions ouvertes à tous (\cle{Égalité des chances équitable}).
    \end{enumerate}
\end{methode}

\begin{exoresolu}[Typologie : Qualifier la peine de prison]
\textbf{Énoncé :} « La prison est-elle une vengeance de la société ou un outil de réinsertion ? » Qualifiez le type de justice concerné et opposez les théories de la peine.

\tcblower
\textbf{Solution détaillée :}
\begin{enumerate}
    \item \textbf{Identification :} La peine de prison relève de la \cle{Justice corrective (ou pénale / répressive)}. Elle sanctionne une violation de la loi pénale (atteinte à l'ordre public).
    \item \textbf{Opposition des finalités (Théories de la peine) :}
    \begin{itemize}
        \item \textbf{Théorie Rétributive (Vengeance légalisée / Expiation) :} \textbf{Kant} (\textit{Métaphysique des mœurs}) : « La peine est un châtiment infligé à un sujet parce qu'il a commis un crime. » Proportionnalité stricte (\textit{lex talionis} symbolique). La société « rend la monnaie de sa pièce ». \textbf{Risque :} Vengeance institutionnelle, déshumanisation.
        \item \textbf{Théorie Utilitaire / Préventive (Beccaria, Bentham) :} La peine n'a de sens que par ses \cle{conséquences futures}.
            \begin{itemize}
                \item \textbf{Prévention générale :} Faire peur aux autres (exemplarité).
                \item \textbf{Prévention spéciale :} Neutraliser le dangereux (incapacitation) / Soigner / Réinsérer.
            \end{itemize}
        \item \textbf{Synthèse moderne (Droit pénal camerounais / Code pénal) :} \cle{Double finalité} (Art. 1 Code pénal / principes directeurs) : \textbf{Répression} (sanctionner la faute, apaiser la victime/société) + \textbf{Réinsertion} (rééducation, formation, liberté conditionnelle). La prison \emph{doit} être un outil de réinsertion (peine individualisée, travail, formation), sans nier sa dimension punitive (justice pour la victime).
    \end{itemize}
    \item \textbf{Exemple concret :} Peine alternative (TIG - Travail d'Intérêt Général) au Cameroun. Évite la prison (école du crime), répare le dommage social (justice restaurative), prépare la réinsertion. Illustrer le passage de la justice rétributive pure à la justice corrective préventive/réparatrice.
\end{enumerate}
\textbf{Phrase de conclusion type :} « La prison, forme moderne de la justice corrective, oscille entre l'exigence rétributive de justice (payer sa dette) et la nécessité utilitaire de réinsertion ; le droit pénal contemporain tente de les concilier par l'individualisation de la peine. »
\end{exoresolu}

\begin{methode}[Problématiser : De la justice relative (positive) à l'idée absolue de Justice]
    \textbf{Objectif :} Traiter le thème classique de la perfectibilité du droit / critique de la justice positive.
    \begin{enumerate}
        \item \textbf{Constater la relativité de la Justice positive :}
            \begin{itemize}
                \item Variabilité spatio-temporelle (coutumes, codes, régimes politiques).
                \item Imperfection technique (lois lacunaires, rigueur du formalisme, erreurs judiciaires).
                \item Instrumentalisation possible (lois injustes, droit comme outil de domination - Marx).
            \end{itemize}
        \item \textbf{Identifier l'Idée absolue (Régulatrice) :}
            \begin{itemize}
                \item \textbf{Kant :} \cle{Idée régulatrice} de la raison pure. « Justice = ensemble des conditions sous lesquelles le choix de l'un peut s'accorder avec le choix de l'autre selon une loi universelle de liberté. » Inatteignable parfaitement, mais guide la législation.
                \item \textbf{Rawls :} \cle{Justice comme équité}. Principes choisis derrière le \cle{voile d'ignorance}. Idéal de justice sociale.
                \item \textbf{Droit naturel / Droits de l'homme :} Standards universels (DUDH 1948, Charte africaine) servant de \cle{critère de validité morale} au droit positif.
            \end{itemize}
        \item \textbf{Articuler le mouvement :} La Justice positive \textbf{tend vers} l'Idée absolue par le \cle{progrès} (abolition esclavage, peine de mort, droits femmes, justice transitionnelle). Le juriste/philosophe critique le droit \emph{au nom} de l'Idée.
        \item \textbf{Mots-clés :} \cle{Idéal régulateur}, \cle{Critique immanente}, \cle{Positivisme vs Jusnaturalisme}, \cle{Progrès moral}, \cle{Désobéissance civile} (Thoreau, Gandhi, King, Rawls) comme tension vers l'absolu.
    \end{enumerate}
\end{methode}

\begin{exoresolu}[L'Idée absolue comme critère de critique du Droit positif]
\textbf{Énoncé :} « La justice des hommes n'est-elle que l'ombre de la justice divine/idéale ? » Discutez en montrant le rôle de l'idée absolue de justice par rapport au droit positif.

\tcblower
\textbf{Solution détaillée :}
\begin{enumerate}
    \item \textbf{Thèse : L'ombre / L'imperfection (Relativisme / Positivisme / Réalisme).}
    \begin{itemize}
        \item \textbf{Kelsen (Positivisme pur) :} Nier l'accès à une « justice absolue » (métaphysique inaccessible). Seul le Droit positif existe. La justice = application de la loi. « Justice » n'est qu'un mot vide ou synonyme de légalité.
        \item \textbf{Marx (Matérialisme historique) :} Le Droit = superstructure reflétant l'infrastructure économique. Justice bourgeoise = justice de classe. L' « idée absolue » est une idéologie masquant la domination.
        \item \textbf{Anthropologie juridique :} Diversité radicale des justices (coutumières, religieuses, modernes). Pas d'universel concret.
    \end{itemize}
    \item \textbf{Antithèse : La lumière / Le guide (Jusnaturalisme / Idéalisme / Kant / Rawls).}
    \begin{itemize}
        \item \textbf{Kant :} L'Idée de Justice (Droit naturel) est une \cle{Idée régulatrice}. Elle ne se réalise jamais pleinement (royaume des fins), mais elle \textbf{commande} la législation. Sans elle, pas de critique possible du nazisme, de l'apartheid, de l'esclavage (légaux en leur temps).
        \item \textbf{Rawls :} Théorie de la justice = recherche de principes \textbf{universalisables} (voile d'ignorance). L'absolu = équité procédurale.
        \item \textbf{Droit international / DUDH :} L' « ombre » s'épaissit : les États signent des traités contraignants (justice pénale internationale, CPI). L'absolu se \cle{positivise} partiellement.
    \end{itemize}
    \item \textbf{Synthèse : Relation dialectique.}
    \begin{itemize}
        \item Le Droit positif est la \textbf{condition d'effectivité} de la Justice (sans loi, pas de sanction, arbitraire).
        \item L'Idée absolue est la \textbf{condition de légitimité et de perfectibilité} (sans idéal, pas de progrès, légalisme aveugle).
        \item \textbf{Exemple camerounais :} Justice coutumière vs Droit écrit vs Droit international (Droits de l'homme). La coutume (ex: exclusion femmes héritage) = justice positive locale. L'Idée absolue (Égalité sexes, CEDAW ratifiée) = levier de réforme (Code de la famille en projet, jurisprudence Cour Suprême).
    \end{itemize}
\end{enumerate}
\textbf{Phrase de conclusion type :} « L'idée absolue de justice n'est pas une chimère inaccessible, mais le phare sans lequel le navire du Droit positif s'échoue sur les récifs de la légalité barbare ; elle est la norme critique qui fait du Droit une histoire, non une nature. »
\end{exoresolu}

\begin{methode}[Appliquer les notions à une situation concrète (Cas pratique / Arrêt / Fait de société)]
    \textbf{Objectif :} Réussir l'exercice de « Dissertation sur texte » ou « Analyse de situation » (Épreuve écrite Bac).
    \begin{enumerate}
        \item \textbf{Lire / Identifier :} Quel est le \cle{conflit} ? Quels acteurs ? Quelle règle de droit (positive) ? Quelle valeur morale (justice) en jeu ?
        \item \textbf{Qualifier juridiquement :} Nommer les concepts : \cle{Responsabilité civile}, \cle{Abus de droit}, \cle{Inégalité de traitement}, \cle{Liberté d'expression vs Ordre public}, \cle{Justice distributive}, etc.
        \item \textbf{Problématiser :} Formuler la tension : « Jusqu'où le Droit positif (loi, contrat, coutume) peut-il s'écarter de la Justice (équité, droits fondamentaux) sans perdre sa légitimité ? »
        \item \textbf{Structurer l'analyse (Plan type) :}
            \begin{enumerate}
                \item \textbf{Exposé des faits / Thèse de l'auteur (si texte).}
                \item \textbf{Analyse conceptuelle :} Droit positif applicable vs Exigence de Justice.
                \item \textbf{Confrontation / Argumentation :} Arguments POUR le Droit (sécurité, prévisibilité) / Arguments POUR la Justice (équité, humanité, supériorité norme supérieure).
                \item \textbf{Résolution / Position personnelle argumentée :} Hiérarchie des normes (Constitution > Loi > Règlement ; Droit international > Droit interne). Proposition de solution juste \emph{et} juridique.
            \end{enumerate}
        \item \textbf{Rédiger :} Style impersonnel, connecteurs logiques (\textit{en effet, toutefois, par conséquent, dès lors}), vocabulaire précis.
    \end{enumerate}
\end{methode}

\begin{exoresolu}[Cas pratique : Licenciement abusif et clause de non-concurrence]
\textbf{Énoncé :} Un ingénieur camerounais signe un contrat avec une clause de non-concurrence de 5 ans sur tout le territoire national, sans contrepartie financière. Licencié, il retrouve un emploi chez un concurrent. L'ex-employeur saisit le tribunal. Analysez au regard du Droit du travail et de la Justice commutative.

\tcblower
\textbf{Solution détaillée :}
\begin{enumerate}
    \item \textbf{Qualification :} Contrat de travail (Droit positif : Code du travail Cameroun, Convention collective). Clause de non-concurrence = restriction de la \cle{liberté du travail} (Droit fondamental, Constitution, Art. 1er DUDH).
    \item \textbf{Analyse du Droit Positif (Légalité) :} Validité clause : 3 conditions cumulatives (Jurisprudence constante Cour Suprême et principes du Code du travail) :
    \begin{enumerate}
        \item Indispensable pour protéger intérêts légitimes entreprise.
        \item Limitée dans l'espace (zone géographique) et le temps (durée maxi souvent 2 ans).
        \item \textbf{Assortie d'une \cle{contrepartie financière}} (indemnité compensatrice).
    \end{enumerate}
    \item \textbf{Application au cas :} Clause = 5 ans (excessif), tout territoire (excessif), \textbf{pas de contrepartie financière}. $\rightarrow$ \textbf{Clause nulle et non avenue} (Droit positif).
    \item \textbf{Analyse par la Justice Commutative (Équité) :} Échange initial : Travail vs Salaire. Clause = charge nouvelle pour salarié (interdiction travailler) sans avantage nouveau (pas de paiement). \textbf{Déséquilibre contractuel} flagrant. Abus de droit (pouvoir de négociation inégal).
    \item \textbf{Justice Distributive / Sociale :} Protection du partie faible (salarié) contre clause léonine. Ordre public social.
    \item \textbf{Conclusion juridique :} Tribunal déboute l'employeur. Salarié libre de travailler. Éventuels dommages-intérêts pour procédure abusive.
    \item \textbf{Conclusion philosophique :} Le Droit positif (nullité clause) \textbf{réalise} ici la Justice commutative (équilibre contractuel) et la Justice sociale (protection faible). Convergence Droit / Justice.
\end{enumerate}
\textbf{Phrase de conclusion type :} « La nullité de la clause illicite illustre la fonction correctrice du Droit positif : il sanctionne le déséquilibre commutatif pour restaurer la justice contractuelle, prouvant que la technique juridique n'est pas étrangère à l'éthique. »
\end{exoresolu}

\begin{methode}[Rédiger une dissertation philosophique (Plan dialectique / Thématique)]
    \textbf{Objectif :} Structurer une réponse de 3h (Bac) : Introduction, Développement (3 parties), Conclusion.
    \begin{enumerate}
        \item \textbf{Analyse du sujet (15 min) :} Définir termes, identifier présupposés, formuler \cle{problématique} (question tension), annoncer \cle{plan}.
        \item \textbf{Introduction (rédaction soignée) :}
            \begin{itemize}
                \item \textbf{Amorce :} Citation, fait d'actualité, définition étymologique, opposition courante.
                \item \textbf{Définitions :} Droit, Justice (selon sujet).
                \item \textbf{Problématique :} « Le Droit positif suffit-il à réaliser la Justice ? » / « L'égalité formelle garantit-elle la justice sociale ? »
                \item \textbf{Annonce du plan :} « Nous verrons d'abord que... (Thèse), puis que... (Antithèse), pour montrer enfin que... (Synthèse). »
            \end{itemize}
        \item \textbf{Développement (3 grandes parties = 3 paragraphes majeurs) :}
            \begin{itemize}
                \item \textbf{Partie I (Thèse) :} Argument 1 (Auteur/Concept) + Exemple. Argument 2 + Exemple. \textit{Mini-transition}.
                \item \textbf{Partie II (Antithèse) :} Objection radicale. Argument 1 + Exemple. Argument 2 + Exemple. \textit{Mini-transition}.
                \item \textbf{Partie III (Synthèse / Dépassement) :} Réconciliation / Hiérarchisation / Condition / Distinction cruciale (ex: Légalité/Légitimité, Formel/Matériel, Positif/Naturel). C'est la partie la plus valorisée.
            \end{itemize}
        \item \textbf{Conclusion :} Réponse directe à la problématique. Ouverture (autre domaine, actualité, auteur).
        \item \textbf{Relecture :} Orthographe, connecteurs, noms d'auteurs exacts, concepts soulignés.
    \end{enumerate}
\end{methode}

\begin{exoresolu}[Dissertation type : « La loi est-elle la mesure de la justice ? »]
\textbf{Énoncé :} Sujet de dissertation. Rédigez l'introduction complète et le plan détaillé avec arguments principaux et exemples.

\tcblower
\textbf{Solution détaillée :}

\textbf{INTRODUCTION}
\textit{[Amorce]} « La loi est la raison sans passion », écrit Aristote (\textit{Politique}). Pourtant, l'histoire retiendra que les pires crimes furent commis « au nom de la loi » : lois raciales de Nuremberg (1935), apartheid légal en Afrique du Sud, Code noir.
\textit{[Définitions]} \textbf{Loi (Droit positif) :} Règle générale, abstraite, impersonnelle, votée par le Parlement, sanctionnée par l'État. Caractère \cle{coercitif} et \cle{formel}. \textbf{Justice :} Idéal moral d'équité, rendre à chacun son dû (\textit{suum cuique}), respect de la dignité. Caractère \cle{matériel} et \cle{universel}.
\textit{[Problématique]} La conformité à la loi (légalité) suffit-elle à garantir la conformité au juste (légitimité) ? En d'autres termes, \textbf{la loi est-elle la mesure de la justice, ou la justice est-elle la mesure de la loi ?}
\textit{[Plan]} Nous montrerons d'abord que la loi \textbf{apparaît} comme la mesure de la justice par sa forme universelle et sa force contraignante (I). Nous verrons ensuite que l'histoire et la philosophie révèlent l'\textbf{insuffisance} de la loi face aux exigences de la justice (II). Enfin, nous établirons que la justice, comme \textbf{idée régulatrice}, doit demeurer le critère suprême de validité morale de la loi (III).

\textbf{DÉVELOPPEMENT (Plan détaillé)}

\textbf{I. La Loi : Mesure formelle et effective de la Justice (Thèse : Légalité = Justice)}
\begin{itemize}
    \item \textbf{A. L'universalité de la loi assure l'égalité formelle (Justice commutative).} Loi générale $\rightarrow$ même traitement pour tous (Art. 6 DDHC 1789). Ex: Code civil, Code pénal. Évite l'arbitraire (justice du prince).
    \item \textbf{B. La loi institue la justice distributive (Justice sociale).} Loi fiscale (progressivité), Code du travail (SMIG, congés), Sécurité sociale. La loi \textbf{réalise} la justice par la redistribution (Rawls : institutions justes).
    \item \textbf{C. La sanction étatique garantit l'effectivité.} Sans loi écrite + force publique, la justice reste vœu pieu (Hobbes, Kelsen). \cle{Sécurité juridique}.
\end{itemize}
$\rightarrow$ \textit{Transition :} Mais la forme légale peut masquer l'injustice matérielle. La loi peut être l'instrument de l'injustice.

\textbf{II. Les limites de la Loi : Quand la légalité contredit la Justice (Antithèse : Légalité $\neq$ Légitimité)}
\begin{itemize}
    \item \textbf{A. Le positivisme juridique peut valider l'injustice (Kelsen / Hart).} Séparation Droit/Morale. « Lex injusta est lex » (loi injuste reste loi valide). Ex: Lois de Vichy, Nuremberg, Ségrégation US. \cle{Légalisme} aveugle.
    \item \textbf{B. La rigidité de la loi vs l'équité du cas particulier (Aristote, \textit{Éthique à Nicomaque} V).} La loi est universelle, les cas sont particuliers. \cle{Équité (epieikeia)} = correction de la loi quand elle est défectueuse par sa généralité. Ex: Juge des référés, pouvoir d'appréciation, circonstances atténuantes.
    \item \textbf{C. La loi comme outil de domination (Marx, Foucault).} Droit = superstructure. Justice bourgeoise protège la propriété (vol du pain puni, vol de salaire impuni). \cle{Justice de classe}.
\end{itemize}
$\rightarrow$ \textit{Transition :} Si la loi ne se suffit pas, faut-il la rejeter ? Non, la corriger par un principe supérieur.

\textbf{III. La Justice Idéale : Critère régulateur et source de légitimité de la Loi (Synthèse : Hiérarchie Normative)}
\begin{itemize}
    \item \textbf{A. L'Idée de Justice comme horizon (Kant, Rawls).} \cle{Idée régulatrice} (Kant) : on ne l'atteint jamais, mais elle guide la législation. \cle{Voile d'ignorance} (Rawls) : test de justice des lois.
    \item \textbf{B. Le contrôle de constitutionnalité et conventionalité.} Hiérarchie des normes : Constitution (bloc de constitutionnalité : DDHC 1789, Préambule 1946, Charte de l'environnement 2004) > Loi. Conseil Constitutionnel. Juge = gardien de la justice vs loi.
    \item \textbf{C. La désobéissance civile (Thoreau, Gandhi, King, Rawls).} Dernier recours quand loi \textbf{manifestement} injuste viole droits fondamentaux. Condition : publique, non-violente, acceptation peine. Preuve que Justice > Loi.
\end{itemize}

\textbf{CONCLUSION}
\textbf{Réponse :} La loi n'est pas \emph{la} mesure de la justice (elle peut être injuste), mais elle en est l'\textbf{instrument nécessaire}. La justice est la \textbf{mesure de la loi} : critère de validité morale, fin ultime du droit.
\textbf{Ouverture :} Aujourd'hui, l'IA juge (justice prédictive) : la loi (algorithme) peut-elle rendre la justice sans l'humain (équité, contexte) ? Le défi de la \cle{justice algorithmique}.
\end{exoresolu}

\begin{attention}[Les 5 erreurs qui coûtent des points au Bac]
\begin{enumerate}
    \item \textbf{Confondre Droit et Justice (ou Légalité et Légitimité).} Dire « C'est juste car c'est la loi » (positivisme naïf) ou « C'est illégal donc injuste » (sans nuance). \textbf{Reflexe :} Toujours distinguer la \cle{validité formelle} (procédure respectée) de la \cle{validité morale} (conformité à l'équité/DDH).
    \item \textbf{Ignorer la distinction Justice Commutative / Distributive.} Traiter un problème de répartition (salaire, bourse, impôt) avec des arguments d'échange (contrat, mérite individuel strict) ou inversement. \textbf{Reflexe :} Identifier : \textit{Échange entre deux ? $\rightarrow$ Commutative (Arithmétique). Partage par la collectivité ? $\rightarrow$ Distributive (Géométrique).}
    \item \textbf{Réduire la Justice Pénale à la Vengeance / Rétribution seule.} Oublier les finalités préventives (réinsertion, prévention générale) et le principe de proportionnalité. \textbf{Reflexe :} Citer \textbf{Beccaria} (utilitarisme) et le \textbf{Code pénal moderne} (double finalité : sanction + réinsertion).
    \item \textbf{Manquer de références auteurs / concepts précis.} « Les philosophes disent que... », « Selon la justice... ». \textbf{Reflexe :} Citer \textbf{noms propres} (Aristote, Kant, Rawls, Kelsen, Marx, Hobbes, Locke, Rousseau, Ulpien) et \textbf{concepts techniques} (Grundnorm, Voile d'ignorance, Volonté générale, Équité, Jusnaturalisme, Positivisme, Contractualisme).
    \item \textbf{Mauvaise structure de la dissertation :} Pas de problématique explicite, plan en « liste de courses » (I. A dit, II. B dit, III. C dit) sans dialectique (Thèse/Antithèse/Synthèse), conclusion qui ne répond pas à la question. \textbf{Reflexe :} \textit{Problématiser $\rightarrow$ Plan dialectique $\rightarrow$ Synthèse qui tranche/hiérarchise $\rightarrow$ Conclusion directe + Ouverture.}
\end{enumerate}
\end{attention}

\newpage

\coursec{Exercices}

\courssub{Je vérifie mes connaissances}

\begin{exercice}[QCM : concepts fondamentaux]
Parmi les propositions suivantes, une seule est exacte. Choisis-la et justifie brièvement ton choix.
\begin{enumerate}
\item Le droit positif désigne l'ensemble des règles morales universelles reconnues par la raison.
\item La justice commutative régit les rapports entre l'individu et la communauté politique.
\item Le droit naturel postule l'existence de principes juridiques antérieurs et supérieurs aux lois écrites.
\item L'égalité proportionnelle consiste à traiter tous les individus de manière identique quelles que soient leurs différences.
\end{enumerate}
\end{exercice}

\begin{exercice}[Vrai ou Faux justifié]
Indique si chaque affirmation est vraie ou fausse, puis justifie ta réponse en une phrase en mobilisant les notions du cours.
\begin{enumerate}
\item « Toute loi promulguée selon la procédure légale est nécessairement juste. »
\item « La justice distributive s'applique uniquement au partage des peines pénales. »
\item « Le droit coutumier n'a aucune valeur juridique face au droit écrit moderne. »
\item « L'idée absolue de justice sert de critère pour juger la validité des lois positives. »
\end{enumerate}
\end{exercice}

\begin{exercice}[Définitions précises]
Donne une définition rigoureuse de chacun des concepts suivants en deux lignes maximum, en distinguant le genre et la différence spécifique :
\begin{enumerate}
\item Droit positif
\item Justice distributive
\item Égalité arithmétique
\item Équité (au sens aristotélicien)
\end{enumerate}
\end{exercice}

\begin{exercice}[Correspondances : formes de justice et situations]
Relie chaque situation à la forme de justice qui lui correspond le mieux (justice commutative, justice distributive, justice réparatrice / corrective). Justifie chaque rapprochement en une phrase.
\begin{enumerate}
\item Un juge condamne un voleur à restituer l'objet volé et à verser des dommages-intérêts.
\item L'État attribue des bourses d'études en fonction des mérites scolaires et des ressources des familles.
\item Deux commerçants concluent un contrat d'achat-vente au prix convenu.
\item Un tribunal ordonne la réparation d'un préjudice corporel causé par un accident de la route.
\end{enumerate}
\end{exercice}

\courssub{Je m'entraîne}

\begin{exercice}[Analyse d'un conflit foncier au Cameroun]
Dans un village de la région du Centre, un litige oppose deux familles pour la propriété d'un champ de cacao. Le chef traditionnel tranche selon la coutume : le champ revient à la famille qui l'exploite depuis trois générations. L'une des familles saisit le tribunal de première instance qui rend un jugement basé sur le titre foncier enregistré au nom de l'autre famille.
\begin{enumerate}
\item Identifie les deux ordres de droit en présence (droit coutumier / droit positif écrit).
\item Explique en quoi la décision du chef traditionnel relève de la justice distributive ou commutative.
\item Discute la légitimité respective des deux décisions : quelle(s) source(s) de légitimité pour chacune ?
\item Propose une articulation possible entre ces deux ordres juridiques pour résoudre le conflit durablement.
\end{enumerate}
\end{exercice}

\begin{exercice}[Classement de situations juridiques]
Classe chacune des dix situations suivantes dans l'une des trois formes de justice (commutative, distributive, réparatrice). Justifie chaque classement en une phrase précise.
\begin{enumerate}
\item Signature d'un contrat de travail entre un employeur et un salarié.
\item Attribution des marchés publics par appel d'offres selon des critères de compétence et de prix.
\item Condamnation d'un conducteur ivre à indemniser sa victime.
\item Répartition de l'impôt sur le revenu selon un barème progressif.
\item Échange d'un bien contre son prix au marché.
\item Versement d'une pension de retraite calculée sur les cotisations versées.
\item Partage d'un héritage entre enfants selon la coutume (part aînée / parts cadets).
\item Sanction disciplinaire (exclusion temporaire) pour un élève ayant frappé un camarade.
\item Remboursement d'un prêt bancaire selon l'échéancier convenu.
\item Attribution de logements sociaux selon les revenus et la taille du ménage.
\end{enumerate}
\end{exercice}

\begin{exercice}[Droit naturel vs droit positif : le débat de l'antigone]
Lis l'extrait suivant (résumé) : « Antigone enterre son frère Polynice malgré l'interdiction du roi Créon. Elle affirme obéir aux "lois non écrites et immuables des dieux" supérieures aux décrets des hommes. »
\begin{enumerate}
\item Identifie la thèse soutenue par Antigone (droit naturel / droit positif).
\item Explique en quoi Créon incarne le positivisme juridique strict.
\item Selon la philosophie du droit naturel, une loi injuste est-elle une loi ? Justifie.
\item Donne un exemple historique ou contemporain (camerounais ou international) où des citoyens ont désobéi à une loi positive au nom d'un droit supérieur.
\end{enumerate}
\end{exercice}

\begin{exercice}[Égalité et justice : le cas des quotas]
L'État camerounais met en place une politique de discrimination positive : 30\% des places dans les grandes écoles sont réservées aux candidates filles et aux élèves des régions défavorisées (Extrême-Nord, Nord, Adamaoua, Est).
\begin{enumerate}
\item Cette mesure relève-t-elle de l'égalité arithmétique ou de l'égalité proportionnelle ? Justifie.
\item Montre en quoi elle peut être considérée comme une application de la justice distributive (critère du mérite, du besoin, de la contribution ?).
\item Un candidat garçon d'une région favorisée, refusé malgré une note supérieure à celle d'une candidate admise au titre du quota, crie à l'injustice. Analyse son argument au prisme de la distinction entre égalité formelle et équité réelle.
\item Ton avis argumenté : cette mesure est-elle juste ? Pourquoi ?
\end{enumerate}
\end{exercice}

\courssub{Je me prépare au Bac}

\begin{exercice}[Dissertation type Baccalauréat technique]
\textbf{Sujet :} « La loi est-elle toujours juste ? »
\begin{enumerate}
\item Rédige une introduction complète (accroche, définition des termes, problématique, annonce du plan).
\item Construis un plan détaillé en trois parties (avec deux sous-parties chacune) qui répond à la problématique. Pour chaque sous-partie, indique l'argument principal, l'auteur ou la référence philosophique mobilisée, et un exemple concret (camerounais si possible).
\item Rédige la conclusion (synthèse, ouverture).
\end{enumerate}
\end{exercice}

\begin{exercice}[Explication de texte : Aristote, Éthique à Nicomaque, Livre V]
\textbf{Texte :} « Il y a donc deux sortes de justice politique : l'une, la justice distributive, concerne la répartition des honneurs, de l'argent et des autres biens que la cité partage entre ceux qui ont part à la constitution ; l'autre, la justice corrective, s'exerce dans les transactions entre particuliers. [...] La justice distributive suppose une égalité proportionnelle : c'est l'égalité selon le mérite. Car tous ne conviennent pas sur ce qu'est le mérite : les démocrates l'identifient à la liberté, les oligarques à la richesse, les aristocrates à la vertu. »
\begin{enumerate}
\item Dégage la thèse principale du texte en une phrase.
\item Identifie et explique les deux formes de justice distinguées par Aristote.
\item Qu'est-ce que l'« égalité proportionnelle » ? Pourquoi Aristote dit-il que « tous ne conviennent pas sur ce qu'est le mérite » ?
\item Illustrer la justice distributive par un exemple camerounais actuel (ex. : attribution des bourses, recrutement à la fonction publique, répartition des revenus pétroliers). Montre comment le critère de mérite choisi influence la justice de la répartition.
\item Le texte suggère-t-il qu'il existe un critère unique et objectif du mérite ? Justifie ta réponse par le texte.
\end{enumerate}
\end{exercice}

\begin{exercice}[Sujet de débat argumenté type Probatoire]
\textbf{Sujet :} « Faut-il traiter tout le monde de la même manière pour être juste ? »
\begin{enumerate}
\item Présente en quelques lignes les deux thèses opposées (thèse de l'égalité formelle / thèse de l'équité proportionnelle).
\item Construis une argumentation personnelle structurée en trois temps (thèse, antithèse, synthèse) pour répondre à la question. Ton développement doit faire environ 20 lignes, mobiliser les concepts de justice commutative, distributive, égalité arithmétique et proportionnelle, et citer au moins un philosophe (Aristote, Rawls, Ricœur, etc.).
\item Conclus en précisant dans quels contextes (famille, école, travail, État) l'un ou l'autre principe te semble prioritaire.
\end{enumerate}
\end{exercice}

\newpage

\coursec{Corrigés des exercices}

\coursec{Le Droit et la Justice}
\courssub{Exercices d'application et sujets types (Probatoire / Baccalauréat technique)}

\begin{corrige}[Exercice 1 — QCM : concepts fondamentaux]
\cle{Bonne réponse : proposition 3.}

\begin{enumerate}
\item \textbf{Faux.} Le \cle{droit positif} désigne l'ensemble des règles de droit \emph{écrites}, promulguées par l'autorité étatique légitime (Parlement, Gouvernement) et en vigueur dans un État à un moment donné. Ce sont des règles \emph{historiques} et \emph{contingentes}, non des règles morales universelles (qui relèvent du droit naturel).
\item \textbf{Faux.} La \cle{justice commutative} (ou corrective) régit les rapports entre \emph{particuliers} considérés comme égaux (contrats, délits, restitution). C'est la \cle{justice distributive} qui régit les rapports entre l'individu et la communauté politique (répartition des honneurs, charges, richesses).
\item \textbf{Vrai.} Le \cle{droit naturel} postule l'existence de principes juridiques universels, immuables et antérieurs à l'État (fondés sur la nature humaine, la raison ou la volonté divine) qui s'imposent au législateur positif. Une loi positive contraire au droit naturel est jugée injuste, voire illégitime (\emph{lex injusta non est lex}).
\item \textbf{Faux.} L'\cle{égalité proportionnelle} (ou géométrique) consiste à traiter les individus \emph{différemment} en fonction de critères pertinents (mérite, besoin, contribution, rang) : « à chacun selon son mérite ». C'est l'\cle{égalité arithmétique} qui consiste à traiter tout le monde de manière identique (un pour un), quelles que soient les différences.
\end{enumerate}
\end{corrige}

\begin{corrige}[Exercice 2 — Vrai ou Faux justifié]
\begin{enumerate}
\item \textbf{Faux.} La \cle{légalité} (respect de la procédure de promulgation) ne garantit pas la \cle{légitimité} morale ni la \cle{justice} de la loi. Le positivisme juridique (Kelsen, Hart) sépare la validité formelle de la valeur morale. Exemple : les lois de ségrégation raciale (USA, Afrique du Sud) ou les lois vichystes (France) étaient légales mais injustes.
\item \textbf{Faux.} La \cle{justice distributive} concerne la répartition des \emph{biens, honneurs, charges et ressources} au sein de la communauté (impôts, bourses, marchés publics, santé). Le partage des \emph{peines pénales} relève de la \cle{justice réparatrice (corrective)} qui sanctionne une violation de la loi pour rétablir l'ordre troublé.
\item \textbf{Faux.} Au Cameroun comme dans de nombreux systèmes mixtes, le \cle{droit coutumier} est une source de droit reconnue par la Constitution (préambule) et le Code civil (art. 1er al. 2 : « La coutume s'applique... pourvu qu'elle ne soit pas contraire à la loi, à l'ordre public et aux bonnes mœurs »). Il a valeur juridique dans les matières de statut personnel, successions, foncier rural, tant qu'il n'est pas contredit par une loi écrite.
\item \textbf{Vrai.} L'\cle{idée absolue de justice} (idéal régulateur, droit naturel, principes supérieurs type DDH) fonctionne comme un \cle{critère critique} : elle permet de juger la validité \emph{morale} des lois positives et d'orienter les réformes. C'est le sens de la maxime d'Augustin : \emph{« Lex injusta non est lex »} (une loi injuste n'est pas une loi).
\end{enumerate}
\end{corrige}

\begin{corrige}[Exercice 3 — Définitions précises]
\begin{enumerate}
\item \textbf{Droit positif :} Ensemble des règles de droit \emph{écrites} et \emph{coercitives}, édictées par l'autorité étatique compétente (législateur, pouvoir réglementaire), en vigueur dans un État déterminé à un moment donné. \textit{(Genre : ensemble de règles juridiques ; Différence : positivité, étaticité, coercitivité, historicité).}
\item \textbf{Justice distributive :} Forme de justice politique qui régit la \emph{répartition proportionnelle} des biens, honneurs, charges et services publics entre les membres de la communauté, selon un critère de mérite, de besoin, de rang ou de contribution. \textit{(Genre : espèce de justice politique ; Différence : objet = répartition communautaire, mode = proportionnalité géométrique).}
\item \textbf{Égalité arithmétique :} Principe d'égalité quantitative qui attribue à chacun la \emph{même part} (un pour un), sans tenir compte des différences de mérite, de besoin, de statut ou de contribution. \textit{(Genre : mode d'égalité ; Différence : identité stricte des parts, aveuglement aux différences pertinentes).}
\item \textbf{Équité (sens aristotélicien) :} Vertu correctrice de la loi générale : elle permet au juge d'\emph{adapter la règle universelle au cas particulier} lorsque l'application littérale de la loi aboutirait à une injustice par excès de généralité ; elle vise l'esprit de la loi plutôt que sa lettre. \textit{(Genre : vertu de justice / correction ; Différence : particularisation, supra-légalité, recherche du juste concret).}
\end{enumerate}
\end{corrige}

\begin{corrige}[Exercice 4 — Correspondances : formes de justice et situations]
\begin{enumerate}
\item \textbf{Justice réparatrice (corrective).} Le vol rompt l'égalité initiale entre la victime et le voleur ; la condamnation à restituer l'objet et verser des dommages-intérêts vise à \emph{rétablir l'égalité arithmétique} lésée (status quo ante) en sanctionnant l'auteur du délit.
\item \textbf{Justice distributive.} L'État, en tant que communauté politique, répartit un bien public rare (les bourses) entre les citoyens selon des critères de \emph{proportionnalité} (mérite scolaire = mérite ; ressources = besoin), conformément à la justice distributive.
\item \textbf{Justice commutative.} Le contrat d'achat-vente est l'archétype de l'échange entre \emph{particuliers égaux} : le prix convenu est l'équivalent de la valeur du bien, réalisant l'\emph{égalité arithmétique} (donner/recevoir de même valeur).
\item \textbf{Justice réparatrice (corrective).} L'accident crée une inégalité (préjudice corporel) entre le responsable et la victime ; l'ordonnance de réparation (indemnisation) vise à \emph{réparer le dommage} et à rétablir l'équilibre patrimonial rompu par le fait dommageable.
\end{enumerate}
\end{corrige}

\begin{corrige}[Exercice 5 — Analyse d'un conflit foncier au Cameroun]
\begin{enumerate}
\item \textbf{Deux ordres de droit :}
      \begin{itemize}
      \item \cle{Droit coutumier} : invoqué par le chef traditionnel, fondé sur la coutume locale (exploitation ancestrale, transmission lignagère, droit de la hache), non écrit, sanctionné par l'autorité traditionnelle.
      \item \cle{Droit positif écrit (droit moderne)} : invoqué devant le tribunal, fondé sur la loi foncière (loi n°74-1 du 6 juillet 1974 modifiée, ordonnance n°74-2), le titre foncier (certificat de propriété), l'immatriculation, sanctionné par l'État.
      \end{itemize}
\item \textbf{Nature de la décision du chef :} Elle relève de la \cle{justice distributive}. Le chef ne juge pas un contrat ou un délit entre deux parties égales (commutative), mais \emph{répartit un bien communautaire} (le champ) entre deux familles en appliquant un critère de \emph{proportionnalité coutumière} : l'ancienneté de l'exploitation (trois générations) et l'appartenance lignagère, considérés comme des « mérites » coutumiers justifiant l'attribution.
\item \textbf{Légitimité respective :}
      \begin{itemize}
      \item \emph{Chef traditionnel :} Légitimité \cle{traditionnelle} (Weber) : ancrage dans l'histoire, acceptation communautaire, efficacité sociale (maintien de la paix lignagère), sacralité de la parole des ancêtres. Source : la coutume vivante.
      \item \emph{Tribunal :} Légitimité \cle{légale-rationnelle} (Weber) : compétence d'attribution (loi), procédure contradictoire, force exécutoire de l'État (monopole de la contrainte), sécurité juridique (publicité foncière). Source : la loi écrite promulguée.
      \end{itemize}
\item \textbf{Articulation possible :} Le droit camerounais prévoit déjà des passerelles :
      \begin{itemize}
      \item \cle{Médiation/conciliation préalable} obligatoire ou facultative devant le chef ou un médiateur assermenté (loi sur l'organisation judiciaire).
      \item \cle{Reconnaissance des droits coutumiers} dans la procédure d'immatriculation (enquête parcellaire, procès-verbal de bornage contradictoire impliquant le chef).
      \item \cle{Titre foncier « coutumier »} ou inscription des droits d'usage coutumiers sur le titre (droit de superficie, servitudes).
      \item \cle{Juridictions de premier degré} (tribunaux coutumiers) dont les décisions sont susceptibles d'appel devant les tribunaux de grande instance, assurant un contrôle de conformité à l'ordre public et aux lois écrites.
      \end{itemize}
      Une solution durable : le tribunal sursoit à statuer, ordonne une médiation coutumière pour identifier les droits réels de chaque famille, puis homologue l'accord ou tranche en intégrant la réalité coutumière dans le titre foncier (ex. : copropriété, droit d'usage).
\end{enumerate}
\end{corrige}

\begin{corrige}[Exercice 6 — Classement de situations juridiques]
\begin{center}
\begin{tabularx}{\linewidth}{|l|X|l|X|}\hline
\textbf{N°} & \textbf{Situation} & \textbf{Forme de justice} & \textbf{Justification} \\ \hline
1 & Signature contrat de travail & \cle{Commutative} & Échange de consentements entre deux parties égales (employeur/salarié) : travail contre salaire (égalité arithmétique valeur/valeur). \\ \hline
2 & Attribution marchés publics & \cle{Distributive} & L'État (communauté) répartit un bien public (marché) entre candidats selon critères de proportionnalité (compétence, prix = mérite/économie). \\ \hline
3 & Condamnation conducteur ivre & \cle{Réparatrice (corrective)} & Rétablit l'égalité rompue par le délit : l'auteur répare le préjudice causé à la victime (indemnisation = correction de l'inégalité). \\ \hline
4 & Répartition impôt progressif & \cle{Distributive} & Répartition de la charge commune (impôt) selon la capacité contributive (proportionnalité géométrique : plus on a, plus on donne). \\ \hline
5 & Échange bien contre prix & \cle{Commutative} & Échange volontaire entre particuliers sur base d'équivalence valeur/prix (égalité arithmétique). \\ \hline
6 & Versement pension retraite (CNPS) & \cle{Distributive} & Système par répartition : solidarité intergénérationnelle, répartition des cotisations actuelles vers retraités selon critères (ancienneté, salaire) = proportionnalité. \\ \hline
7 & Partage héritage coutumier & \cle{Distributive} & Répartition d'un patrimoine entre ayants droit selon critères coutumiers de proportionnalité (rang, sexe, rôle : part aînée vs cadets). \\ \hline
8 & Sanction exclusion élève & \cle{Réparatrice (corrective)} & Réparation du trouble à l'ordre scolaire et de l'atteinte à l'intégrité de la victime ; sanction proportionnée à la faute (corrective). \\ \hline
9 & Remboursement prêt bancaire & \cle{Commutative} & Exécution d'un contrat de prêt : restitution du capital + intérêts selon échéancier = respect de l'égalité arithmétique contractuelle. \\ \hline
10 & Attribution logements sociaux & \cle{Distributive} & Répartition d'un bien rare (logement) par la collectivité selon critères de besoin (revenus, taille ménage) = proportionnalité géométrique. \\ \hline
\end{tabularx}
\end{center}
\end{corrige}

\begin{corrige}[Exercice 7 — Droit naturel vs droit positif : le débat de l'Antigone]
\begin{enumerate}
\item \textbf{Thèse d'Antigone :} \cle{Droit naturel}. Elle affirme l'existence de « lois non écrites et immuables des dieux » (lois divines, lois de la nature, droit naturel) qui sont \emph{antérieures} et \emph{supérieures} aux décrets des hommes (droit positif). Elle leur obéit par devoir moral absolu.
\item \textbf{Créon et le positivisme :} Créon incarne le \cle{positivisme juridique strict} : la loi (\emph{nomos}) émane de la volonté du souverain (le roi) ; elle est la seule source du droit ; sa validité tient à sa promulgation par l'autorité légitime, non à sa conformité morale ; l'obéissance est due \emph{inconditionnellement} pour assurer l'ordre et la stabilité de la Cité (\emph{« Nul n'est assez puissant pour être au-dessus de la loi »} — mais ici la loi = volonté du roi).
\item \textbf{Droit naturel et loi injuste :} Pour le droit naturel classique (Saint Augustin, Saint Thomas d'Aquin), \textbf{une loi injuste n'est pas une loi} (\emph{lex injusta non est lex}). Une loi positive qui contredit le droit naturel (loi éternelle, loi divine, raison droite) est une \emph{corruption de loi} (\emph{lex corrupta}) ; elle ne lie pas en conscience (\emph{non obligat in foro conscientiae}), sauf à l'obéissance par prudence pour éviter un plus grand désordre (scandale, anarchie).
\item \textbf{Exemple historique/contemporain :}
      \begin{itemize}
      \item \emph{International :} \textbf{Rosa Parks / Martin Luther King} (USA, 1955-68) : désobéissance civile aux lois de ségrégation raciale (lois Jim Crow) au nom de la Constitution (14e amendement) et de la loi morale divine/universelle (droits de l'homme).
      \item \emph{Cameroun :} \textbf{Les maquisards de l'UPC (années 1950-60)} : lutte armée et désobéissance à l'ordre colonial français (lois d'exception, code de l'indigénat) au nom du droit des peuples à l'autodétermination (droit naturel/universel) et de l'indépendance.
      \end{itemize}
\end{enumerate}
\end{corrige}

\begin{corrige}[Exercice 8 — Égalité et justice : le cas des quotas]
\begin{enumerate}
\item \textbf{Égalité proportionnelle.} La mesure ne traite pas tous les candidats de manière identique (même note = même sort) : elle \emph{différencie} le traitement en fonction de critères jugés pertinents (sexe, origine géographique défavorisée) pour corriger des inégalités structurelles de départ. C'est l'essence de l'égalité proportionnelle : « à chacun selon ses besoins / selon sa situation de désavantage ».
\item \textbf{Justice distributive : critère du \cle{besoin} et de la \cle{réparation}.} La justice distributive répartit des biens rares (places en grandes écoles) selon un critère de mérite. Ici, le « mérite » est redéfini : ce n'est pas seulement la note brute (mérite scolaire acquis), mais le \emph{potentiel méritoire dans un contexte défavorisé} (besoin d'équité) et la \emph{réparation d'une injustice historique} (sous-représentation des filles et des régions Nord/Est/Adamaoua/Extrême-Nord). C'est une application du principe de \cle{différence} de Rawls : les inégalités sont justes si elles profitent aux plus défavorisés.
\item \textbf{Analyse de l'argument du candidat :} Il invoque l'\cle{égalité formelle} (égalité arithmétique) : « même note = même droit », la loi doit être aveugle aux différences (sexe, région). \emph{Analyse :} Cette égalité formelle ignore les \emph{inégalités réelles de départ} (qualité de l'encadrement scolaire, stéréotypes de genre, pauvreté, distance). L'\cle{équité réelle} (justice proportionnelle) exige de traiter différemment des situations différentes pour leur donner une chance égale de réussite (égalité des chances \emph{réelle} vs formelle). Le candidat confond « justice = même traitement » et « justice = traitement proportionné à la situation ».
\item \textbf{Avis argumenté :} \textbf{Oui, cette mesure est juste \emph{provisoirement} et \emph{conditionnellement}.}
      \begin{itemize}
      \item \emph{Pour :} Elle corrige une injustice structurelle avérée (sous-représentation), favorise la diversité (enrichissement mutuel), respecte le principe de différence (Rawls), vise l'égalité réelle des chances (but constitutionnel).
      \item \emph{Contre (vigilance) :} Risque de stigmatisation des bénéficiaires (« quota »), sentiment d'injustice chez les exclus (perte de légitimité), ne dispense pas de réformes structurelles (écoles primaires au Nord, lutte contre stéréotypes).
      \item \emph{Condition :} Doit être \emph{temporaire} (mesure transitoire), \emph{ciblée} (seuils révisables), \emph{accompagnée} de mesures amont (investissement éducatif dans les régions concernées). Sans cela, elle devient une injustice nouvelle (discrimination positive inversée permanente).
      \end{itemize}
\end{enumerate}
\end{corrige}

\begin{corrige}[Exercice 9 — Dissertation type Baccalauréat technique : « La loi est-elle toujours juste ? »]
\textbf{Barème indicatif (sur 20) :} Introduction (3) + Développement 3 parties $\times$ 2 sous-parties (12) + Conclusion (3) + Qualité langue/structure (2).

\begin{enumerate}
\item \textbf{Introduction complète :}
      \begin{itemize}
      \item \emph{Accroche :} « Antigone enterre son frère malgré l'interdiction de Créon : elle obéit à une loi supérieure, non écrite, que la loi de l'État ne peut abolir. » (Sophocle, \emph{Antigone}).
      \item \emph{Définitions :} \cle{Loi} = règle générale, abstraite, impersonnelle, obligatoire, sanctionnée par la force publique, émanant de l'autorité légitime (droit positif). \cle{Juste} = ce qui est conforme au droit, à l'équité, à la morale, à l'idée de justice (donner à chacun ce qui lui revient).
      \item \emph{Problématique :} La loi positive possède la force coercitive de l'État, mais sa validité formelle (légalité) suffit-elle à garantir sa valeur morale (légitimité, justice) ? Autrement dit, \emph{l'obéissance à la loi dispense-t-elle de l'examen de sa justice ?}
      \item \emph{Annonce du plan :} Nous verrons d'abord que la loi \emph{revendique} d'être juste par sa forme et sa procédure (I), puis que l'histoire et la philosophie montrent qu'elle \emph{peut être radicalement injuste} (II), enfin que l'État de droit moderne tente d'\emph{articuler légalité et justice} par des garanties institutionnelles (III).
      \end{itemize}

\item \textbf{Plan détaillé :}

\cle{I. La loi porte en elle une prétention à la justice (légalité formelle comme condition de la justice).}
\begin{itemize}
\item \textbf{A. La procédure législative comme garantie d'impartialité et de généralité.}
      \textit{Argument :} La loi s'applique à tous (généralité), sans viser personne (abstraction), votée par les représentants du peuple (légitimité démocratique). Cela évite l'arbitraire du juge ou du roi. \textit{Auteur :} \textbf{Kelsen} (théorie pure du droit : validité = respect de la norme fondamentale), \textbf{Hart} (règle de reconnaissance). \textit{Exemple camerounais :} Vote de la Loi de Finances à l'Assemblée Nationale : débat public, amendements, promulgation par le Président $\to$ sécurité juridique pour les contribuables.
\item \textbf{B. Le positivisme : séparation du droit et de la morale (validité $\neq$ valeur).}
      \textit{Argument :} Pour le positivisme, une loi est « juste » juridiquement si elle est valide (promulguée selon la procédure). La justice est un idéal extra-juridique. \textit{Auteur :} \textbf{Kelsen} (« Le droit est ce qui est posé »), \textbf{Hart} (thèse de la séparation). \textit{Exemple :} Les lois d'apartheid en Afrique du Sud ou ségrégationnistes aux USA étaient \emph{juridiquement valides} (procédure respectée) mais \emph{moralement injustes}.
\end{itemize}

\cle{II. La loi positive peut être substantiellement injuste (critique par le droit naturel et l'histoire).}
\begin{itemize}
\item \textbf{A. Les lois iniques : légalité sans légitimité.}
      \textit{Argument :} Des lois respectant la forme peuvent violer les droits fondamentaux (droit à la vie, dignité, non-discrimination). \textit{Auteur :} \textbf{Saint Augustin} (\emph{De libero arbitrio} : \emph{« Lex injusta non est lex »}), \textbf{Saint Thomas} (loi humaine contraire à la loi naturelle = corruption). \textit{Exemple camerounais :} Le \emph{Code de l'indigénat} (colonial) ou le \emph{travail forcé} (lois 1920-1945) : lois positives votées par l'autorité coloniale, légalement applicables, mais crimes contre l'humanité (injustes par nature).
\item \textbf{B. Le droit naturel et les droits de l'homme comme critères supérieurs.}
      \textit{Argument :} Il existe des principes supérieurs (dignité, droits inaliénables) qui jugent la loi positive. \textit{Auteur :} \textbf{Locke} (droits naturels vie/liberté/propriété antérieurs à l'État), \textbf{DDHC 1789 art. 2}, \textbf{Déclaration universelle 1948}. \textit{Exemple :} Loi n°2016/007 du 12 juillet 2016 (Code pénal camerounais) art. 347-1 (homosexualité) : critiquée par ONG internationales comme violant droit à la vie privée et non-discrimination (droit naturel/DH) bien que légale positivement.
\end{itemize}

\cle{III. Vers la conformité de la loi à la justice : l'État de droit constitutionnel.}
\begin{itemize}
\item \textbf{A. Le contrôle de constitutionnalité : juge de la conformité loi/justice.}
      \textit{Argument :} La Constitution (norme suprême) inscrit les droits fondamentaux ; le juge constitutionnel censure les lois contraires. \textit{Auteur :} \textbf{Kelsen} (justice constitutionnelle), \textbf{Favoreu} (juridictionnalisation du politique). \textit{Exemple camerounais :} \textbf{Conseil Constitutionnel} (loi n°96/06 du 18 janvier 1996 modifiée) : contrôle a priori (saisine Président/PM/1/10 députés) et a posteriori (QPC via exception d'inconstitutionnalité depuis loi 2018). Ex. : Décision n°001/2019 sur la loi électorale.
\item \textbf{B. La démocratie participative et l'État de droit comme garde-fous permanents.}
      \textit{Argument :} La justice de la loi ne tient pas qu'au juge ; elle suppose un processus législatif transparent, consultation société civile, respect procédural (débat, amendements), évaluation d'impact. \textit{Auteur :} \textbf{Habermas} (droit comme medium d'intégration sociale par le discours), \textbf{Rawls} (devoir de civisme, overlapping consensus). \textit{Exemple camerounais :} \emph{Grand Dialogue National (2019)} et lois sur le statut spécial des régions NOSO : tentative d'articulation légalité (loi) / justice (réponse aux griefs identitaires, décentralisation).
\end{itemize}

\item \textbf{Conclusion :}
      \begin{itemize}
      \item \emph{Synthèse :} La loi n'est \emph{pourtant} pas \emph{toujours} juste : sa validité formelle (légalité) est nécessaire mais non suffisante. L'histoire (lois coloniales, ségrégationnistes) et la philosophie (droit naturel) prouvent le divorce possible. L'État de droit moderne (constitutionnalisme, contrôle juridictionnel, démocratie délibérative) tente de réduire cet écart en soumettant la loi à des normes supérieures (Constitution, DH).
      \item \emph{Ouverture :} Mais le juge constitutionnel lui-même interprète ; le risque de « gouvernement des juges » ou d'instrumentalisation politique demeure. La justice de la loi reste un \emph{idéal régulateur} (Kant) jamais pleinement atteint, exigeant une vigilance citoyenne permanente. Peut-on aller jusqu'à une \emph{justice transitionnelle} pour juger les lois du passé ?
      \end{itemize}
\end{enumerate}

\textbf{Points de vigilance (Bac) :}
\begin{itemize}
\item Ne pas confondre « loi injuste » et « loi inconstitutionnelle » (toute loi inconstitutionnelle est injuste, mais une loi peut être constitutionnelle et moralement discutable).
\item Citer explicitement la distinction \cle{légalité / légitimité} (Weber) ou \cle{validité / valeur} (Kelsen).
\item Au moins un exemple camerounais précis (loi, décision, événement) par grande partie.
\item Structure dialectique claire : Thèse (loi = justice formelle) / Antithèse (loi = injustice possible) / Synthèse (institutions pour rapprocher loi et justice).
\end{itemize}
\end{corrige}

\begin{corrige}[Exercice 10 — Explication de texte : Aristote, Éthique à Nicomaque, Livre V]
\begin{enumerate}
\item \textbf{Thèse principale :} La justice politique se divise en deux espèces distinctes — la \cle{justice distributive} (répartition proportionnelle des biens communs selon le mérite) et la \cle{justice corrective} (rétablissement de l'égalité arithmétique dans les transactions entre particuliers) — mais le critère du « mérite » fondant la distribution est \emph{relatif au régime politique} (démocratie, oligarchie, aristocratie), ce qui empêche tout consensus universel sur la justice distributive.
\item \textbf{Deux formes de justice distinguées :}
      \begin{itemize}
      \item \cle{Justice distributive (géométrique) :} Concerne la relation \emph{Communauté $\to$ Particuliers}. Elle répartit les biens communs (honneurs, argent, charges) \emph{proportionnellement} au mérite de chacun (A/B = C/D). C'est une égalité \emph{proportionnelle} : « à chacun selon sa valeur ».
      \item \cle{Justice corrective (arithmétique) :} Concerne les relations \emph{Particuliers $\leftrightarrow$ Particuliers} (contrats volontaires : vente, prêt ; ou involontaires : vol, coups). Elle traite les parties comme \emph{égales} (indifférence au mérite, rang, richesse) et vise à \emph{rétablir l'égalité arithmétique} (même valeur avant/après) lésée par le délit ou le contrat rompu.
      \end{itemize}
\item \textbf{Égalité proportionnelle et désaccord sur le mérite :}
      \begin{itemize}
      \item \emph{Égalité proportionnelle :} Répartition où les parts sont proportionnelles aux mérites : $\frac{\text{Part}_A}{\text{Mérite}_A} = \frac{\text{Part}_B}{\text{Mérite}_B}$. Ex : si A a 2 fois plus de mérite que B, A reçoit 2 fois plus.
      \item \emph{Pourquoi « tous ne conviennent pas » :} Le « mérite » (\emph{axia}) n'est pas une donnée objective unique ; c'est une \emph{construction politique}. Pour les \textbf{démocrates}, le mérite = \cle{liberté} (citoyenneté égale) $\to$ partage égalitaire. Pour les \textbf{oligarques}, le mérite = \cle{richesse} (capacité contributive) $\to$ partage inégalitaire favorisant les riches. Pour les \textbf{aristocrates}, le mérite = \cle{vertu} (excellence morale/intellectuelle) $\to$ partage aux meilleurs. \emph{Le critère de justice distributive dépend donc de la définition de la justice que se donne le régime (sa constitution).}
      \end{itemize}
\item \textbf{Exemple camerounais : Attribution des bourses d'excellence (MINESUP).}
      \begin{itemize}
      \item \emph{Critère « mérite scolaire pur » (notes) :} Approche \emph{aristocratique/oligarchique} (excellence académique). Avantage : récompense l'effort, l'intelligence. Inconvénient : favorise les lycées urbains bien dotés, creuse les inégalités territoriales.
      \item \emph{Critère « mérite + besoin / origine régionale » (quotas) :} Approche \emph{démocratique/rawlsienne} (équité, correction des inégalités de départ). Avantage : justice sociale, représentation nationale. Inconvénient : sentiment d'injustice chez les excellents exclus, risque de baisse du niveau.
      \item \emph{Analyse :} Le choix du critère \emph{transforme la justice de la répartition}. Il n'y a pas de « justice de la bourse » unique ; il y a un choix politique du type de mérite valorisé.
      \end{itemize}
\item \textbf{Critère unique et objectif ?} \textbf{Non.} Le texte affirme explicitement : \emph{« Tous ne conviennent pas sur ce qu'est le mérite »} et énumère trois définitions rivales (liberté, richesse, vertu) liées à trois régimes. Pour Aristote, le mérite est \emph{relatif à la constitution de la cité} (\emph{politeia}). Il n'existe pas de critère unique, objectif et universel du mérite dans le texte ; la justice distributive est \emph{politiquement située}.
\end{enumerate}

\textbf{Points de vigilance (Bac) :}
\begin{itemize}
\item Distinguer clairement \emph{justice distributive} (géométrique, communauté/particuliers) et \emph{justice corrective} (arithmétique, particuliers/particuliers).
\item Expliquer la formule mathématique de la proportionnalité géométrique (A/B = C/D).
\item Ne pas dire « Aristote préfère l'aristocratie » : il \emph{constate} la relativité du mérite aux régimes (même s'il a sa préférence pour la vertu en politique).
\item L'exemple camerounais doit montrer \emph{concrètement} comment le changement de critère change la liste des bénéficiaires.
\end{itemize}
\end{corrige}

\begin{corrige}[Exercice 11 — Sujet de débat argumenté type Probatoire : « Faut-il traiter tout le monde de la même manière pour être juste ? »]
\begin{enumerate}
\item \textbf{Deux thèses opposées :}
      \begin{itemize}
      \item \cle{Thèse de l'égalité formelle (arithmétique) :} \textbf{OUI.} La justice exige l'impartialité aveugle : même règle, même traitement pour tous, quelles que soient les différences (origine, sexe, richesse, talent). C'est la condition de la sécurité juridique, de la non-discrimination et de l'État de droit (Art. 1 et 6 DDHC 1789 : « Les hommes naissent et demeurent libres et égaux en droits », « La loi doit être la même pour tous »).
      \item \cle{Thèse de l'équité proportionnelle (géométrique) :} \textbf{NON.} Traiter identiquement des situations \emph{différentes} produit de l'injustice (inégalité réelle). La justice consiste à « donner à chacun ce qui lui revient » (\emph{suum cuique}) selon des critères pertinents : mérite, besoin, contribution, situation de handicap. C'est l'égalité \emph{proportionnelle} (Aristote, \emph{Éthique à Nicomaque} V) ou la \cle{différence} (Rawls, \emph{Théorie de la justice}) : les inégalités de traitement sont justes si elles profitent aux plus défavorisés.
      \end{itemize}
\item \textbf{Argumentation personnelle (synthèse $\sim$20 lignes) :}
      \textbf{Thèse (Égalité formelle) :} Dans un État de droit, la \cle{justice commutative} et la \cle{procédure} exigent l'égalité arithmétique : même peine pour même délit (Code pénal), même salaire pour même travail (Code du travail), anonymat des concours. C'est la protection contre l'arbitraire, le favoritisme, la corruption. \textbf{Antithèse (Équité proportionnelle) :} Mais l'égalité formelle devient injustice \emph{réelle} face aux inégalités de départ. \cle{Justice distributive} (Aristote) et \cle{principe de différence} (Rawls) imposent de traiter \emph{inégalement} pour établir une égalité des chances effective : aménagement d'épreuves pour handicapés, quotas filles/régions défavorisées (Cameroun), impôt progressif, discrimination positive. \textbf{Synthèse (Ricœur / Articulation) :} La justice n'est pas \emph{soit} l'égalité \emph{soit} l'équité, mais leur \emph{articulation contextuelle} (Ricœur, \emph{Le Juste}). L'égalité formelle est le \emph{plancher} infranchissable (droits fondamentaux, procédure) ; l'équité proportionnelle est le \emph{plafond} visé pour la répartition des ressources et chances (social, éducatif, fiscal). La loi doit être générale dans sa forme (égalité), mais prévoir des \emph{critères de différenciation justifiés} (équité) pour corriger les effets de la loi uniforme sur des situations hétérogènes.
\item \textbf{Contextes de priorité :}
      \begin{itemize}
      \item \emph{Famille :} \textbf{Équité proportionnelle} (besoins selon âge, santé, capacité ; « à chacun selon ses besoins »).
      \item \emph{École :} \textbf{Égalité des chances (équité) :} même programme (formel) + pédagogie différenciée / aides ciblées (proportionnel) pour corriger inégalités sociales.
      \item \emph{Travail :} \textbf{Égalité formelle} (salaire égal travail égal, non-discrimination) + \textbf{Équité} (formation, promotion selon mérite/ancienneté, aménagement handicap).
      \item \emph{État (loi pénale, procédure, vote) :} \textbf{Égalité formelle stricte} (un citoyen = une voix, même loi pour tous). \emph{État (politiques publiques, fiscalité, santé, éducation) :} \textbf{Équité proportionnelle} (redistribution, discrimination positive, universalisme ciblé).
      \end{itemize}
\end{enumerate}

\textbf{Points de vigilance (Probatoire) :}
\begin{itemize}
\item Ne pas opposer bêtement « égalité » et « justice » : l'égalité \emph{est} une composante de la justice (commutative), mais pas la seule (distributive).
\item Citer \textbf{Aristote} (distributive/commutative, proportionnelle/arithmétique) et \textbf{Rawls} (différence) ou \textbf{Ricœur} (articulation) comme exigé.
\item La synthèse doit être \emph{contextuelle} : dire « ça dépend » en précisant \emph{de quoi} ça dépend (domaine : droits / ressources ; nature de la règle : procédure / répartition).
\item Respecter la longueur ($\sim$20 lignes = $\sim$200-250 mots).
\end{itemize}
\end{corrige}

\newpage

\coursec{Fiche bilan}

\coursec{Le Droit et la Justice}

\begin{popfigure}
\centering
\begin{tikzpicture}[
  mindmap,
  concept color=popblue,
  level 1/.style={level distance=3.5cm, sibling angle=60},
  level 2/.style={level distance=2.5cm, sibling angle=45},
  every node/.style={font=\small, text width=2.8cm, align=center},
  branch1/.style={concept color=popblue, edge from parent path={(\tikzparentnode) -- (\tikzchildnode)}},
  branch2/.style={concept color=poporange, edge from parent path={(\tikzparentnode) -- (\tikzchildnode)}},
  branch3/.style={concept color=popgreen, edge from parent path={(\tikzparentnode) -- (\tikzchildnode)}},
  branch4/.style={concept color=poppurple, edge from parent path={(\tikzparentnode) -- (\tikzchildnode)}},
  branch5/.style={concept color=poppink, edge from parent path={(\tikzparentnode) -- (\tikzchildnode)}},
  branch6/.style={concept color=popgold, edge from parent path={(\tikzparentnode) -- (\tikzchildnode)}}
]
\node[concept, text width=3.2cm] {Le Droit et la Justice}
  child[grow=90, branch1] {node[concept, align=center]{Définitions\\Droit, Justice, Loi}}
  child[grow=30, branch2] {node[concept, align=center]{Fondements\\Naturel, Positif, Social, Divin}}
  child[grow=-30, branch3] {node[concept, align=center]{Justice \& Égalité\\Arithmétique, Géométrique, Équité}}
  child[grow=-90, branch4] {node[concept, align=center]{Formes de Justice\\Commutative, Distributive, Réparatrice, Pénale}}
  child[grow=-150, branch5] {node[concept, align=center]{Idée Absolue\\Idéal régulateur, Critique du positif}}
  child[grow=150, branch6] {node[concept, align=center]{Méthode Bac\\Dissertation, Commentaire, Exemples CM}};
\end{tikzpicture}
\legende{Carte mentale de la leçon : concepts clés, fondements, formes de justice et démarche méthodologique.}
\end{popfigure}

\begin{definition}[Définitions fondamentales]
Le \cle{Droit} est l'ensemble des règles de conduite obligatoires, sanctionnées par la force publique (l'État), qui organisent la vie en société. On distingue :
\begin{itemize}
  \item le \cle{Droit objectif} : l'ordre juridique considéré comme un système de normes (ex. : le Code pénal camerounais) ;
  \item les \cle{Droits subjectifs} : les prérogatives reconnues à chaque individu (ex. : droit de propriété, liberté d'expression).
\end{itemize}
La \cle{Justice} est la vertu morale consistant à rendre à chacun ce qui lui est dû (\emph{suum cuique tribuere}) ; elle désigne la conformité au droit et à l'équité. C'est une valeur idéale qui sert de critère de jugement au droit positif.
La \cle{Loi} est une règle générale, abstraite, impersonnelle et obligatoire, votée par le Parlement (pouvoir législatif). Hiérarchie des normes au Cameroun : Constitution (1996, révisée 2008) $>$ Traités et accords internationaux $>$ Lois $>$ Règlements (décrets, arrêtés).
\cle{Légalité} vs \cle{Légitimité} : la légalité est la conformité à la loi écrite ; la légitimité est la conformité à la justice, au consentement du peuple ou à la morale. Exemple : les lois d'apartheid ou les lois esclavagistes étaient \emph{légales} mais \emph{illégitimes}.
\end{definition}

\begin{propriete}[Les fondements du droit : tableau récapitulatif]
\begin{center}
\begin{tabularx}{\linewidth}{|l|X|X|}\hline
\rowcolor{popblueL}
\textbf{Fondement} & \textbf{Thèse centrale} & \textbf{Auteurs / Références clés} \\ \hline
\cle{Droit naturel} & Existence d'un droit universel, antérieur aux lois humaines, fondé sur la nature humaine et la raison. Il sert de critère de validité au droit positif. & Antigone (Sophocle), Stoïciens, \emph{Thomas d'Aquin} (Somme théologique), \emph{John Locke}, \emph{Emmanuel Kant}, Déclarations des droits (1789, 1948). \\ \hline
\cle{Droit positif / Légalisme} & Le droit se confond avec la loi écrite, posée par le souverain ou l'État. Pas de droit hors la loi. La validité juridique dépend uniquement de la procédure d'élaboration (pureté de la norme). & \emph{Thomas Hobbes} (Léviathan), \emph{Hans Kelsen} (Théorie pure du droit, norme fondamentale), \emph{H.L.A. Hart}, Positivisme juridique (École de l'exégèse). \\ \hline
\cle{Fondement social / Sociologique} & Le droit émerge des faits sociaux, des coutumes et des besoins de cohésion du groupe. L'État ne fait que sanctionner des règles préexistantes (droit vivant). & \emph{Émile Durkheim} (solidarité mécanique/organique), \emph{Eugen Ehrlich} (droit vivant), École historique allemande (\emph{Friedrich Carl von Savigny}). \\ \hline
\cle{Fondement divin / Théologique} & Le droit découle de la volonté de Dieu ; c'est la source suprême de l'obligation. La loi humaine n'est valide que si elle ne contredit pas la loi divine. & Droit canonique, \emph{Saint Augustin} (Cité de Dieu), \emph{Thomas d'Aquin} (loi éternelle $>$ loi naturelle $>$ loi humaine). \\ \hline
\end{tabularx}
\end{center}
\end{propriete}

\begin{propriete}[Justice et égalité : les trois registres]
\begin{itemize}
  \item \cle{Égalité arithmétique} : même traitement pour tous (même part, même peine, même droit de vote). Elle fonde la \emph{justice commutative} (échanges, contrats, réparations entre particuliers). \textit{Exemple camerounais :} force obligatoire du contrat (Code civil camerounais, Livre III / Acte uniforme OHADA sur le droit des contrats).
  \item \cle{Égalité géométrique (proportionnelle)} : traiter chacun proportionnellement à son mérite, à son besoin, à sa contribution ou à son rang. Elle fonde la \emph{justice distributive} (répartition des honneurs, richesses, charges publiques). \textit{Exemples camerounais :} Impôt sur le Revenu des Personnes Physiques (IRPP) progressif ; quotas de genre (30\,\%) dans les listes électorales (Code électoral) ; bourses d'excellence MINESUP sur critères de mérite.
  \item \cle{Équité (\emph{Epikeia})} : correction de la rigueur de la loi générale par le juge face à un cas particulier imprévu par le législateur (Aristote, \emph{Éthique à Nicomaque}, Livre V). Elle permet d'atteindre la justice matérielle là où la loi formelle est aveugle ou injuste. \textit{Exemple :} pouvoir d'appréciation du juge pour les circonstances atténuantes (Code de procédure pénale).
\end{itemize}
\end{propriete}

\begin{propriete}[Les quatre formes de justice (Aristote, Thomas d'Aquin, Droit moderne)]
\begin{itemize}
  \item \textbf{Justice commutative} : régit les rapports entre particuliers (contrats, délits civils, responsabilité délictuelle). Principe : égalité arithmétique (restitution \emph{in integrum}). \textit{Exemple :} achat d'un terrain à Yaoundé, respect des clauses contractuelles, réparation en cas de vice caché ou de dol.
  \item \textbf{Justice distributive} : régit le rapport entre la société (État, collectivité) et l'individu (répartition des biens, honneurs, charges, services publics). Principe : égalité géométrique selon un critère pertinent (mérite, besoin, capacité contributive). \textit{Exemple :} recrutement à la Fonction publique par concours (mérite) ; allocation de bourses sociales (besoin) ; fiscalité progressive (capacité).
  \item \textbf{Justice réparatrice (ou corrective)} : vise à rétablir l'équilibre rompu par un dommage injuste (responsabilité civile). Elle concerne la réparation du préjudice (dommages-intérêts). \textit{Exemple :} accident de la route sur l'axe Douala-Yaoundé, indemnisation de la victime par l'assurance du responsable (Code des assurances / Code civil).
  \item \textbf{Justice pénale (répressive)} : sanction par l'État de l'infraction à l'ordre public (infraction : contravention, délit, crime). Buts : réparation sociale, dissuasion (prévention générale), réinsertion du condamné (prévention spéciale). \textit{Exemple :} procès devant le Tribunal de Grande Instance (TGI) pour vol ou homicide, peine de prison et/ou amende (Code pénal camerounais).
\end{itemize}
\end{propriete}

\begin{propriete}[De la justice relative à l'idée absolue de justice]
\begin{itemize}
  \item \cle{Justice relative (positive)} : celle qui s'applique \emph{hic et nunc}, selon les lois en vigueur dans un État donné. Elle est historique, imparfaite, perfectible, parfois injuste. \textit{« La justice des hommes est imparfaite » (Kant, Métaphysique des mœurs).}
  \item \cle{Idée absolue de justice (Idéal régulateur)} : horizon indépassable (Droit naturel, Droits de l'Homme, Impératif catégorique kantien : « Agis de telle sorte que tu traites l'humanité [...] toujours comme une fin, jamais simplement comme un moyen »). Elle est le critère de critique permanente du droit positif.
  \item \cle{Tension dynamique} : le droit positif doit tendre vers la justice absolue sans jamais l'atteindre pleinement. Rôle du juge (interprétation conforme à la Constitution), du législateur (réformes), du citoyen (résistance civile, vote) : réduire l'écart. Au Cameroun : contrôle de constitutionnalité par le Conseil Constitutionnel (Loi n°96/06 du 18 janvier 1996), juridictions administratives et judiciaires.
\end{itemize}
\end{propriete}

\begin{methode}[Méthode dissertation et commentaire (Bac Probatoire / Baccalauréat Technique)]
\begin{enumerate}
  \item \textbf{Analyser le sujet} : définir rigoureusement les termes (Droit, Justice, Loi, Égalité, Légitimité...), identifier la tension philosophique sous-jacente (Légalité vs Légitimité, Droit positif vs Droit naturel, Justice formelle vs Justice matérielle, Relatif vs Absolu).
  \item \textbf{Problématiser} : transformer le sujet en question philosophique ouverte. Exemples types : « La loi suffit-elle à faire la justice ? » / « L'égalité formelle masque-t-elle des inégalités réelles ? » / « Le droit positif peut-il se passer de fondement moral ? » / « Obéir aux lois, est-ce être juste ? »
  \item \textbf{Construire un plan dialectique (Thèse / Antithèse / Synthèse)} :
    \begin{itemize}
      \item \textbf{Thèse} : Le droit positif assure l'ordre, la sécurité juridique, la prévisibilité et l'égalité formelle (Kelsen, Hobbes, Durkheim : fait social). La loi est la source unique du droit.
      \item \textbf{Antithèse} : Le droit positif peut être injuste (lois racistes, esclavagistes, discriminatoires) ; la justice vraie exige l'équité, le droit naturel, les Droits de l'Homme (Antigone, Kant, Rawls, Kelsen lui-même admet une norme fondamentale présupposée). La légalité $\neq$ légitimité.
      \item \textbf{Synthèse} : L'État de droit moderne = droit positif \textbf{contrôlé} par la justice (bloc de constitutionnalité, contrôle de conventionalité, Conseil Constitutionnel camerounais, CEDH). La loi est \emph{nécessaire} (instrument) mais \emph{non suffisante} (fin) ; la justice est l'horizon critique et régulateur.
    \end{itemize}
  \item \textbf{Illustrer par des exemples camerounais (OBLIGATOIRE)} :
    \begin{itemize}
      \item Code pénal : criminalisation de l'homosexualité (art. 347 bis), encadrement de l'avortement (art. 339) $\to$ débat légalité/légitimité, droit naturel/DDH.
      \item Législation familiale : dot coutumière vs dot civile, successions (Ordonnances 74-1, 74-2, 81-02) $\to$ pluralisme juridique, justice distributive/réparatrice.
      \item Jurisprudence : arrêts de la Cour Suprême (chambre judiciaire/administrative) et du Conseil Constitutionnel (contrôle lois électorales, libertés publiques).
      \item Vie courante : corruption (justice à deux vitesses), justice populaire / expéditive (\"justice mob\"), fiscalité (IRPP, TVA), conflits fonciers (droit écrit vs droit coutumier).
    \end{itemize}
  \item \textbf{Rédiger} : introduction (accroche, définitions, problématique, annonce de plan), développement (arguments structurés, auteurs cités, exemples intégrés), conclusion (réponse à la problématique, ouverture). Connecteurs logiques : \emph{donc, cependant, par ailleurs, en somme, en définitive}.
\end{enumerate}
\end{methode}

\begin{aretenir}[Checklist avant le Bac Probatoire / Baccalauréat Technique]
\begin{itemize}
  \item $\square$ Je définis précisément : Droit (objectif/subjectif), Justice, Loi, Légalité, Légitimité, Équité.
  \item $\square$ Je distingue les 4 fondements du droit (Naturel, Positif, Social, Divin) et cite au moins un auteur majeur par fondement.
  \item $\square$ J'oppose Égalité arithmétique (justice commutative) et Égalité géométrique (justice distributive) avec exemples concrets camerounais.
  \item $\square$ Je caractérise les 4 formes de justice : Commutative, Distributive, Réparatrice, Pénale.
  \item $\square$ J'explique la distinction Justice relative (positive) / Idée absolue de justice (idéal régulateur kantien / DDH).
  \item $\square$ Je sais argumenter le débat classique : « Une loi injuste est-elle une loi ? » (Thèses : Positivisme / Jusnaturalisme / École sociologique).
  \item $\square$ Je maîtrise la structure dissertation : Analyse $\to$ Problématique $\to$ Plan dialectique (Thèse/Antithèse/Synthèse) $\to$ Conclusion ouverte.
  \item $\square$ Je possède 3 à 5 exemples \textbf{camerounais} récents ou classiques (Code pénal, Constitution 1996/2008, Jurisprudence Conseil Constitutionnel/Cour Suprême, Coutume, Fiscalité, Fonction publique).
  \item $\square$ Je sais commenter un texte : Identifier thèse, concepts, arguments, structure, portée critique, et confronter à mes connaissances.
  \item $\square$ Je rédige en français correct, phrases complètes, connecteurs logiques, sans familier ni anglicismes.
  \item $\square$ Je gère mon temps : 3h-4h dissertation (brouillon soigné), 2h-3h commentaire (lecture active, analyse, rédaction).
  \item $\square$ Je relis : orthographe, noms d'auteurs, concepts soulignés (\cle{}), exemples présents, réponse explicite à la problématique.
\end{itemize}
\end{aretenir}

\findecours

\end{document}
